\makeatletter
\def\input@path{{./}{../}{../../}{../../../}{../../../../}{preamble/}{../preamble/}{../../preamble/}{../../../preamble/}{../../../../preamble/}}
\makeatother
\input{preamble_G10_T1_C2C3.tex}
\input{graphs/highlights.tex}
% \GTenHighlightsfalse % Uncomment to disable instructional highlights.
\input{graphs/specific_highlights.tex}

\title{Grade 10 Term 1 Lesson 1 Practice Activity 2\\Revenue, Costs, Profit, and Payoff Tables}
\author{}
\date{}
\begin{document}
\maketitle
\setcounter{tocdepth}{1}
\setcounter{secnumdepth}{2}
\phantomsection\label{toc}
\tableofcontents

% ================================================================
\section*{Evaluated criteria}
\begin{enumerate}[label=\textbf{C\arabic*}]
\item Characterises decision-making problems under uncertainty by identifying decision alternatives, states of nature, consequences, and the economic or financial context of the decision.
\item[\textbf{C3}] Calculates profit as revenue minus costs across states of nature and organises the resulting payoffs in payoff tables.
\end{enumerate}

% ================================================================
\section{Introduction}\label{sec:introduction}
\phantomsection\label{tone:legend}
\PracticeHighlightLegend{G10-T1-C1-C3}


\SectionProblemsTableOfContents{
	\SectionProblemLink{sec:problem1-introductory-general-c1-c3}{Introductory General Problem (C1--C3)} \\
	\SectionProblemLink{sec:problem2-advanced-general-c1-c3}{Advanced General Problem (C1--C3)}
}\vspace{8mm}

A decision problem under uncertainty requires a choice of an
\emph{\GTenToneHL{alternative}{alternative}} \(\GTenToneHL{alternative}{A_i}\) before an uncertain event reveals a
\emph{\GTenToneHL{state}{state} \GTenToneHL{state}{of} \GTenToneHL{state}{nature}} \(\GTenToneHL{state}{S_j}\). The \GTenToneHL{alternative}{alternative} is a row of the table and the
state is a column. Their combination produces a \emph{consequence}; its
monetary \emph{\GTenToneHL{payoff}{payoff}} is \GTenToneHL{payoff}{profit} \(\GTenToneHL{payoff}{P}(\GTenToneHL{alternative}{A_i},\GTenToneHL{state}{S_j})\). \emph{\GTenToneHL{revenue}{Revenue}}
\(\GTenToneHL{revenue}{R}(\GTenToneHL{alternative}{A_i},\GTenToneHL{state}{S_j})\) is income earned, \emph{\GTenToneHL{cost}{cost}} \(\GTenToneHL{cost}{C}(\GTenToneHL{alternative}{A_i},\GTenToneHL{state}{S_j})\) is the value of
resources used, and
\[
\GTenToneHL{payoff}{P}(\GTenToneHL{alternative}{A_i},\GTenToneHL{state}{S_j})=\GTenToneHL{revenue}{R}(\GTenToneHL{alternative}{A_i},\GTenToneHL{state}{S_j})-\GTenToneHL{cost}{C}(\GTenToneHL{alternative}{A_i},\GTenToneHL{state}{S_j}).
\]
A \GTenToneHL{payoff}{payoff} table places every \(\GTenToneHL{payoff}{P}(\GTenToneHL{alternative}{A_i},\GTenToneHL{state}{S_j})\) at row \(\GTenToneHL{alternative}{A_i}\) and column \(\GTenToneHL{state}{S_j}\).
The objective here is to construct that complete table, not to choose the best
\GTenToneHL{alternative}{alternative}.

\GTenToneHL{revenue}{Revenue} and \GTenToneHL{cost}{cost} are independent axes. Their general models are
\[
\begin{aligned}
	\GTenToneHL{revenue}{R}(\GTenToneHL{alternative}{A_i},\GTenToneHL{state}{S_j})
	={}&
	\GTenToneHL{quantity}{q}(\GTenToneHL{alternative}{A_i},\GTenToneHL{state}{S_j})
	\left[
	\GTenToneHL{revenue}{r}(\GTenToneHL{alternative}{A_i})+\GTenToneHL{revenue}{r}(\GTenToneHL{state}{S_j})
	\right]
	+
	\GTenToneHL{revenue}{f_r}(\GTenToneHL{alternative}{A_i})
	+
	\GTenToneHL{revenue}{f_r}(\GTenToneHL{state}{S_j})
	\\
	\GTenToneHL{cost}{C}(\GTenToneHL{alternative}{A_i},\GTenToneHL{state}{S_j})
	={}&
	\GTenToneHL{quantity}{q}(\GTenToneHL{alternative}{A_i},\GTenToneHL{state}{S_j})
	\left[
	\GTenToneHL{cost}{c}(\GTenToneHL{alternative}{A_i})+\GTenToneHL{cost}{c}(\GTenToneHL{state}{S_j})
	\right]
	+
	\GTenToneHL{cost}{f_c}(\GTenToneHL{alternative}{A_i})
	+
	\GTenToneHL{cost}{f_c}(\GTenToneHL{state}{S_j})
\end{aligned}
\]
Therefore,
\[
\begin{aligned}
	\GTenToneHL{payoff}{P}(\GTenToneHL{alternative}{A_i},\GTenToneHL{state}{S_j})
	={}&
	\GTenToneHL{quantity}{q}(\GTenToneHL{alternative}{A_i},\GTenToneHL{state}{S_j})
	\left[
	\GTenToneHL{revenue}{r}(\GTenToneHL{alternative}{A_i})+\GTenToneHL{revenue}{r}(\GTenToneHL{state}{S_j})
	\right]
	+
	\GTenToneHL{revenue}{f_r}(\GTenToneHL{alternative}{A_i})
	+
	\GTenToneHL{revenue}{f_r}(\GTenToneHL{state}{S_j})
	\\
	&-
	\left\{
	\GTenToneHL{quantity}{q}(\GTenToneHL{alternative}{A_i},\GTenToneHL{state}{S_j})
	\left[
	\GTenToneHL{cost}{c}(\GTenToneHL{alternative}{A_i})+\GTenToneHL{cost}{c}(\GTenToneHL{state}{S_j})
	\right]
	+
	\GTenToneHL{cost}{f_c}(\GTenToneHL{alternative}{A_i})
	+
	\GTenToneHL{cost}{f_c}(\GTenToneHL{state}{S_j})
	\right\}.
\end{aligned}
\]

Here a \emph{variable component} is \GTenToneHL{quantity}{quantity} \(\GTenToneHL{quantity}{q}(\GTenToneHL{alternative}{A_i},\GTenToneHL{state}{S_j})\) multiplied by a
per-unit rate; a \emph{fixed component} is a lump sum independent of \GTenToneHL{quantity}{quantity}.
The argument \(\GTenToneHL{alternative}{A_i}\) identifies an \GTenToneHL{alternative}{alternative}-based component, while the
argument \(\GTenToneHL{state}{S_j}\) identifies a state-based component. The \GTenToneHL{quantity}{quantity}
\(\GTenToneHL{quantity}{q}(\GTenToneHL{alternative}{A_i},\GTenToneHL{state}{S_j})\) depends on both arguments and may
change in every cell. Thus, inspect the arguments: \(\GTenToneHL{alternative}{A_i}\) means that the value
depends on the \GTenToneHL{alternative}{alternative}, \(\GTenToneHL{state}{S_j}\) means that it depends on the state, and
\((\GTenToneHL{alternative}{A_i},\GTenToneHL{state}{S_j})\) means that it may depend on both. A fixed lump sum is not
multiplied by \GTenToneHL{quantity}{quantity}. A component absent from one model is zero. A directly
supplied \(\GTenToneHL{revenue}{R}(\GTenToneHL{alternative}{A_i},\GTenToneHL{state}{S_j})\) or \(\GTenToneHL{cost}{C}(\GTenToneHL{alternative}{A_i},\GTenToneHL{state}{S_j})\) needs no decomposition.

\textbf{Reusable procedure.}
\begin{enumerate}[label=\arabic*.]
	\item Identify the \GTenToneHL{alternative}{alternatives}, uncertain event, states, consequences, and
	units.
	
	\item Separate all \GTenToneHL{revenue}{revenue} data from all \GTenToneHL{cost}{cost} data.
	
	\item Classify each item as variable or fixed and as \GTenToneHL{alternative}{alternative}-based or
	state-based; treat absent components as zero.
	
	\item For every cell, calculate \(\GTenToneHL{revenue}{R}(\GTenToneHL{alternative}{A_i},\GTenToneHL{state}{S_j})\) from its \GTenToneHL{revenue}{revenue} model and
	\(\GTenToneHL{cost}{C}(\GTenToneHL{alternative}{A_i},\GTenToneHL{state}{S_j})\) independently from its \GTenToneHL{cost}{cost} model.
	
	\item Subtract \GTenToneHL{revenue}{revenue} minus \GTenToneHL{cost}{costs}:
	\[
	\GTenToneHL{payoff}{P}(\GTenToneHL{alternative}{A_i},\GTenToneHL{state}{S_j})=\GTenToneHL{revenue}{R}(\GTenToneHL{alternative}{A_i},\GTenToneHL{state}{S_j})-\GTenToneHL{cost}{C}(\GTenToneHL{alternative}{A_i},\GTenToneHL{state}{S_j}).
	\]
	
	\item Put \GTenToneHL{alternative}{alternatives} in rows and states in columns; check every cell and
	all units.
\end{enumerate}

\textbf{Paired ladder overview.} AV, SV, AF, and SF denote
\GTenToneHL{alternative}{alternative}-variable, state-variable, \GTenToneHL{alternative}{alternative}-fixed, and state-fixed
components. Within each tier, the \GTenToneHL{cost}{cost} list is a deliberate permutation of the
\GTenToneHL{revenue}{revenue} list; matching structures are not assumed.

\begin{center}
	\small
	\begin{tabularx}{\textwidth}{@{}lXX@{}}
		\toprule
		Tier
		& \GTenToneHL{revenue}{Revenue} structures, in problem order
		& \GTenToneHL{cost}{Cost} structures, in problem order \\
		\midrule
		Direct
		& direct
		& direct \\
		
		One component
		& AV; SV; AF; SF
		& SV; AF; SF; AV \\
		
		Two components
		& AV+SV; AV+AF; AV+SF; SV+AF; SV+SF; AF+SF
		& AV+SF; SV+AF; SV+SF; AF+SF; AV+SV; AV+AF \\
		
		Three components
		& AV+SV+AF; AV+SV+SF; AV+AF+SF; SV+AF+SF
		& AV+SV+SF; AV+AF+SF; SV+AF+SF; AV+SV+AF \\
		
		Complete
		& AV+SV+AF+SF
		& AV+SV+AF+SF \\
		\bottomrule
	\end{tabularx}
\end{center}

% ----------------------------------------------------------------
\subsection{Introductory General Problem (C1--C3)}
\label{sec:problem1-introductory-general-c1-c3}

A school enterprise chooses \(\GTenToneHL{alternative}{A_1}\) or \(\GTenToneHL{alternative}{A_2}\) before demand is \(\GTenToneHL{state}{S_1}\) or
\(\GTenToneHL{state}{S_2}\). The consequence is \GTenToneHL{payoff}{profit} in thousand USD. Total revenues are
\[
\begin{aligned}
	\GTenToneHL{revenue}{R}(\GTenToneHL{alternative}{A_1},\GTenToneHL{state}{S_1})&=\GTenToneHL{revenue}{20},
	&
	\GTenToneHL{revenue}{R}(\GTenToneHL{alternative}{A_1},\GTenToneHL{state}{S_2})&=\GTenToneHL{revenue}{13},
	\\
	\GTenToneHL{revenue}{R}(\GTenToneHL{alternative}{A_2},\GTenToneHL{state}{S_1})&=\GTenToneHL{revenue}{18},
	&
	\GTenToneHL{revenue}{R}(\GTenToneHL{alternative}{A_2},\GTenToneHL{state}{S_2})&=15.
\end{aligned}
\]
The directly supplied total \GTenToneHL{cost}{costs} are
\[
\begin{aligned}
	\GTenToneHL{cost}{C}(\GTenToneHL{alternative}{A_1},\GTenToneHL{state}{S_1})&=\GTenToneHL{cost}{9},
	&
	\GTenToneHL{cost}{C}(\GTenToneHL{alternative}{A_1},\GTenToneHL{state}{S_2})&=\GTenToneHL{cost}{8},
	\\
	\GTenToneHL{cost}{C}(\GTenToneHL{alternative}{A_2},\GTenToneHL{state}{S_1})&=\GTenToneHL{cost}{10},
	&
	\GTenToneHL{cost}{C}(\GTenToneHL{alternative}{A_2},\GTenToneHL{state}{S_2})&=9.
\end{aligned}
\]

\textbf{Questions/Tasks.}
\begin{enumerate}
	\CriterionTask
	{subsec:problem1-introductory-general-c1}{1}
	{Interpret the decision elements.}

	\CriterionTask
	{subsec:problem1-introductory-general-c3}{3}
	{Build the payoff table.}
\end{enumerate}

\subsubsection*{C1 - Interpreting the decision.}
\phantomsection
\label{subsec:problem1-introductory-general-c1}

\DecisionElementsTable
{\(\GTenToneHL{alternative}{A_1}\); \(\GTenToneHL{alternative}{A_2}\)}
{\(\GTenToneHL{state}{S_1}\); \(\GTenToneHL{state}{S_2}\)}
{Demand revealed after the enterprise choice}
{Profit in thousand USD for each \GTenToneHL{alternative}{alternative}--state pair}

\secInicio[sec:problem1-introductory-general-c1-c3][sec:introduction]

\subsubsection*{C3 - Calculating revenue, cost, and payoff.}
\phantomsection
\label{subsec:problem1-introductory-general-c3}

\begin{center}
	\small
	\begin{tabular}{lrrr}
		\toprule
		Cell
		& \(\GTenToneHL{revenue}{R}(\GTenToneHL{alternative}{A_i},\GTenToneHL{state}{S_j})\)
		& \(\GTenToneHL{cost}{C}(\GTenToneHL{alternative}{A_i},\GTenToneHL{state}{S_j})\)
		& \(\GTenToneHL{payoff}{P}(\GTenToneHL{alternative}{A_i},\GTenToneHL{state}{S_j})=\GTenToneHL{revenue}{R}(\GTenToneHL{alternative}{A_i},\GTenToneHL{state}{S_j})-\GTenToneHL{cost}{C}(\GTenToneHL{alternative}{A_i},\GTenToneHL{state}{S_j})\) \\
		\midrule
		\((\GTenToneHL{alternative}{A_1},\GTenToneHL{state}{S_1})\) & \GTenToneHL{revenue}{20} & \GTenToneHL{cost}{9}  & \(20-9=11\) \\
		\((\GTenToneHL{alternative}{A_1},\GTenToneHL{state}{S_2})\) & \GTenToneHL{revenue}{13} & \GTenToneHL{cost}{8}  & \(13-8=5\) \\
		\((\GTenToneHL{alternative}{A_2},\GTenToneHL{state}{S_1})\) & \GTenToneHL{revenue}{18} & \GTenToneHL{cost}{10} & \(18-10=8\) \\
		\((\GTenToneHL{alternative}{A_2},\GTenToneHL{state}{S_2})\) & \GTenToneHL{revenue}{15} & \GTenToneHL{cost}{9}  & \(15-9=6\) \\
		\bottomrule
	\end{tabular}
\end{center}

\begin{center}
	\begin{tabular}{lrr}
		\toprule
		Alternative & \(\GTenToneHL{state}{S_1}\) & \(\GTenToneHL{state}{S_2}\) \\
		\midrule
		\(\GTenToneHL{alternative}{A_1}\) & \GTenToneHL{payoff}{11} & \GTenToneHL{payoff}{5} \\
		\(\GTenToneHL{alternative}{A_2}\) & \GTenToneHL{payoff}{8}  & \GTenToneHL{payoff}{6} \\
		\bottomrule
	\end{tabular}
\end{center}

\secInicio[sec:problem1-introductory-general-c1-c3][sec:introduction]

% ----------------------------------------------------------------
\subsection{Advanced General Problem (C1--C3)}
\label{sec:problem2-advanced-general-c1-c3}

A community business chooses \(\GTenToneHL{alternative}{A_1}\) or \(\GTenToneHL{alternative}{A_2}\) before demand is \(\GTenToneHL{state}{S_1}\) or
\(\GTenToneHL{state}{S_2}\). The quantities are
\[
\begin{aligned}
	\GTenToneHL{quantity}{q}(\GTenToneHL{alternative}{A_1},\GTenToneHL{state}{S_1})&=\GTenToneHL{quantity}{2},
	&
	\GTenToneHL{quantity}{q}(\GTenToneHL{alternative}{A_1},\GTenToneHL{state}{S_2})&=\GTenToneHL{quantity}{1},
	\\
	\GTenToneHL{quantity}{q}(\GTenToneHL{alternative}{A_2},\GTenToneHL{state}{S_1})&=\GTenToneHL{quantity}{3},
	&
	\GTenToneHL{quantity}{q}(\GTenToneHL{alternative}{A_2},\GTenToneHL{state}{S_2})&=2.
\end{aligned}
\]

In thousand USD, the \GTenToneHL{alternative}{alternative}-based \GTenToneHL{revenue}{revenue} rates are
\[
\GTenToneHL{revenue}{r}(\GTenToneHL{alternative}{A_1})=\GTenToneHL{revenue}{10},
\qquad
\GTenToneHL{revenue}{r}(\GTenToneHL{alternative}{A_2})=\GTenToneHL{revenue}{11},
\]
and the state-based \GTenToneHL{revenue}{revenue} rates are
\[
\GTenToneHL{revenue}{r}(\GTenToneHL{state}{S_1})=\GTenToneHL{revenue}{5},
\qquad
\GTenToneHL{revenue}{r}(\GTenToneHL{state}{S_2})=3.
\]
The \GTenToneHL{alternative}{alternative}-based fixed revenues are
\[
\GTenToneHL{revenue}{f_r}(\GTenToneHL{alternative}{A_1})=\GTenToneHL{revenue}{7},
\qquad
\GTenToneHL{revenue}{f_r}(\GTenToneHL{alternative}{A_2})=\GTenToneHL{revenue}{9},
\]
and the state-based fixed revenues are
\[
\GTenToneHL{revenue}{f_r}(\GTenToneHL{state}{S_1})=\GTenToneHL{revenue}{6},
\qquad
\GTenToneHL{revenue}{f_r}(\GTenToneHL{state}{S_2})=7.
\]

The \GTenToneHL{alternative}{alternative}-based \GTenToneHL{cost}{cost} rates are
\[
\GTenToneHL{cost}{c}(\GTenToneHL{alternative}{A_1})=\GTenToneHL{cost}{3},
\qquad
\GTenToneHL{cost}{c}(\GTenToneHL{alternative}{A_2})=\GTenToneHL{cost}{4},
\]
and the state-based \GTenToneHL{cost}{cost} rates are
\[
\GTenToneHL{cost}{c}(\GTenToneHL{state}{S_1})=\GTenToneHL{cost}{1},
\qquad
\GTenToneHL{cost}{c}(\GTenToneHL{state}{S_2})=2.
\]
The \GTenToneHL{alternative}{alternative}-based fixed \GTenToneHL{cost}{costs} are
\[
\GTenToneHL{cost}{f_c}(\GTenToneHL{alternative}{A_1})=\GTenToneHL{cost}{8},
\qquad
\GTenToneHL{cost}{f_c}(\GTenToneHL{alternative}{A_2})=\GTenToneHL{cost}{10},
\]
and the state-based fixed \GTenToneHL{cost}{costs} are
\[
\GTenToneHL{cost}{f_c}(\GTenToneHL{state}{S_1})=\GTenToneHL{cost}{5},
\qquad
\GTenToneHL{cost}{f_c}(\GTenToneHL{state}{S_2})=3.
\]

\textbf{Questions/Tasks.}
\begin{enumerate}
	\CriterionTask
	{subsec:problem2-advanced-general-c1}{1}
	{Interpret the decision elements.}

	\CriterionTask
	{subsec:problem2-advanced-general-c3}{3}
	{Build the payoff table.}
\end{enumerate}

\subsubsection*{C1 - Interpreting the decision.}
\phantomsection
\label{subsec:problem2-advanced-general-c1}

\DecisionElementsTable
{\(\GTenToneHL{alternative}{A_1}\); \(\GTenToneHL{alternative}{A_2}\)}
{\(\GTenToneHL{state}{S_1}\); \(\GTenToneHL{state}{S_2}\)}
{Demand revealed after the business choice}
{Profit in thousand USD from the complete \GTenToneHL{revenue}{revenue} and \GTenToneHL{cost}{cost} models}

For every \GTenToneHL{alternative}{alternative}--state pair, \GTenToneHL{revenue}{revenue} and \GTenToneHL{cost}{cost} are calculated independently:
\[
\begin{aligned}
	\GTenToneHL{revenue}{R}(\GTenToneHL{alternative}{A_i},\GTenToneHL{state}{S_j})
	={}&
	\GTenToneHL{quantity}{q}(\GTenToneHL{alternative}{A_i},\GTenToneHL{state}{S_j})
	\left[
	\GTenToneHL{revenue}{r}(\GTenToneHL{alternative}{A_i})+\GTenToneHL{revenue}{r}(\GTenToneHL{state}{S_j})
	\right]
	+
	\GTenToneHL{revenue}{f_r}(\GTenToneHL{alternative}{A_i})
	+
	\GTenToneHL{revenue}{f_r}(\GTenToneHL{state}{S_j}),
	\\
	\GTenToneHL{cost}{C}(\GTenToneHL{alternative}{A_i},\GTenToneHL{state}{S_j})
	={}&
	\GTenToneHL{quantity}{q}(\GTenToneHL{alternative}{A_i},\GTenToneHL{state}{S_j})
	\left[
	\GTenToneHL{cost}{c}(\GTenToneHL{alternative}{A_i})+\GTenToneHL{cost}{c}(\GTenToneHL{state}{S_j})
	\right]
	+
	\GTenToneHL{cost}{f_c}(\GTenToneHL{alternative}{A_i})
	+
	\GTenToneHL{cost}{f_c}(\GTenToneHL{state}{S_j}),
	\\
	\GTenToneHL{payoff}{P}(\GTenToneHL{alternative}{A_i},\GTenToneHL{state}{S_j})
	={}&
	\GTenToneHL{revenue}{R}(\GTenToneHL{alternative}{A_i},\GTenToneHL{state}{S_j})-\GTenToneHL{cost}{C}(\GTenToneHL{alternative}{A_i},\GTenToneHL{state}{S_j}).
\end{aligned}
\]

\secInicio[sec:problem2-advanced-general-c1-c3][sec:introduction]

\subsubsection*{C3 - Calculating revenue, cost, and payoff.}
\phantomsection
\label{subsec:problem2-advanced-general-c3}

\begin{center}
	\begin{tabular}{llll}
		\toprule
		Cell
		& \(\GTenToneHL{revenue}{R}(\GTenToneHL{alternative}{A_i},\GTenToneHL{state}{S_j})\)
		& \(\GTenToneHL{cost}{C}(\GTenToneHL{alternative}{A_i},\GTenToneHL{state}{S_j})\)
		& \(\GTenToneHL{payoff}{P}(\GTenToneHL{alternative}{A_i},\GTenToneHL{state}{S_j})=\GTenToneHL{revenue}{R}(\GTenToneHL{alternative}{A_i},\GTenToneHL{state}{S_j})-\GTenToneHL{cost}{C}(\GTenToneHL{alternative}{A_i},\GTenToneHL{state}{S_j})\) \\
		\midrule
		\((\GTenToneHL{alternative}{A_1},\GTenToneHL{state}{S_1})\)
		& \(2(10+5)+7+6=43\)
		& \(2(3+1)+8+5=21\)
		& \(43-21=22\) \\
		
		\((\GTenToneHL{alternative}{A_1},\GTenToneHL{state}{S_2})\)
		& \(1(10+3)+7+7=27\)
		& \(1(3+2)+8+3=16\)
		& \(27-16=11\) \\
		
		\((\GTenToneHL{alternative}{A_2},\GTenToneHL{state}{S_1})\)
		& \(3(11+5)+9+6=63\)
		& \(3(4+1)+10+5=30\)
		& \(63-30=33\) \\
		
		\((\GTenToneHL{alternative}{A_2},\GTenToneHL{state}{S_2})\)
		& \(2(11+3)+9+7=44\)
		& \(2(4+2)+10+3=25\)
		& \(44-25=19\) \\
		\bottomrule
	\end{tabular}
\end{center}

\begin{center}
	\begin{tabular}{lrr}
		\toprule
		Alternative & \(\GTenToneHL{state}{S_1}\) & \(\GTenToneHL{state}{S_2}\) \\
		\midrule
		\(\GTenToneHL{alternative}{A_1}\) & \GTenToneHL{payoff}{22} & \GTenToneHL{payoff}{11} \\
		\(\GTenToneHL{alternative}{A_2}\) & \GTenToneHL{payoff}{33} & \GTenToneHL{payoff}{19} \\
		\bottomrule
	\end{tabular}
\end{center}

\secInicio[sec:problem2-advanced-general-c1-c3][sec:introduction]

% ================================================================

\phantomsection\label{tone:reference}
\subsection*{Notation reference}
\GTenToneReference{C1-C3}
\section{Direct Revenue and Cost Models}
\label{sec:direct-revenue-cost-models}

\SectionProblemsTableOfContents{
	\SectionProblemLink{sec:problem3-local-cafe-investment-c1-c3}{Local Café Investment (C1--C3)} \\
	\SectionProblemLink{sec:problem4-green-energy-production-mix-c1-c3}{Green Energy Production Mix (C1--C3)}
}

\vspace{0.8cm}

In a direct model, the total \GTenToneHL{revenue}{revenue} \(\GTenToneHL{revenue}{R}(\GTenToneHL{alternative}{A_i},\GTenToneHL{state}{S_j})\) and total \GTenToneHL{cost}{cost}
\(\GTenToneHL{cost}{C}(\GTenToneHL{alternative}{A_i},\GTenToneHL{state}{S_j})\) are supplied for every \GTenToneHL{alternative}{alternative}--state pair. For
each cell, the \GTenToneHL{payoff}{payoff} is calculated directly as
\[
\GTenToneHL{payoff}{P}(\GTenToneHL{alternative}{A_i},\GTenToneHL{state}{S_j})=\GTenToneHL{revenue}{R}(\GTenToneHL{alternative}{A_i},\GTenToneHL{state}{S_j})-\GTenToneHL{cost}{C}(\GTenToneHL{alternative}{A_i},\GTenToneHL{state}{S_j}).
\]

% ----------------------------------------------------------------
\subsection{Local Café Investment (C1--C3)}
\label{sec:problem3-local-cafe-investment-c1-c3}

A small café is deciding whether to launch a new line of artisan desserts.
The transaction is daily café sales, where \GTenToneHL{revenue}{revenue} is earned from sales and
\GTenToneHL{cost}{costs} include the resources required to operate each menu option. The owner
must choose between:

\begin{multicols}{2}
	\begin{enumerate}[label=\arabic*.]
		\item \(\GTenToneHL{alternative}{A_1}\): \textbf{Launch the dessert line}
		\item \(\GTenToneHL{alternative}{A_2}\): \textbf{Keep the current menu}
	\end{enumerate}
	
	The possible \GTenToneHL{state}{states} \GTenToneHL{state}{of} \GTenToneHL{state}{nature} are:
	\begin{enumerate}[label=\arabic*.]
		\item \(\GTenToneHL{state}{S_1}\): High traffic
		\item \(\GTenToneHL{state}{S_2}\): Low traffic
	\end{enumerate}
	
	\columnbreak
	
	The directly estimated revenues and \GTenToneHL{cost}{costs}, in thousand USD, are:
	\begin{enumerate}[label=\arabic*.]
		\item Launch dessert line:
		\begin{enumerate}[label=\roman*)]
			\item High traffic:
			\(\GTenToneHL{revenue}{R}(\GTenToneHL{alternative}{A_1},\GTenToneHL{state}{S_1})=\GTenToneHL{revenue}{58}\), \(\GTenToneHL{cost}{C}(\GTenToneHL{alternative}{A_1},\GTenToneHL{state}{S_1})=\GTenToneHL{cost}{16}\)
			
			\item Low traffic:
			\(\GTenToneHL{revenue}{R}(\GTenToneHL{alternative}{A_1},\GTenToneHL{state}{S_2})=\GTenToneHL{revenue}{27}\), \(\GTenToneHL{cost}{C}(\GTenToneHL{alternative}{A_1},\GTenToneHL{state}{S_2})=\GTenToneHL{cost}{19}\)
		\end{enumerate}
		
		\item Keep current menu:
		\begin{enumerate}[label=\roman*)]
			\item High traffic:
			\(\GTenToneHL{revenue}{R}(\GTenToneHL{alternative}{A_2},\GTenToneHL{state}{S_1})=\GTenToneHL{revenue}{49}\), \(\GTenToneHL{cost}{C}(\GTenToneHL{alternative}{A_2},\GTenToneHL{state}{S_1})=\GTenToneHL{cost}{21}\)
			
			\item Low traffic:
			\(\GTenToneHL{revenue}{R}(\GTenToneHL{alternative}{A_2},\GTenToneHL{state}{S_2})=\GTenToneHL{revenue}{41}\), \(\GTenToneHL{cost}{C}(\GTenToneHL{alternative}{A_2},\GTenToneHL{state}{S_2})=\GTenToneHL{cost}{21}\)
		\end{enumerate}
	\end{enumerate}
\end{multicols}

\textbf{Questions/Tasks.}
\begin{enumerate}
	\CriterionTask
	{subsec:problem3-local-cafe-investment-c1}{1}
	{Interpreting decision alternatives, events, consequences, and states.}

	\CriterionTask
	{subsec:problem3-local-cafe-investment-c3}{3}
	{Calculating profit and building the payoff table.}
\end{enumerate}

\subsubsection*{C1 - Interpreting decision alternatives, events, consequences, and states.}
\phantomsection
\label{subsec:problem3-local-cafe-investment-c1}

\DecisionElementsTable
{Launch dessert line; Keep current menu}
{High traffic; Low traffic}
{Actual foot traffic after the decision}
{Profit in thousand USD for each \GTenToneHL{alternative}{alternative}--state pair}

\secInicio[sec:problem3-local-cafe-investment-c1-c3][sec:direct-revenue-cost-models]

\subsubsection*{C3 - Calculating profit and building the payoff table.}
\phantomsection
\label{subsec:problem3-local-cafe-investment-c3}

For every \GTenToneHL{alternative}{alternative}--state pair, calculate
\[
\GTenToneHL{payoff}{P}(\GTenToneHL{alternative}{A_i},\GTenToneHL{state}{S_j})=\GTenToneHL{revenue}{R}(\GTenToneHL{alternative}{A_i},\GTenToneHL{state}{S_j})-\GTenToneHL{cost}{C}(\GTenToneHL{alternative}{A_i},\GTenToneHL{state}{S_j}).
\]

\begin{center}
	\small
	\begin{tabular}{lrrr}
		\toprule
		Cell
		& \(\GTenToneHL{revenue}{R}(\GTenToneHL{alternative}{A_i},\GTenToneHL{state}{S_j})\)
		& \(\GTenToneHL{cost}{C}(\GTenToneHL{alternative}{A_i},\GTenToneHL{state}{S_j})\)
		& \(\GTenToneHL{payoff}{P}(\GTenToneHL{alternative}{A_i},\GTenToneHL{state}{S_j})=\GTenToneHL{revenue}{R}(\GTenToneHL{alternative}{A_i},\GTenToneHL{state}{S_j})-\GTenToneHL{cost}{C}(\GTenToneHL{alternative}{A_i},\GTenToneHL{state}{S_j})\) \\
		\midrule
		\((\GTenToneHL{alternative}{A_1},\GTenToneHL{state}{S_1})\) & \GTenToneHL{revenue}{58} & \GTenToneHL{cost}{16} & \(58-16=42\) \\
		\((\GTenToneHL{alternative}{A_1},\GTenToneHL{state}{S_2})\) & \GTenToneHL{revenue}{27} & \GTenToneHL{cost}{19} & \(27-19=8\) \\
		\((\GTenToneHL{alternative}{A_2},\GTenToneHL{state}{S_1})\) & \GTenToneHL{revenue}{49} & \GTenToneHL{cost}{21} & \(49-21=28\) \\
		\((\GTenToneHL{alternative}{A_2},\GTenToneHL{state}{S_2})\) & \GTenToneHL{revenue}{41} & \GTenToneHL{cost}{21} & \(41-21=20\) \\
		\bottomrule
	\end{tabular}
\end{center}

The resulting \GTenToneHL{payoff}{payoff} table is:

\begin{center}
	\begin{tabular}{lrr}
		\toprule
		Alternative & High traffic & Low traffic \\
		\midrule
		Launch dessert line & \GTenToneHL{payoff}{42} & \GTenToneHL{payoff}{8} \\
		Keep current menu & \GTenToneHL{payoff}{28} & \GTenToneHL{payoff}{20} \\
		\bottomrule
	\end{tabular}
\end{center}

\secInicio[sec:problem3-local-cafe-investment-c1-c3][sec:direct-revenue-cost-models]

% ----------------------------------------------------------------
\subsection{Green Energy Production Mix (C1--C3)}
\label{sec:problem4-green-energy-production-mix-c1-c3}

A renewable-energy company must select an electricity-generation mix.
The transaction is annual electricity sales, where \GTenToneHL{revenue}{revenue} is earned from
electricity generated and sold, while \GTenToneHL{cost}{costs} arise from operating the selected
generation system.

\begin{multicols}{2}
	The company is comparing three \GTenToneHL{alternative}{alternatives}:
	\begin{enumerate}[label=\arabic*.]
		\item \(\GTenToneHL{alternative}{A_1}\): Build wind turbines
		\item \(\GTenToneHL{alternative}{A_2}\): Build solar farms
		\item \(\GTenToneHL{alternative}{A_3}\): Build a balanced hybrid system
	\end{enumerate}
	
	Three weather states are possible:
	\begin{enumerate}[label=\arabic*.]
		\item \(\GTenToneHL{state}{S_1}\): Windy
		\item \(\GTenToneHL{state}{S_2}\): Sunny
		\item \(\GTenToneHL{state}{S_3}\): Cloudy
	\end{enumerate}
	
	\columnbreak
	
	The directly estimated revenues and \GTenToneHL{cost}{costs}, in thousand USD, are:
	\begin{enumerate}[label=\arabic*.]
		\item Wind turbines:
		\begin{enumerate}[label=\roman*)]
			\item Windy:
			\(\GTenToneHL{revenue}{R}(\GTenToneHL{alternative}{A_1},\GTenToneHL{state}{S_1})=\GTenToneHL{revenue}{132}\), \(\GTenToneHL{cost}{C}(\GTenToneHL{alternative}{A_1},\GTenToneHL{state}{S_1})=\GTenToneHL{cost}{42}\)
			
			\item Sunny:
			\(\GTenToneHL{revenue}{R}(\GTenToneHL{alternative}{A_1},\GTenToneHL{state}{S_2})=\GTenToneHL{revenue}{78}\), \(\GTenToneHL{cost}{C}(\GTenToneHL{alternative}{A_1},\GTenToneHL{state}{S_2})=\GTenToneHL{cost}{33}\)
			
			\item Cloudy:
			\(\GTenToneHL{revenue}{R}(\GTenToneHL{alternative}{A_1},\GTenToneHL{state}{S_3})=\GTenToneHL{revenue}{62}\), \(\GTenToneHL{cost}{C}(\GTenToneHL{alternative}{A_1},\GTenToneHL{state}{S_3})=\GTenToneHL{cost}{32}\)
		\end{enumerate}
		
		\item Solar farms:
		\begin{enumerate}[label=\roman*)]
			\item Windy:
			\(\GTenToneHL{revenue}{R}(\GTenToneHL{alternative}{A_2},\GTenToneHL{state}{S_1})=\GTenToneHL{revenue}{77}\), \(\GTenToneHL{cost}{C}(\GTenToneHL{alternative}{A_2},\GTenToneHL{state}{S_1})=\GTenToneHL{cost}{42}\)
			
			\item Sunny:
			\(\GTenToneHL{revenue}{R}(\GTenToneHL{alternative}{A_2},\GTenToneHL{state}{S_2})=\GTenToneHL{revenue}{143}\), \(\GTenToneHL{cost}{C}(\GTenToneHL{alternative}{A_2},\GTenToneHL{state}{S_2})=\GTenToneHL{cost}{48}\)
			
			\item Cloudy:
			\(\GTenToneHL{revenue}{R}(\GTenToneHL{alternative}{A_2},\GTenToneHL{state}{S_3})=\GTenToneHL{revenue}{63}\), \(\GTenToneHL{cost}{C}(\GTenToneHL{alternative}{A_2},\GTenToneHL{state}{S_3})=\GTenToneHL{cost}{38}\)
		\end{enumerate}
		
		\item Hybrid system:
		\begin{enumerate}[label=\roman*)]
			\item Windy:
			\(\GTenToneHL{revenue}{R}(\GTenToneHL{alternative}{A_3},\GTenToneHL{state}{S_1})=\GTenToneHL{revenue}{112}\), \(\GTenToneHL{cost}{C}(\GTenToneHL{alternative}{A_3},\GTenToneHL{state}{S_1})=\GTenToneHL{cost}{42}\)
			
			\item Sunny:
			\(\GTenToneHL{revenue}{R}(\GTenToneHL{alternative}{A_3},\GTenToneHL{state}{S_2})=\GTenToneHL{revenue}{107}\), \(\GTenToneHL{cost}{C}(\GTenToneHL{alternative}{A_3},\GTenToneHL{state}{S_2})=\GTenToneHL{cost}{42}\)
			
			\item Cloudy:
			\(\GTenToneHL{revenue}{R}(\GTenToneHL{alternative}{A_3},\GTenToneHL{state}{S_3})=\GTenToneHL{revenue}{96}\), \(\GTenToneHL{cost}{C}(\GTenToneHL{alternative}{A_3},\GTenToneHL{state}{S_3})=\GTenToneHL{cost}{36}\)
		\end{enumerate}
	\end{enumerate}
\end{multicols}

\textbf{Questions/Tasks.}
\begin{enumerate}
	\CriterionTask
	{subsec:problem4-green-energy-production-mix-c1}{1}
	{Interpreting decision alternatives, events, consequences, and states.}

	\CriterionTask
	{subsec:problem4-green-energy-production-mix-c3}{3}
	{Calculating profit and building the payoff table.}
\end{enumerate}

\subsubsection*{C1 - Interpreting decision alternatives, events, consequences, and states.}
\phantomsection
\label{subsec:problem4-green-energy-production-mix-c1}

\DecisionElementsTable
{Wind turbines; Solar farms; Hybrid system}
{Windy; Sunny; Cloudy}
{Weather outcome affecting electricity generation}
{Profit in thousand USD for each strategy--weather combination}

\secInicio[sec:problem4-green-energy-production-mix-c1-c3][sec:direct-revenue-cost-models]

\subsubsection*{C3 - Calculating profit and building the payoff table.}
\phantomsection
\label{subsec:problem4-green-energy-production-mix-c3}

For every \GTenToneHL{alternative}{alternative}--state pair, calculate
\[
\GTenToneHL{payoff}{P}(\GTenToneHL{alternative}{A_i},\GTenToneHL{state}{S_j})=\GTenToneHL{revenue}{R}(\GTenToneHL{alternative}{A_i},\GTenToneHL{state}{S_j})-\GTenToneHL{cost}{C}(\GTenToneHL{alternative}{A_i},\GTenToneHL{state}{S_j}).
\]

\begin{center}
	\small
	\begin{tabular}{lrrr}
		\toprule
		Cell
		& \(\GTenToneHL{revenue}{R}(\GTenToneHL{alternative}{A_i},\GTenToneHL{state}{S_j})\)
		& \(\GTenToneHL{cost}{C}(\GTenToneHL{alternative}{A_i},\GTenToneHL{state}{S_j})\)
		& \(\GTenToneHL{payoff}{P}(\GTenToneHL{alternative}{A_i},\GTenToneHL{state}{S_j})=\GTenToneHL{revenue}{R}(\GTenToneHL{alternative}{A_i},\GTenToneHL{state}{S_j})-\GTenToneHL{cost}{C}(\GTenToneHL{alternative}{A_i},\GTenToneHL{state}{S_j})\) \\
		\midrule
		\((\GTenToneHL{alternative}{A_1},\GTenToneHL{state}{S_1})\) & \GTenToneHL{revenue}{132} & \GTenToneHL{cost}{42} & \(132-42=90\) \\
		\((\GTenToneHL{alternative}{A_1},\GTenToneHL{state}{S_2})\) & \GTenToneHL{revenue}{78}  & \GTenToneHL{cost}{33} & \(78-33=45\) \\
		\((\GTenToneHL{alternative}{A_1},\GTenToneHL{state}{S_3})\) & \GTenToneHL{revenue}{62}  & \GTenToneHL{cost}{32} & \(62-32=30\) \\
		\((\GTenToneHL{alternative}{A_2},\GTenToneHL{state}{S_1})\) & \GTenToneHL{revenue}{77}  & \GTenToneHL{cost}{42} & \(77-42=35\) \\
		\((\GTenToneHL{alternative}{A_2},\GTenToneHL{state}{S_2})\) & \GTenToneHL{revenue}{143} & \GTenToneHL{cost}{48} & \(143-48=95\) \\
		\((\GTenToneHL{alternative}{A_2},\GTenToneHL{state}{S_3})\) & \GTenToneHL{revenue}{63}  & \GTenToneHL{cost}{38} & \(63-38=25\) \\
		\((\GTenToneHL{alternative}{A_3},\GTenToneHL{state}{S_1})\) & \GTenToneHL{revenue}{112} & \GTenToneHL{cost}{42} & \(112-42=70\) \\
		\((\GTenToneHL{alternative}{A_3},\GTenToneHL{state}{S_2})\) & \GTenToneHL{revenue}{107} & \GTenToneHL{cost}{42} & \(107-42=65\) \\
		\((\GTenToneHL{alternative}{A_3},\GTenToneHL{state}{S_3})\) & \GTenToneHL{revenue}{96}  & \GTenToneHL{cost}{36} & \(96-36=60\) \\
		\bottomrule
	\end{tabular}
\end{center}

The resulting \GTenToneHL{payoff}{payoff} table is:

\begin{center}
	\begin{tabular}{lrrr}
		\toprule
		Strategy & Windy & Sunny & Cloudy \\
		\midrule
		Wind turbines & \GTenToneHL{payoff}{90} & \GTenToneHL{payoff}{45} & \GTenToneHL{payoff}{30} \\
		Solar farms & \GTenToneHL{payoff}{35} & \GTenToneHL{payoff}{95} & \GTenToneHL{payoff}{25} \\
		Hybrid system & \GTenToneHL{payoff}{70} & \GTenToneHL{payoff}{65} & \GTenToneHL{payoff}{60} \\
		\bottomrule
	\end{tabular}
\end{center}

\secInicio[sec:problem4-green-energy-production-mix-c1-c3][sec:direct-revenue-cost-models]

% ================================================================
% ================================================================
\section{One-Component Models}
\label{sec:one-component-models}

\SectionProblemsTableOfContents{
	\SectionProblemLink{sec:problem5-custom-printing-product-line-c1-c3}{Custom Printing Product Line (C1--C3)} \\
	\SectionProblemLink{sec:problem6-event-equipment-fleet-c1-c3}{Event Equipment Fleet (C1--C3)} \\
	\SectionProblemLink{sec:problem7-community-arts-sponsorship-c1-c3}{Community Arts Sponsorship (C1--C3)} \\
	\SectionProblemLink{sec:problem8-distribution-contract-fulfilment-c1-c3}{Distribution Contract Fulfilment (C1--C3)}
}

\vspace{0.8cm}

In a one-component model, \GTenToneHL{revenue}{revenue} contains exactly one nonzero component
and \GTenToneHL{cost}{cost} contains exactly one nonzero component. Each component may be
variable or fixed and may depend on the \GTenToneHL{alternative}{alternative} or the state. The four paired patterns in this section are:

\begin{center}
	\small
	\begin{tabular}{cll}
		\toprule
		Problem & \GTenToneHL{revenue}{Revenue} model & \GTenToneHL{cost}{Cost} model \\
		\midrule
		5
		& \(\GTenToneHL{revenue}{R}(\GTenToneHL{alternative}{A_i},\GTenToneHL{state}{S_j})=\GTenToneHL{quantity}{q}(\GTenToneHL{alternative}{A_i},\GTenToneHL{state}{S_j})\GTenToneHL{revenue}{r}(\GTenToneHL{alternative}{A_i})\)
		& \(\GTenToneHL{cost}{C}(\GTenToneHL{alternative}{A_i},\GTenToneHL{state}{S_j})=\GTenToneHL{quantity}{q}(\GTenToneHL{alternative}{A_i},\GTenToneHL{state}{S_j})\GTenToneHL{cost}{c}(\GTenToneHL{state}{S_j})\) \\
		
		6
		& \(\GTenToneHL{revenue}{R}(\GTenToneHL{alternative}{A_i},\GTenToneHL{state}{S_j})=\GTenToneHL{quantity}{q}(\GTenToneHL{alternative}{A_i},\GTenToneHL{state}{S_j})\GTenToneHL{revenue}{r}(\GTenToneHL{state}{S_j})\)
		& \(\GTenToneHL{cost}{C}(\GTenToneHL{alternative}{A_i},\GTenToneHL{state}{S_j})=\GTenToneHL{cost}{f_c}(\GTenToneHL{alternative}{A_i})\) \\
		
		7
		& \(\GTenToneHL{revenue}{R}(\GTenToneHL{alternative}{A_i},\GTenToneHL{state}{S_j})=\GTenToneHL{revenue}{f_r}(\GTenToneHL{alternative}{A_i})\)
		& \(\GTenToneHL{cost}{C}(\GTenToneHL{alternative}{A_i},\GTenToneHL{state}{S_j})=\GTenToneHL{cost}{f_c}(\GTenToneHL{state}{S_j})\) \\
		
		8
		& \(\GTenToneHL{revenue}{R}(\GTenToneHL{alternative}{A_i},\GTenToneHL{state}{S_j})=\GTenToneHL{revenue}{f_r}(\GTenToneHL{state}{S_j})\)
		& \(\GTenToneHL{cost}{C}(\GTenToneHL{alternative}{A_i},\GTenToneHL{state}{S_j})=\GTenToneHL{quantity}{q}(\GTenToneHL{alternative}{A_i},\GTenToneHL{state}{S_j})\GTenToneHL{cost}{c}(\GTenToneHL{alternative}{A_i})\) \\
		\bottomrule
	\end{tabular}
\end{center}

Thus, the four one-component structures---AV, SV, AF, and SF---appear
exactly once as a \GTenToneHL{revenue}{revenue} structure and exactly once as a \GTenToneHL{cost}{cost} structure.

% ----------------------------------------------------------------
\subsection{Custom Printing Product Line (C1--C3)}
\label{sec:problem5-custom-printing-product-line-c1-c3}

A custom-printing cooperative must select the product line that it will
produce before market conditions are known. The transaction is the sale of
production lots. \GTenToneHL{revenue}{Revenue} per lot depends on the selected product, while \GTenToneHL{cost}{cost}
per lot depends on the market state.

\begin{multicols}{2}
	The cooperative is comparing two \GTenToneHL{alternative}{alternatives}:
	\begin{enumerate}[label=\arabic*.]
		\item \(\GTenToneHL{alternative}{A_1}\): Produce premium notebooks
		\item \(\GTenToneHL{alternative}{A_2}\): Produce canvas bags
	\end{enumerate}
	
	The possible \GTenToneHL{state}{states} \GTenToneHL{state}{of} \GTenToneHL{state}{nature} are:
	\begin{enumerate}[label=\arabic*.]
		\item \(\GTenToneHL{state}{S_1}\): Strong demand and stable material prices
		\item \(\GTenToneHL{state}{S_2}\): Weak demand and higher material prices
	\end{enumerate}
	
	\columnbreak
	
	The quantities, measured in production lots, are
	\[
	\begin{aligned}
		\GTenToneHL{quantity}{q}(\GTenToneHL{alternative}{A_1},\GTenToneHL{state}{S_1})&=\GTenToneHL{quantity}{6},
		&
		\GTenToneHL{quantity}{q}(\GTenToneHL{alternative}{A_1},\GTenToneHL{state}{S_2})&=\GTenToneHL{quantity}{4},
		\\
		\GTenToneHL{quantity}{q}(\GTenToneHL{alternative}{A_2},\GTenToneHL{state}{S_1})&=\GTenToneHL{quantity}{5},
		&
		\GTenToneHL{quantity}{q}(\GTenToneHL{alternative}{A_2},\GTenToneHL{state}{S_2})&=3.
	\end{aligned}
	\]
	
	The \GTenToneHL{alternative}{alternative}-based \GTenToneHL{revenue}{revenue} rates, in thousand USD per lot, are
	\[
	\GTenToneHL{revenue}{r}(\GTenToneHL{alternative}{A_1})=\GTenToneHL{revenue}{12},
	\qquad
	\GTenToneHL{revenue}{r}(\GTenToneHL{alternative}{A_2})=15.
	\]
	
	The state-based \GTenToneHL{cost}{cost} rates, in thousand USD per lot, are
	\[
	\GTenToneHL{cost}{c}(\GTenToneHL{state}{S_1})=\GTenToneHL{cost}{5},
	\qquad
	\GTenToneHL{cost}{c}(\GTenToneHL{state}{S_2})=8.
	\]
\end{multicols}

\textbf{Questions/Tasks.}
\begin{enumerate}
	\CriterionTask
	{subsec:problem5-custom-printing-product-line-c1}{1}
	{Interpreting decision alternatives, events, consequences, and states.}

	\CriterionTask
	{subsec:problem5-custom-printing-product-line-c3}{3}
	{Calculating profit and building the payoff table.}
\end{enumerate}

\subsubsection*{C1 - Interpreting decision alternatives, events, consequences, and states.}
\phantomsection
\label{subsec:problem5-custom-printing-product-line-c1}

\DecisionElementsTable
{Produce premium notebooks; Produce canvas bags}
{Strong demand and stable material prices; Weak demand and higher material prices}
{Market conditions revealed after the product-line decision}
{Profit in thousand USD for each product-line--market-state pair}

The only \GTenToneHL{revenue}{revenue} component is an \GTenToneHL{alternative}{alternative}-based variable component:
\[
\GTenToneHL{revenue}{R}(\GTenToneHL{alternative}{A_i},\GTenToneHL{state}{S_j})=\GTenToneHL{quantity}{q}(\GTenToneHL{alternative}{A_i},\GTenToneHL{state}{S_j})\GTenToneHL{revenue}{r}(\GTenToneHL{alternative}{A_i}).
\]
The only \GTenToneHL{cost}{cost} component is a state-based variable component:
\[
\GTenToneHL{cost}{C}(\GTenToneHL{alternative}{A_i},\GTenToneHL{state}{S_j})=\GTenToneHL{quantity}{q}(\GTenToneHL{alternative}{A_i},\GTenToneHL{state}{S_j})\GTenToneHL{cost}{c}(\GTenToneHL{state}{S_j}).
\]
Therefore,
\[
\GTenToneHL{payoff}{P}(\GTenToneHL{alternative}{A_i},\GTenToneHL{state}{S_j})
=
\GTenToneHL{quantity}{q}(\GTenToneHL{alternative}{A_i},\GTenToneHL{state}{S_j})\GTenToneHL{revenue}{r}(\GTenToneHL{alternative}{A_i})
-
\GTenToneHL{quantity}{q}(\GTenToneHL{alternative}{A_i},\GTenToneHL{state}{S_j})\GTenToneHL{cost}{c}(\GTenToneHL{state}{S_j}).
\]

\secInicio[sec:problem5-custom-printing-product-line-c1-c3][sec:one-component-models]

\subsubsection*{C3 - Calculating profit and building the payoff table.}
\phantomsection
\label{subsec:problem5-custom-printing-product-line-c3}

For every \GTenToneHL{alternative}{alternative}--state pair, calculate \GTenToneHL{revenue}{revenue} and \GTenToneHL{cost}{cost} independently:
\[
\begin{aligned}
	\GTenToneHL{revenue}{R}(\GTenToneHL{alternative}{A_i},\GTenToneHL{state}{S_j})&=\GTenToneHL{quantity}{q}(\GTenToneHL{alternative}{A_i},\GTenToneHL{state}{S_j})\GTenToneHL{revenue}{r}(\GTenToneHL{alternative}{A_i}),
	\\
	\GTenToneHL{cost}{C}(\GTenToneHL{alternative}{A_i},\GTenToneHL{state}{S_j})&=\GTenToneHL{quantity}{q}(\GTenToneHL{alternative}{A_i},\GTenToneHL{state}{S_j})\GTenToneHL{cost}{c}(\GTenToneHL{state}{S_j}),
	\\
	\GTenToneHL{payoff}{P}(\GTenToneHL{alternative}{A_i},\GTenToneHL{state}{S_j})&=\GTenToneHL{revenue}{R}(\GTenToneHL{alternative}{A_i},\GTenToneHL{state}{S_j})-\GTenToneHL{cost}{C}(\GTenToneHL{alternative}{A_i},\GTenToneHL{state}{S_j}).
\end{aligned}
\]

\begin{center}
	\small
	\begin{tabular}{llll}
		\toprule
		Cell
		& \(\GTenToneHL{revenue}{R}(\GTenToneHL{alternative}{A_i},\GTenToneHL{state}{S_j})\)
		& \(\GTenToneHL{cost}{C}(\GTenToneHL{alternative}{A_i},\GTenToneHL{state}{S_j})\)
		& \(\GTenToneHL{payoff}{P}(\GTenToneHL{alternative}{A_i},\GTenToneHL{state}{S_j})=\GTenToneHL{revenue}{R}(\GTenToneHL{alternative}{A_i},\GTenToneHL{state}{S_j})-\GTenToneHL{cost}{C}(\GTenToneHL{alternative}{A_i},\GTenToneHL{state}{S_j})\) \\
		\midrule
		\((\GTenToneHL{alternative}{A_1},\GTenToneHL{state}{S_1})\)
		& \(6(12)=72\)
		& \(6(5)=30\)
		& \(72-30=42\) \\
		
		\((\GTenToneHL{alternative}{A_1},\GTenToneHL{state}{S_2})\)
		& \(4(12)=48\)
		& \(4(8)=32\)
		& \(48-32=16\) \\
		
		\((\GTenToneHL{alternative}{A_2},\GTenToneHL{state}{S_1})\)
		& \(5(15)=75\)
		& \(5(5)=25\)
		& \(75-25=50\) \\
		
		\((\GTenToneHL{alternative}{A_2},\GTenToneHL{state}{S_2})\)
		& \(3(15)=45\)
		& \(3(8)=24\)
		& \(45-24=21\) \\
		\bottomrule
	\end{tabular}
\end{center}

The resulting \GTenToneHL{payoff}{payoff} table is:

\begin{center}
	\begin{tabular}{lrr}
		\toprule
		Alternative
		& Strong demand and stable prices
		& Weak demand and higher prices \\
		\midrule
		Premium notebooks & \GTenToneHL{payoff}{42} & \GTenToneHL{payoff}{16} \\
		Canvas bags & \GTenToneHL{payoff}{50} & \GTenToneHL{payoff}{21} \\
		\bottomrule
	\end{tabular}
\end{center}

\secInicio[sec:problem5-custom-printing-product-line-c1-c3][sec:one-component-models]

% ----------------------------------------------------------------
\subsection{Event Equipment Fleet (C1--C3)}
\label{sec:problem6-event-equipment-fleet-c1-c3}

An event-equipment provider must choose the size of the fleet that it will
prepare before rental demand is known. \GTenToneHL{revenue}{Revenue} per rental contract depends on
the demand state, while the preparation \GTenToneHL{cost}{cost} is a fixed amount determined by
the selected fleet.

\begin{multicols}{2}
	The provider is comparing two \GTenToneHL{alternative}{alternatives}:
	\begin{enumerate}[label=\arabic*.]
		\item \(\GTenToneHL{alternative}{A_1}\): Prepare a compact fleet
		\item \(\GTenToneHL{alternative}{A_2}\): Prepare an expanded fleet
	\end{enumerate}
	
	The possible \GTenToneHL{state}{states} \GTenToneHL{state}{of} \GTenToneHL{state}{nature} are:
	\begin{enumerate}[label=\arabic*.]
		\item \(\GTenToneHL{state}{S_1}\): High rental demand
		\item \(\GTenToneHL{state}{S_2}\): Low rental demand
	\end{enumerate}
	
	\columnbreak
	
	The numbers of rental contracts are
	\[
	\begin{aligned}
		\GTenToneHL{quantity}{q}(\GTenToneHL{alternative}{A_1},\GTenToneHL{state}{S_1})&=\GTenToneHL{quantity}{6},
		&
		\GTenToneHL{quantity}{q}(\GTenToneHL{alternative}{A_1},\GTenToneHL{state}{S_2})&=\GTenToneHL{quantity}{3},
		\\
		\GTenToneHL{quantity}{q}(\GTenToneHL{alternative}{A_2},\GTenToneHL{state}{S_1})&=\GTenToneHL{quantity}{8},
		&
		\GTenToneHL{quantity}{q}(\GTenToneHL{alternative}{A_2},\GTenToneHL{state}{S_2})&=4.
	\end{aligned}
	\]
	
	The state-based \GTenToneHL{revenue}{revenue} rates, in thousand USD per contract, are
	\[
	\GTenToneHL{revenue}{r}(\GTenToneHL{state}{S_1})=\GTenToneHL{revenue}{9},
	\qquad
	\GTenToneHL{revenue}{r}(\GTenToneHL{state}{S_2})=6.
	\]
	
	The \GTenToneHL{alternative}{alternative}-based fixed \GTenToneHL{cost}{costs}, in thousand USD, are
	\[
	\GTenToneHL{cost}{f_c}(\GTenToneHL{alternative}{A_1})=\GTenToneHL{cost}{22},
	\qquad
	\GTenToneHL{cost}{f_c}(\GTenToneHL{alternative}{A_2})=34.
	\]
\end{multicols}

\textbf{Questions/Tasks.}
\begin{enumerate}
	\CriterionTask
	{subsec:problem6-event-equipment-fleet-c1}{1}
	{Interpreting decision alternatives, events, consequences, and states.}

	\CriterionTask
	{subsec:problem6-event-equipment-fleet-c3}{3}
	{Calculating profit and building the payoff table.}
\end{enumerate}

\subsubsection*{C1 - Interpreting decision alternatives, events, consequences, and states.}
\phantomsection
\label{subsec:problem6-event-equipment-fleet-c1}

\DecisionElementsTable
{Prepare a compact fleet; Prepare an expanded fleet}
{High rental demand; Low rental demand}
{Rental demand revealed after the fleet-size decision}
{Profit in thousand USD for each fleet--demand pair}

The only \GTenToneHL{revenue}{revenue} component is a state-based variable component:
\[
\GTenToneHL{revenue}{R}(\GTenToneHL{alternative}{A_i},\GTenToneHL{state}{S_j})=\GTenToneHL{quantity}{q}(\GTenToneHL{alternative}{A_i},\GTenToneHL{state}{S_j})\GTenToneHL{revenue}{r}(\GTenToneHL{state}{S_j}).
\]
The only \GTenToneHL{cost}{cost} component is an \GTenToneHL{alternative}{alternative}-based fixed component:
\[
\GTenToneHL{cost}{C}(\GTenToneHL{alternative}{A_i},\GTenToneHL{state}{S_j})=\GTenToneHL{cost}{f_c}(\GTenToneHL{alternative}{A_i}).
\]
Therefore,
\[
\GTenToneHL{payoff}{P}(\GTenToneHL{alternative}{A_i},\GTenToneHL{state}{S_j})
=
\GTenToneHL{quantity}{q}(\GTenToneHL{alternative}{A_i},\GTenToneHL{state}{S_j})\GTenToneHL{revenue}{r}(\GTenToneHL{state}{S_j})
-
\GTenToneHL{cost}{f_c}(\GTenToneHL{alternative}{A_i}).
\]

\secInicio[sec:problem6-event-equipment-fleet-c1-c3][sec:one-component-models]

\subsubsection*{C3 - Calculating profit and building the payoff table.}
\phantomsection
\label{subsec:problem6-event-equipment-fleet-c3}

For every \GTenToneHL{alternative}{alternative}--state pair, calculate
\[
\begin{aligned}
	\GTenToneHL{revenue}{R}(\GTenToneHL{alternative}{A_i},\GTenToneHL{state}{S_j})&=\GTenToneHL{quantity}{q}(\GTenToneHL{alternative}{A_i},\GTenToneHL{state}{S_j})\GTenToneHL{revenue}{r}(\GTenToneHL{state}{S_j}),
	\\
	\GTenToneHL{cost}{C}(\GTenToneHL{alternative}{A_i},\GTenToneHL{state}{S_j})&=\GTenToneHL{cost}{f_c}(\GTenToneHL{alternative}{A_i}),
	\\
	\GTenToneHL{payoff}{P}(\GTenToneHL{alternative}{A_i},\GTenToneHL{state}{S_j})&=\GTenToneHL{revenue}{R}(\GTenToneHL{alternative}{A_i},\GTenToneHL{state}{S_j})-\GTenToneHL{cost}{C}(\GTenToneHL{alternative}{A_i},\GTenToneHL{state}{S_j}).
\end{aligned}
\]

\begin{center}
	\small
	\begin{tabular}{llll}
		\toprule
		Cell
		& \(\GTenToneHL{revenue}{R}(\GTenToneHL{alternative}{A_i},\GTenToneHL{state}{S_j})\)
		& \(\GTenToneHL{cost}{C}(\GTenToneHL{alternative}{A_i},\GTenToneHL{state}{S_j})\)
		& \(\GTenToneHL{payoff}{P}(\GTenToneHL{alternative}{A_i},\GTenToneHL{state}{S_j})=\GTenToneHL{revenue}{R}(\GTenToneHL{alternative}{A_i},\GTenToneHL{state}{S_j})-\GTenToneHL{cost}{C}(\GTenToneHL{alternative}{A_i},\GTenToneHL{state}{S_j})\) \\
		\midrule
		\((\GTenToneHL{alternative}{A_1},\GTenToneHL{state}{S_1})\)
		& \(6(9)=54\)
		& \(22\)
		& \(54-22=32\) \\
		
		\((\GTenToneHL{alternative}{A_1},\GTenToneHL{state}{S_2})\)
		& \(3(6)=18\)
		& \(22\)
		& \(18-22=-4\) \\
		
		\((\GTenToneHL{alternative}{A_2},\GTenToneHL{state}{S_1})\)
		& \(8(9)=72\)
		& \(34\)
		& \(72-34=38\) \\
		
		\((\GTenToneHL{alternative}{A_2},\GTenToneHL{state}{S_2})\)
		& \(4(6)=24\)
		& \(34\)
		& \(24-34=-10\) \\
		\bottomrule
	\end{tabular}
\end{center}

The resulting \GTenToneHL{payoff}{payoff} table is:

\begin{center}
	\begin{tabular}{lrr}
		\toprule
		Alternative & High rental demand & Low rental demand \\
		\midrule
		Compact fleet & \GTenToneHL{payoff}{32} & \(-4\) \\
		Expanded fleet & \GTenToneHL{payoff}{38} & \(-10\) \\
		\bottomrule
	\end{tabular}
\end{center}

The negative \GTenToneHL{payoff}{payoffs} in the low-demand state represent losses because the
fixed preparation \GTenToneHL{cost}{cost} is greater than the \GTenToneHL{revenue}{revenue} earned.

\secInicio[sec:problem6-event-equipment-fleet-c1-c3][sec:one-component-models]

% ----------------------------------------------------------------
\subsection{Community Arts Sponsorship (C1--C3)}
\label{sec:problem7-community-arts-sponsorship-c1-c3}

A community arts organisation must select a sponsorship agreement before the
event requirements are confirmed. Each agreement provides a fixed sponsorship
payment, while the event incurs a fixed regulatory and safety \GTenToneHL{cost}{cost} determined
by the \GTenToneHL{state}{state} \GTenToneHL{state}{of} \GTenToneHL{state}{nature}.

\begin{multicols}{2}
	The organisation is comparing two \GTenToneHL{alternative}{alternatives}:
	\begin{enumerate}[label=\arabic*.]
		\item \(\GTenToneHL{alternative}{A_1}\): Accept the standard sponsorship agreement
		\item \(\GTenToneHL{alternative}{A_2}\): Accept the premium sponsorship agreement
	\end{enumerate}
	
	The possible \GTenToneHL{state}{states} \GTenToneHL{state}{of} \GTenToneHL{state}{nature} are:
	\begin{enumerate}[label=\arabic*.]
		\item \(\GTenToneHL{state}{S_1}\): Routine permit requirements
		\item \(\GTenToneHL{state}{S_2}\): Enhanced safety requirements
	\end{enumerate}
	
	\columnbreak
	
	The \GTenToneHL{alternative}{alternative}-based fixed revenues, in thousand USD, are
	\[
	\GTenToneHL{revenue}{f_r}(\GTenToneHL{alternative}{A_1})=\GTenToneHL{revenue}{52},
	\qquad
	\GTenToneHL{revenue}{f_r}(\GTenToneHL{alternative}{A_2})=68.
	\]
	
	The state-based fixed \GTenToneHL{cost}{costs}, in thousand USD, are
	\[
	\GTenToneHL{cost}{f_c}(\GTenToneHL{state}{S_1})=\GTenToneHL{cost}{18},
	\qquad
	\GTenToneHL{cost}{f_c}(\GTenToneHL{state}{S_2})=31.
	\]
	
	Neither model contains a variable component, so \GTenToneHL{quantity}{quantity} is not required.
\end{multicols}

\textbf{Questions/Tasks.}
\begin{enumerate}
	\CriterionTask
	{subsec:problem7-community-arts-sponsorship-c1}{1}
	{Interpreting decision alternatives, events, consequences, and states.}

	\CriterionTask
	{subsec:problem7-community-arts-sponsorship-c3}{3}
	{Calculating profit and building the payoff table.}
\end{enumerate}

\subsubsection*{C1 - Interpreting decision alternatives, events, consequences, and states.}
\phantomsection
\label{subsec:problem7-community-arts-sponsorship-c1}

\DecisionElementsTable
{Accept the standard sponsorship agreement; Accept the premium sponsorship agreement}
{Routine permit requirements; Enhanced safety requirements}
{Event requirements confirmed after the sponsorship decision}
{Profit in thousand USD for each sponsorship--requirement pair}

The only \GTenToneHL{revenue}{revenue} component is an \GTenToneHL{alternative}{alternative}-based fixed component:
\[
\GTenToneHL{revenue}{R}(\GTenToneHL{alternative}{A_i},\GTenToneHL{state}{S_j})=\GTenToneHL{revenue}{f_r}(\GTenToneHL{alternative}{A_i}).
\]
The only \GTenToneHL{cost}{cost} component is a state-based fixed component:
\[
\GTenToneHL{cost}{C}(\GTenToneHL{alternative}{A_i},\GTenToneHL{state}{S_j})=\GTenToneHL{cost}{f_c}(\GTenToneHL{state}{S_j}).
\]
Therefore,
\[
\GTenToneHL{payoff}{P}(\GTenToneHL{alternative}{A_i},\GTenToneHL{state}{S_j})
=
\GTenToneHL{revenue}{f_r}(\GTenToneHL{alternative}{A_i})
-
\GTenToneHL{cost}{f_c}(\GTenToneHL{state}{S_j}).
\]

\secInicio[sec:problem7-community-arts-sponsorship-c1-c3][sec:one-component-models]

\subsubsection*{C3 - Calculating profit and building the payoff table.}
\phantomsection
\label{subsec:problem7-community-arts-sponsorship-c3}

For every \GTenToneHL{alternative}{alternative}--state pair, calculate
\[
\begin{aligned}
	\GTenToneHL{revenue}{R}(\GTenToneHL{alternative}{A_i},\GTenToneHL{state}{S_j})&=\GTenToneHL{revenue}{f_r}(\GTenToneHL{alternative}{A_i}),
	\\
	\GTenToneHL{cost}{C}(\GTenToneHL{alternative}{A_i},\GTenToneHL{state}{S_j})&=\GTenToneHL{cost}{f_c}(\GTenToneHL{state}{S_j}),
	\\
	\GTenToneHL{payoff}{P}(\GTenToneHL{alternative}{A_i},\GTenToneHL{state}{S_j})&=\GTenToneHL{revenue}{R}(\GTenToneHL{alternative}{A_i},\GTenToneHL{state}{S_j})-\GTenToneHL{cost}{C}(\GTenToneHL{alternative}{A_i},\GTenToneHL{state}{S_j}).
\end{aligned}
\]

\begin{center}
	\small
	\begin{tabular}{lrrr}
		\toprule
		Cell
		& \(\GTenToneHL{revenue}{R}(\GTenToneHL{alternative}{A_i},\GTenToneHL{state}{S_j})\)
		& \(\GTenToneHL{cost}{C}(\GTenToneHL{alternative}{A_i},\GTenToneHL{state}{S_j})\)
		& \(\GTenToneHL{payoff}{P}(\GTenToneHL{alternative}{A_i},\GTenToneHL{state}{S_j})=\GTenToneHL{revenue}{R}(\GTenToneHL{alternative}{A_i},\GTenToneHL{state}{S_j})-\GTenToneHL{cost}{C}(\GTenToneHL{alternative}{A_i},\GTenToneHL{state}{S_j})\) \\
		\midrule
		\((\GTenToneHL{alternative}{A_1},\GTenToneHL{state}{S_1})\) & \GTenToneHL{revenue}{52} & \GTenToneHL{cost}{18} & \(52-18=34\) \\
		\((\GTenToneHL{alternative}{A_1},\GTenToneHL{state}{S_2})\) & \GTenToneHL{revenue}{52} & \GTenToneHL{cost}{31} & \(52-31=21\) \\
		\((\GTenToneHL{alternative}{A_2},\GTenToneHL{state}{S_1})\) & \GTenToneHL{revenue}{68} & \GTenToneHL{cost}{18} & \(68-18=50\) \\
		\((\GTenToneHL{alternative}{A_2},\GTenToneHL{state}{S_2})\) & \GTenToneHL{revenue}{68} & \GTenToneHL{cost}{31} & \(68-31=37\) \\
		\bottomrule
	\end{tabular}
\end{center}

The resulting \GTenToneHL{payoff}{payoff} table is:

\begin{center}
	\begin{tabular}{lrr}
		\toprule
		Alternative
		& Routine requirements
		& Enhanced requirements \\
		\midrule
		Standard sponsorship & \GTenToneHL{payoff}{34} & \GTenToneHL{payoff}{21} \\
		Premium sponsorship & \GTenToneHL{payoff}{50} & \GTenToneHL{payoff}{37} \\
		\bottomrule
	\end{tabular}
\end{center}

\secInicio[sec:problem7-community-arts-sponsorship-c1-c3][sec:one-component-models]

% ----------------------------------------------------------------
\subsection{Distribution Contract Fulfilment (C1--C3)}
\label{sec:problem8-distribution-contract-fulfilment-c1-c3}

A distribution company must choose a fulfilment method before the volume of a
new contract is known. The contract provides a fixed payment determined by the
realised volume state, while the processing \GTenToneHL{cost}{cost} per unit depends on the
selected fulfilment method.

\begin{multicols}{2}
	The company is comparing two \GTenToneHL{alternative}{alternatives}:
	\begin{enumerate}[label=\arabic*.]
		\item \(\GTenToneHL{alternative}{A_1}\): Use manual packing
		\item \(\GTenToneHL{alternative}{A_2}\): Use automated packing
	\end{enumerate}
	
	The possible \GTenToneHL{state}{states} \GTenToneHL{state}{of} \GTenToneHL{state}{nature} are:
	\begin{enumerate}[label=\arabic*.]
		\item \(\GTenToneHL{state}{S_1}\): Standard-volume contract
		\item \(\GTenToneHL{state}{S_2}\): High-volume contract
	\end{enumerate}
	
	\columnbreak
	
	The quantities, measured in hundreds of processed orders, are
	\[
	\begin{aligned}
		\GTenToneHL{quantity}{q}(\GTenToneHL{alternative}{A_1},\GTenToneHL{state}{S_1})&=\GTenToneHL{quantity}{4},
		&
		\GTenToneHL{quantity}{q}(\GTenToneHL{alternative}{A_1},\GTenToneHL{state}{S_2})&=\GTenToneHL{quantity}{7},
		\\
		\GTenToneHL{quantity}{q}(\GTenToneHL{alternative}{A_2},\GTenToneHL{state}{S_1})&=\GTenToneHL{quantity}{5},
		&
		\GTenToneHL{quantity}{q}(\GTenToneHL{alternative}{A_2},\GTenToneHL{state}{S_2})&=8.
	\end{aligned}
	\]
	
	The state-based fixed revenues, in thousand USD, are
	\[
	\GTenToneHL{revenue}{f_r}(\GTenToneHL{state}{S_1})=\GTenToneHL{revenue}{70},
	\qquad
	\GTenToneHL{revenue}{f_r}(\GTenToneHL{state}{S_2})=104.
	\]
	
	The \GTenToneHL{alternative}{alternative}-based \GTenToneHL{cost}{cost} rates, in thousand USD per hundred orders, are
	\[
	\GTenToneHL{cost}{c}(\GTenToneHL{alternative}{A_1})=\GTenToneHL{cost}{7},
	\qquad
	\GTenToneHL{cost}{c}(\GTenToneHL{alternative}{A_2})=5.
	\]
\end{multicols}

\textbf{Questions/Tasks.}
\begin{enumerate}
	\CriterionTask
	{subsec:problem8-distribution-contract-fulfilment-c1}{1}
	{Interpreting decision alternatives, events, consequences, and states.}

	\CriterionTask
	{subsec:problem8-distribution-contract-fulfilment-c3}{3}
	{Calculating profit and building the payoff table.}
\end{enumerate}

\subsubsection*{C1 - Interpreting decision alternatives, events, consequences, and states.}
\phantomsection
\label{subsec:problem8-distribution-contract-fulfilment-c1}

\DecisionElementsTable
{Use manual packing; Use automated packing}
{Standard-volume contract; High-volume contract}
{Contract volume revealed after the fulfilment-method decision}
{Profit in thousand USD for each fulfilment-method--volume pair}

The only \GTenToneHL{revenue}{revenue} component is a state-based fixed component:
\[
\GTenToneHL{revenue}{R}(\GTenToneHL{alternative}{A_i},\GTenToneHL{state}{S_j})=\GTenToneHL{revenue}{f_r}(\GTenToneHL{state}{S_j}).
\]
The only \GTenToneHL{cost}{cost} component is an \GTenToneHL{alternative}{alternative}-based variable component:
\[
\GTenToneHL{cost}{C}(\GTenToneHL{alternative}{A_i},\GTenToneHL{state}{S_j})=\GTenToneHL{quantity}{q}(\GTenToneHL{alternative}{A_i},\GTenToneHL{state}{S_j})\GTenToneHL{cost}{c}(\GTenToneHL{alternative}{A_i}).
\]
Therefore,
\[
\GTenToneHL{payoff}{P}(\GTenToneHL{alternative}{A_i},\GTenToneHL{state}{S_j})
=
\GTenToneHL{revenue}{f_r}(\GTenToneHL{state}{S_j})
-
\GTenToneHL{quantity}{q}(\GTenToneHL{alternative}{A_i},\GTenToneHL{state}{S_j})\GTenToneHL{cost}{c}(\GTenToneHL{alternative}{A_i}).
\]

\secInicio[sec:problem8-distribution-contract-fulfilment-c1-c3][sec:one-component-models]

\subsubsection*{C3 - Calculating profit and building the payoff table.}
\phantomsection
\label{subsec:problem8-distribution-contract-fulfilment-c3}

For every \GTenToneHL{alternative}{alternative}--state pair, calculate
\[
\begin{aligned}
	\GTenToneHL{revenue}{R}(\GTenToneHL{alternative}{A_i},\GTenToneHL{state}{S_j})&=\GTenToneHL{revenue}{f_r}(\GTenToneHL{state}{S_j}),
	\\
	\GTenToneHL{cost}{C}(\GTenToneHL{alternative}{A_i},\GTenToneHL{state}{S_j})&=\GTenToneHL{quantity}{q}(\GTenToneHL{alternative}{A_i},\GTenToneHL{state}{S_j})\GTenToneHL{cost}{c}(\GTenToneHL{alternative}{A_i}),
	\\
	\GTenToneHL{payoff}{P}(\GTenToneHL{alternative}{A_i},\GTenToneHL{state}{S_j})&=\GTenToneHL{revenue}{R}(\GTenToneHL{alternative}{A_i},\GTenToneHL{state}{S_j})-\GTenToneHL{cost}{C}(\GTenToneHL{alternative}{A_i},\GTenToneHL{state}{S_j}).
\end{aligned}
\]

\begin{center}
	\small
	\begin{tabular}{llll}
		\toprule
		Cell
		& \(\GTenToneHL{revenue}{R}(\GTenToneHL{alternative}{A_i},\GTenToneHL{state}{S_j})\)
		& \(\GTenToneHL{cost}{C}(\GTenToneHL{alternative}{A_i},\GTenToneHL{state}{S_j})\)
		& \(\GTenToneHL{payoff}{P}(\GTenToneHL{alternative}{A_i},\GTenToneHL{state}{S_j})=\GTenToneHL{revenue}{R}(\GTenToneHL{alternative}{A_i},\GTenToneHL{state}{S_j})-\GTenToneHL{cost}{C}(\GTenToneHL{alternative}{A_i},\GTenToneHL{state}{S_j})\) \\
		\midrule
		\((\GTenToneHL{alternative}{A_1},\GTenToneHL{state}{S_1})\)
		& \(70\)
		& \(4(7)=28\)
		& \(70-28=42\) \\
		
		\((\GTenToneHL{alternative}{A_1},\GTenToneHL{state}{S_2})\)
		& \(104\)
		& \(7(7)=49\)
		& \(104-49=55\) \\
		
		\((\GTenToneHL{alternative}{A_2},\GTenToneHL{state}{S_1})\)
		& \(70\)
		& \(5(5)=25\)
		& \(70-25=45\) \\
		
		\((\GTenToneHL{alternative}{A_2},\GTenToneHL{state}{S_2})\)
		& \(104\)
		& \(8(5)=40\)
		& \(104-40=64\) \\
		\bottomrule
	\end{tabular}
\end{center}

The resulting \GTenToneHL{payoff}{payoff} table is:

\begin{center}
	\begin{tabular}{lrr}
		\toprule
		Alternative
		& Standard-volume contract
		& High-volume contract \\
		\midrule
		Manual packing & \GTenToneHL{payoff}{42} & \GTenToneHL{payoff}{55} \\
		Automated packing & \GTenToneHL{payoff}{45} & \GTenToneHL{payoff}{64} \\
		\bottomrule
	\end{tabular}
\end{center}

\secInicio[sec:problem8-distribution-contract-fulfilment-c1-c3][sec:one-component-models]

% ================================================================
\section{Two-Component Models}
\label{sec:two-component-models}

\SectionProblemsTableOfContents{
	\SectionProblemLink{sec:problem9-farm-produce-subscription-c1-c3}{Farm Produce Subscription Plans (C1--C3)} \\
	\SectionProblemLink{sec:problem10-mobile-catering-configuration-c1-c3}{Mobile Catering Configuration (C1--C3)} \\
	\SectionProblemLink{sec:problem11-online-training-programme-c1-c3}{Online Training Programme (C1--C3)} \\
	\SectionProblemLink{sec:problem12-community-venue-membership-c1-c3}{Community Venue Membership (C1--C3)} \\
	\SectionProblemLink{sec:problem13-urban-delivery-fleet-c1-c3}{Urban Delivery Fleet (C1--C3)} \\
	\SectionProblemLink{sec:problem14-data-service-contract-c1-c3}{Data-Service Contract (C1--C3)}
}

\vspace{0.8cm}

In a two-component model, \GTenToneHL{revenue}{revenue} contains exactly two nonzero components
and \GTenToneHL{cost}{cost} contains exactly two nonzero components. Each model is constructed
independently, so the \GTenToneHL{revenue}{revenue} and \GTenToneHL{cost}{cost} structures do not need to match.

The six paired patterns in this section are:

\begin{center}
	\small
	\begin{tabular}{cll}
		\toprule
		Problem & \GTenToneHL{revenue}{Revenue} model & \GTenToneHL{cost}{Cost} model \\
		\midrule
		9
		& \(\GTenToneHL{revenue}{R}(\GTenToneHL{alternative}{A_i},\GTenToneHL{state}{S_j})=\GTenToneHL{quantity}{q}(\GTenToneHL{alternative}{A_i},\GTenToneHL{state}{S_j})\bigl[\GTenToneHL{revenue}{r}(\GTenToneHL{alternative}{A_i})+\GTenToneHL{revenue}{r}(\GTenToneHL{state}{S_j})\bigr]\)
		& \(\GTenToneHL{cost}{C}(\GTenToneHL{alternative}{A_i},\GTenToneHL{state}{S_j})=\GTenToneHL{quantity}{q}(\GTenToneHL{alternative}{A_i},\GTenToneHL{state}{S_j})\GTenToneHL{cost}{c}(\GTenToneHL{alternative}{A_i})+\GTenToneHL{cost}{f_c}(\GTenToneHL{state}{S_j})\) \\
		
		10
		& \(\GTenToneHL{revenue}{R}(\GTenToneHL{alternative}{A_i},\GTenToneHL{state}{S_j})=\GTenToneHL{quantity}{q}(\GTenToneHL{alternative}{A_i},\GTenToneHL{state}{S_j})\GTenToneHL{revenue}{r}(\GTenToneHL{alternative}{A_i})+\GTenToneHL{revenue}{f_r}(\GTenToneHL{alternative}{A_i})\)
		& \(\GTenToneHL{cost}{C}(\GTenToneHL{alternative}{A_i},\GTenToneHL{state}{S_j})=\GTenToneHL{quantity}{q}(\GTenToneHL{alternative}{A_i},\GTenToneHL{state}{S_j})\GTenToneHL{cost}{c}(\GTenToneHL{state}{S_j})+\GTenToneHL{cost}{f_c}(\GTenToneHL{alternative}{A_i})\) \\
		
		11
		& \(\GTenToneHL{revenue}{R}(\GTenToneHL{alternative}{A_i},\GTenToneHL{state}{S_j})=\GTenToneHL{quantity}{q}(\GTenToneHL{alternative}{A_i},\GTenToneHL{state}{S_j})\GTenToneHL{revenue}{r}(\GTenToneHL{alternative}{A_i})+\GTenToneHL{revenue}{f_r}(\GTenToneHL{state}{S_j})\)
		& \(\GTenToneHL{cost}{C}(\GTenToneHL{alternative}{A_i},\GTenToneHL{state}{S_j})=\GTenToneHL{quantity}{q}(\GTenToneHL{alternative}{A_i},\GTenToneHL{state}{S_j})\GTenToneHL{cost}{c}(\GTenToneHL{state}{S_j})+\GTenToneHL{cost}{f_c}(\GTenToneHL{state}{S_j})\) \\
		
		12
		& \(\GTenToneHL{revenue}{R}(\GTenToneHL{alternative}{A_i},\GTenToneHL{state}{S_j})=\GTenToneHL{quantity}{q}(\GTenToneHL{alternative}{A_i},\GTenToneHL{state}{S_j})\GTenToneHL{revenue}{r}(\GTenToneHL{state}{S_j})+\GTenToneHL{revenue}{f_r}(\GTenToneHL{alternative}{A_i})\)
		& \(\GTenToneHL{cost}{C}(\GTenToneHL{alternative}{A_i},\GTenToneHL{state}{S_j})=\GTenToneHL{cost}{f_c}(\GTenToneHL{alternative}{A_i})+\GTenToneHL{cost}{f_c}(\GTenToneHL{state}{S_j})\) \\
		
		13
		& \(\GTenToneHL{revenue}{R}(\GTenToneHL{alternative}{A_i},\GTenToneHL{state}{S_j})=\GTenToneHL{quantity}{q}(\GTenToneHL{alternative}{A_i},\GTenToneHL{state}{S_j})\GTenToneHL{revenue}{r}(\GTenToneHL{state}{S_j})+\GTenToneHL{revenue}{f_r}(\GTenToneHL{state}{S_j})\)
		& \(\GTenToneHL{cost}{C}(\GTenToneHL{alternative}{A_i},\GTenToneHL{state}{S_j})=\GTenToneHL{quantity}{q}(\GTenToneHL{alternative}{A_i},\GTenToneHL{state}{S_j})\bigl[\GTenToneHL{cost}{c}(\GTenToneHL{alternative}{A_i})+\GTenToneHL{cost}{c}(\GTenToneHL{state}{S_j})\bigr]\) \\
		
		14
		& \(\GTenToneHL{revenue}{R}(\GTenToneHL{alternative}{A_i},\GTenToneHL{state}{S_j})=\GTenToneHL{revenue}{f_r}(\GTenToneHL{alternative}{A_i})+\GTenToneHL{revenue}{f_r}(\GTenToneHL{state}{S_j})\)
		& \(\GTenToneHL{cost}{C}(\GTenToneHL{alternative}{A_i},\GTenToneHL{state}{S_j})=\GTenToneHL{quantity}{q}(\GTenToneHL{alternative}{A_i},\GTenToneHL{state}{S_j})\GTenToneHL{cost}{c}(\GTenToneHL{alternative}{A_i})+\GTenToneHL{cost}{f_c}(\GTenToneHL{alternative}{A_i})\) \\
		\bottomrule
	\end{tabular}
\end{center}

Thus, all six possible two-component structures---AV+SV, AV+AF,
AV+SF, SV+AF, SV+SF, and AF+SF---appear exactly once as a \GTenToneHL{revenue}{revenue}
structure and exactly once as a \GTenToneHL{cost}{cost} structure.

% ----------------------------------------------------------------
\subsection{Farm Produce Subscription Plans (C1--C3)}
\label{sec:problem9-farm-produce-subscription-c1-c3}

A farm cooperative must choose a produce-subscription plan before seasonal
market conditions are known. \GTenToneHL{revenue}{Revenue} per subscription consists of a rate
determined by the selected plan and an additional rate determined by the
market state. Costs consist of a plan-based variable \GTenToneHL{cost}{cost} and a fixed
distribution \GTenToneHL{cost}{cost} determined by the state.

\begin{multicols}{2}
	The cooperative is comparing two \GTenToneHL{alternative}{alternatives}:
	\begin{enumerate}[label=\arabic*.]
		\item \(\GTenToneHL{alternative}{A_1}\): Offer a standard produce box
		\item \(\GTenToneHL{alternative}{A_2}\): Offer a premium produce box
	\end{enumerate}
	
	The possible \GTenToneHL{state}{states} \GTenToneHL{state}{of} \GTenToneHL{state}{nature} are:
	\begin{enumerate}[label=\arabic*.]
		\item \(\GTenToneHL{state}{S_1}\): Strong seasonal demand
		\item \(\GTenToneHL{state}{S_2}\): Weak seasonal demand
	\end{enumerate}
	
	\columnbreak
	
	The quantities, measured in hundreds of subscriptions, are
	\[
	\begin{aligned}
		\GTenToneHL{quantity}{q}(\GTenToneHL{alternative}{A_1},\GTenToneHL{state}{S_1})&=\GTenToneHL{quantity}{7},
		&
		\GTenToneHL{quantity}{q}(\GTenToneHL{alternative}{A_1},\GTenToneHL{state}{S_2})&=\GTenToneHL{quantity}{4},
		\\
		\GTenToneHL{quantity}{q}(\GTenToneHL{alternative}{A_2},\GTenToneHL{state}{S_1})&=\GTenToneHL{quantity}{6},
		&
		\GTenToneHL{quantity}{q}(\GTenToneHL{alternative}{A_2},\GTenToneHL{state}{S_2})&=3.
	\end{aligned}
	\]
	
	The \GTenToneHL{alternative}{alternative}-based \GTenToneHL{revenue}{revenue} rates, in thousand USD per hundred
	subscriptions, are
	\[
	\GTenToneHL{revenue}{r}(\GTenToneHL{alternative}{A_1})=\GTenToneHL{revenue}{8},
	\qquad
	\GTenToneHL{revenue}{r}(\GTenToneHL{alternative}{A_2})=11.
	\]
	
	The state-based \GTenToneHL{revenue}{revenue} rates are
	\[
	\GTenToneHL{revenue}{r}(\GTenToneHL{state}{S_1})=\GTenToneHL{revenue}{3},
	\qquad
	\GTenToneHL{revenue}{r}(\GTenToneHL{state}{S_2})=1.
	\]
	
	The \GTenToneHL{alternative}{alternative}-based \GTenToneHL{cost}{cost} rates are
	\[
	\GTenToneHL{cost}{c}(\GTenToneHL{alternative}{A_1})=\GTenToneHL{cost}{4},
	\qquad
	\GTenToneHL{cost}{c}(\GTenToneHL{alternative}{A_2})=6.
	\]
	
	The state-based fixed \GTenToneHL{cost}{costs}, in thousand USD, are
	\[
	\GTenToneHL{cost}{f_c}(\GTenToneHL{state}{S_1})=\GTenToneHL{cost}{9},
	\qquad
	\GTenToneHL{cost}{f_c}(\GTenToneHL{state}{S_2})=13.
	\]
\end{multicols}

\textbf{Questions/Tasks.}
\begin{enumerate}
	\CriterionTask
	{subsec:problem9-farm-produce-subscription-c1}{1}
	{Interpreting decision alternatives, events, consequences, and states.}

	\CriterionTask
	{subsec:problem9-farm-produce-subscription-c3}{3}
	{Calculating profit and building the payoff table.}
\end{enumerate}

\subsubsection*{C1 - Interpreting decision alternatives, events, consequences, and states.}
\phantomsection
\label{subsec:problem9-farm-produce-subscription-c1}

\DecisionElementsTable
{Offer a standard produce box; Offer a premium produce box}
{Strong seasonal demand; Weak seasonal demand}
{Seasonal market conditions revealed after the subscription-plan decision}
{Profit in thousand USD for each plan--market-state pair}

\GTenToneHL{revenue}{Revenue} contains \GTenToneHL{alternative}{alternative}-based and state-based variable components:
\[
\GTenToneHL{revenue}{R}(\GTenToneHL{alternative}{A_i},\GTenToneHL{state}{S_j})
=
\GTenToneHL{quantity}{q}(\GTenToneHL{alternative}{A_i},\GTenToneHL{state}{S_j})
\left[
\GTenToneHL{revenue}{r}(\GTenToneHL{alternative}{A_i})+\GTenToneHL{revenue}{r}(\GTenToneHL{state}{S_j})
\right].
\]

\GTenToneHL{cost}{Cost} contains an \GTenToneHL{alternative}{alternative}-based variable component and a state-based
fixed component:
\[
\GTenToneHL{cost}{C}(\GTenToneHL{alternative}{A_i},\GTenToneHL{state}{S_j})
=
\GTenToneHL{quantity}{q}(\GTenToneHL{alternative}{A_i},\GTenToneHL{state}{S_j})\GTenToneHL{cost}{c}(\GTenToneHL{alternative}{A_i})+\GTenToneHL{cost}{f_c}(\GTenToneHL{state}{S_j}).
\]

Therefore,
\[
\GTenToneHL{payoff}{P}(\GTenToneHL{alternative}{A_i},\GTenToneHL{state}{S_j})
=
\GTenToneHL{quantity}{q}(\GTenToneHL{alternative}{A_i},\GTenToneHL{state}{S_j})
\left[
\GTenToneHL{revenue}{r}(\GTenToneHL{alternative}{A_i})+\GTenToneHL{revenue}{r}(\GTenToneHL{state}{S_j})
\right]
-
\left[
\GTenToneHL{quantity}{q}(\GTenToneHL{alternative}{A_i},\GTenToneHL{state}{S_j})\GTenToneHL{cost}{c}(\GTenToneHL{alternative}{A_i})+\GTenToneHL{cost}{f_c}(\GTenToneHL{state}{S_j})
\right].
\]

\secInicio[sec:problem9-farm-produce-subscription-c1-c3][sec:two-component-models]

\subsubsection*{C3 - Calculating profit and building the payoff table.}
\phantomsection
\label{subsec:problem9-farm-produce-subscription-c3}

For every \GTenToneHL{alternative}{alternative}--state pair, calculate
\[
\begin{aligned}
	\GTenToneHL{revenue}{R}(\GTenToneHL{alternative}{A_i},\GTenToneHL{state}{S_j})
	&=
	\GTenToneHL{quantity}{q}(\GTenToneHL{alternative}{A_i},\GTenToneHL{state}{S_j})
	\left[
	\GTenToneHL{revenue}{r}(\GTenToneHL{alternative}{A_i})+\GTenToneHL{revenue}{r}(\GTenToneHL{state}{S_j})
	\right],
	\\
	\GTenToneHL{cost}{C}(\GTenToneHL{alternative}{A_i},\GTenToneHL{state}{S_j})
	&=
	\GTenToneHL{quantity}{q}(\GTenToneHL{alternative}{A_i},\GTenToneHL{state}{S_j})\GTenToneHL{cost}{c}(\GTenToneHL{alternative}{A_i})+\GTenToneHL{cost}{f_c}(\GTenToneHL{state}{S_j}),
	\\
	\GTenToneHL{payoff}{P}(\GTenToneHL{alternative}{A_i},\GTenToneHL{state}{S_j})
	&=
	\GTenToneHL{revenue}{R}(\GTenToneHL{alternative}{A_i},\GTenToneHL{state}{S_j})-\GTenToneHL{cost}{C}(\GTenToneHL{alternative}{A_i},\GTenToneHL{state}{S_j}).
\end{aligned}
\]

\begin{center}
	\small
	\begin{tabular}{llll}
		\toprule
		Cell
		& \(\GTenToneHL{revenue}{R}(\GTenToneHL{alternative}{A_i},\GTenToneHL{state}{S_j})\)
		& \(\GTenToneHL{cost}{C}(\GTenToneHL{alternative}{A_i},\GTenToneHL{state}{S_j})\)
		& \(\GTenToneHL{payoff}{P}(\GTenToneHL{alternative}{A_i},\GTenToneHL{state}{S_j})=\GTenToneHL{revenue}{R}(\GTenToneHL{alternative}{A_i},\GTenToneHL{state}{S_j})-\GTenToneHL{cost}{C}(\GTenToneHL{alternative}{A_i},\GTenToneHL{state}{S_j})\) \\
		\midrule
		\((\GTenToneHL{alternative}{A_1},\GTenToneHL{state}{S_1})\)
		& \(7(8+3)=77\)
		& \(7(4)+9=37\)
		& \(77-37=40\) \\
		
		\((\GTenToneHL{alternative}{A_1},\GTenToneHL{state}{S_2})\)
		& \(4(8+1)=36\)
		& \(4(4)+13=29\)
		& \(36-29=7\) \\
		
		\((\GTenToneHL{alternative}{A_2},\GTenToneHL{state}{S_1})\)
		& \(6(11+3)=84\)
		& \(6(6)+9=45\)
		& \(84-45=39\) \\
		
		\((\GTenToneHL{alternative}{A_2},\GTenToneHL{state}{S_2})\)
		& \(3(11+1)=36\)
		& \(3(6)+13=31\)
		& \(36-31=5\) \\
		\bottomrule
	\end{tabular}
\end{center}

The resulting \GTenToneHL{payoff}{payoff} table is:

\begin{center}
	\begin{tabular}{lrr}
		\toprule
		Alternative & Strong seasonal demand & Weak seasonal demand \\
		\midrule
		Standard produce box & \GTenToneHL{payoff}{40} & \GTenToneHL{payoff}{7} \\
		Premium produce box & \GTenToneHL{payoff}{39} & \GTenToneHL{payoff}{5} \\
		\bottomrule
	\end{tabular}
\end{center}

\secInicio[sec:problem9-farm-produce-subscription-c1-c3][sec:two-component-models]

% ----------------------------------------------------------------
\subsection{Mobile Catering Configuration (C1--C3)}
\label{sec:problem10-mobile-catering-configuration-c1-c3}

A mobile-catering company must choose a service configuration before event
demand and ingredient-market conditions are known. \GTenToneHL{revenue}{Revenue} consists of a
configuration-based rate per service order and a fixed booking payment
determined by the selected configuration. \GTenToneHL{cost}{Cost} consists of a state-based
rate per order and a fixed setup \GTenToneHL{cost}{cost} determined by the configuration.

\begin{multicols}{2}
	The company is comparing two \GTenToneHL{alternative}{alternatives}:
	\begin{enumerate}[label=\arabic*.]
		\item \(\GTenToneHL{alternative}{A_1}\): Prepare a compact catering station
		\item \(\GTenToneHL{alternative}{A_2}\): Prepare a full mobile kitchen
	\end{enumerate}
	
	The possible \GTenToneHL{state}{states} \GTenToneHL{state}{of} \GTenToneHL{state}{nature} are:
	\begin{enumerate}[label=\arabic*.]
		\item \(\GTenToneHL{state}{S_1}\): Busy event with stable ingredient prices
		\item \(\GTenToneHL{state}{S_2}\): Quiet event with higher ingredient prices
	\end{enumerate}
	
	\columnbreak
	
	The quantities, measured in hundreds of service orders, are
	\[
	\begin{aligned}
		\GTenToneHL{quantity}{q}(\GTenToneHL{alternative}{A_1},\GTenToneHL{state}{S_1})&=\GTenToneHL{quantity}{8},
		&
		\GTenToneHL{quantity}{q}(\GTenToneHL{alternative}{A_1},\GTenToneHL{state}{S_2})&=\GTenToneHL{quantity}{4},
		\\
		\GTenToneHL{quantity}{q}(\GTenToneHL{alternative}{A_2},\GTenToneHL{state}{S_1})&=\GTenToneHL{quantity}{10},
		&
		\GTenToneHL{quantity}{q}(\GTenToneHL{alternative}{A_2},\GTenToneHL{state}{S_2})&=5.
	\end{aligned}
	\]
	
	The \GTenToneHL{alternative}{alternative}-based \GTenToneHL{revenue}{revenue} rates, in thousand USD per hundred
	orders, are
	\[
	\GTenToneHL{revenue}{r}(\GTenToneHL{alternative}{A_1})=\GTenToneHL{revenue}{7},
	\qquad
	\GTenToneHL{revenue}{r}(\GTenToneHL{alternative}{A_2})=9.
	\]
	
	The \GTenToneHL{alternative}{alternative}-based fixed revenues, in thousand USD, are
	\[
	\GTenToneHL{revenue}{f_r}(\GTenToneHL{alternative}{A_1})=\GTenToneHL{revenue}{18},
	\qquad
	\GTenToneHL{revenue}{f_r}(\GTenToneHL{alternative}{A_2})=25.
	\]
	
	The state-based \GTenToneHL{cost}{cost} rates, in thousand USD per hundred orders, are
	\[
	\GTenToneHL{cost}{c}(\GTenToneHL{state}{S_1})=\GTenToneHL{cost}{3},
	\qquad
	\GTenToneHL{cost}{c}(\GTenToneHL{state}{S_2})=5.
	\]
	
	The \GTenToneHL{alternative}{alternative}-based fixed \GTenToneHL{cost}{costs}, in thousand USD, are
	\[
	\GTenToneHL{cost}{f_c}(\GTenToneHL{alternative}{A_1})=\GTenToneHL{cost}{20},
	\qquad
	\GTenToneHL{cost}{f_c}(\GTenToneHL{alternative}{A_2})=32.
	\]
\end{multicols}

\textbf{Questions/Tasks.}
\begin{enumerate}
	\CriterionTask
	{subsec:problem10-mobile-catering-configuration-c1}{1}
	{Interpreting decision alternatives, events, consequences, and states.}

	\CriterionTask
	{subsec:problem10-mobile-catering-configuration-c3}{3}
	{Calculating profit and building the payoff table.}
\end{enumerate}

\subsubsection*{C1 - Interpreting decision alternatives, events, consequences, and states.}
\phantomsection
\label{subsec:problem10-mobile-catering-configuration-c1}

\DecisionElementsTable
{Prepare a compact catering station; Prepare a full mobile kitchen}
{Busy event with stable ingredient prices; Quiet event with higher ingredient prices}
{Event demand and ingredient-market conditions revealed after the configuration decision}
{Profit in thousand USD for each configuration--market-state pair}

\GTenToneHL{revenue}{Revenue} contains \GTenToneHL{alternative}{alternative}-based variable and fixed components:
\[
\GTenToneHL{revenue}{R}(\GTenToneHL{alternative}{A_i},\GTenToneHL{state}{S_j})
=
\GTenToneHL{quantity}{q}(\GTenToneHL{alternative}{A_i},\GTenToneHL{state}{S_j})\GTenToneHL{revenue}{r}(\GTenToneHL{alternative}{A_i})+\GTenToneHL{revenue}{f_r}(\GTenToneHL{alternative}{A_i}).
\]

\GTenToneHL{cost}{Cost} contains a state-based variable component and an \GTenToneHL{alternative}{alternative}-based
fixed component:
\[
\GTenToneHL{cost}{C}(\GTenToneHL{alternative}{A_i},\GTenToneHL{state}{S_j})
=
\GTenToneHL{quantity}{q}(\GTenToneHL{alternative}{A_i},\GTenToneHL{state}{S_j})\GTenToneHL{cost}{c}(\GTenToneHL{state}{S_j})+\GTenToneHL{cost}{f_c}(\GTenToneHL{alternative}{A_i}).
\]

Therefore,
\[
\GTenToneHL{payoff}{P}(\GTenToneHL{alternative}{A_i},\GTenToneHL{state}{S_j})
=
\left[
\GTenToneHL{quantity}{q}(\GTenToneHL{alternative}{A_i},\GTenToneHL{state}{S_j})\GTenToneHL{revenue}{r}(\GTenToneHL{alternative}{A_i})+\GTenToneHL{revenue}{f_r}(\GTenToneHL{alternative}{A_i})
\right]
-
\left[
\GTenToneHL{quantity}{q}(\GTenToneHL{alternative}{A_i},\GTenToneHL{state}{S_j})\GTenToneHL{cost}{c}(\GTenToneHL{state}{S_j})+\GTenToneHL{cost}{f_c}(\GTenToneHL{alternative}{A_i})
\right].
\]

\secInicio[sec:problem10-mobile-catering-configuration-c1-c3][sec:two-component-models]

\subsubsection*{C3 - Calculating profit and building the payoff table.}
\phantomsection
\label{subsec:problem10-mobile-catering-configuration-c3}

For every \GTenToneHL{alternative}{alternative}--state pair, calculate
\[
\begin{aligned}
	\GTenToneHL{revenue}{R}(\GTenToneHL{alternative}{A_i},\GTenToneHL{state}{S_j})
	&=
	\GTenToneHL{quantity}{q}(\GTenToneHL{alternative}{A_i},\GTenToneHL{state}{S_j})\GTenToneHL{revenue}{r}(\GTenToneHL{alternative}{A_i})+\GTenToneHL{revenue}{f_r}(\GTenToneHL{alternative}{A_i}),
	\\
	\GTenToneHL{cost}{C}(\GTenToneHL{alternative}{A_i},\GTenToneHL{state}{S_j})
	&=
	\GTenToneHL{quantity}{q}(\GTenToneHL{alternative}{A_i},\GTenToneHL{state}{S_j})\GTenToneHL{cost}{c}(\GTenToneHL{state}{S_j})+\GTenToneHL{cost}{f_c}(\GTenToneHL{alternative}{A_i}),
	\\
	\GTenToneHL{payoff}{P}(\GTenToneHL{alternative}{A_i},\GTenToneHL{state}{S_j})
	&=
	\GTenToneHL{revenue}{R}(\GTenToneHL{alternative}{A_i},\GTenToneHL{state}{S_j})-\GTenToneHL{cost}{C}(\GTenToneHL{alternative}{A_i},\GTenToneHL{state}{S_j}).
\end{aligned}
\]

\begin{center}
	\small
	\begin{tabular}{llll}
		\toprule
		Cell
		& \(\GTenToneHL{revenue}{R}(\GTenToneHL{alternative}{A_i},\GTenToneHL{state}{S_j})\)
		& \(\GTenToneHL{cost}{C}(\GTenToneHL{alternative}{A_i},\GTenToneHL{state}{S_j})\)
		& \(\GTenToneHL{payoff}{P}(\GTenToneHL{alternative}{A_i},\GTenToneHL{state}{S_j})=\GTenToneHL{revenue}{R}(\GTenToneHL{alternative}{A_i},\GTenToneHL{state}{S_j})-\GTenToneHL{cost}{C}(\GTenToneHL{alternative}{A_i},\GTenToneHL{state}{S_j})\) \\
		\midrule
		\((\GTenToneHL{alternative}{A_1},\GTenToneHL{state}{S_1})\)
		& \(8(7)+18=74\)
		& \(8(3)+20=44\)
		& \(74-44=30\) \\
		
		\((\GTenToneHL{alternative}{A_1},\GTenToneHL{state}{S_2})\)
		& \(4(7)+18=46\)
		& \(4(5)+20=40\)
		& \(46-40=6\) \\
		
		\((\GTenToneHL{alternative}{A_2},\GTenToneHL{state}{S_1})\)
		& \(10(9)+25=115\)
		& \(10(3)+32=62\)
		& \(115-62=53\) \\
		
		\((\GTenToneHL{alternative}{A_2},\GTenToneHL{state}{S_2})\)
		& \(5(9)+25=70\)
		& \(5(5)+32=57\)
		& \(70-57=13\) \\
		\bottomrule
	\end{tabular}
\end{center}

The resulting \GTenToneHL{payoff}{payoff} table is:

\begin{center}
	\begin{tabular}{lrr}
		\toprule
		Alternative & Busy event & Quiet event \\
		\midrule
		Compact catering station & \GTenToneHL{payoff}{30} & \GTenToneHL{payoff}{6} \\
		Full mobile kitchen & \GTenToneHL{payoff}{53} & \GTenToneHL{payoff}{13} \\
		\bottomrule
	\end{tabular}
\end{center}

\secInicio[sec:problem10-mobile-catering-configuration-c1-c3][sec:two-component-models]

% ----------------------------------------------------------------
\subsection{Online Training Programme (C1--C3)}
\label{sec:problem11-online-training-programme-c1-c3}

An online-training provider must select a programme before enrolment conditions
are known. \GTenToneHL{revenue}{Revenue} consists of a programme-based rate per enrolment group and
a fixed payment determined by the enrolment state. \GTenToneHL{cost}{Cost} consists of a
state-based delivery rate and a fixed platform \GTenToneHL{cost}{cost} determined by the state.

\begin{multicols}{2}
	The provider is comparing two \GTenToneHL{alternative}{alternatives}:
	\begin{enumerate}[label=\arabic*.]
		\item \(\GTenToneHL{alternative}{A_1}\): Offer the foundation programme
		\item \(\GTenToneHL{alternative}{A_2}\): Offer the advanced programme
	\end{enumerate}
	
	The possible \GTenToneHL{state}{states} \GTenToneHL{state}{of} \GTenToneHL{state}{nature} are:
	\begin{enumerate}[label=\arabic*.]
		\item \(\GTenToneHL{state}{S_1}\): Corporate enrolment
		\item \(\GTenToneHL{state}{S_2}\): Individual enrolment
	\end{enumerate}
	
	\columnbreak
	
	The quantities, measured in enrolment groups, are
	\[
	\begin{aligned}
		\GTenToneHL{quantity}{q}(\GTenToneHL{alternative}{A_1},\GTenToneHL{state}{S_1})&=\GTenToneHL{quantity}{5},
		&
		\GTenToneHL{quantity}{q}(\GTenToneHL{alternative}{A_1},\GTenToneHL{state}{S_2})&=\GTenToneHL{quantity}{3},
		\\
		\GTenToneHL{quantity}{q}(\GTenToneHL{alternative}{A_2},\GTenToneHL{state}{S_1})&=\GTenToneHL{quantity}{4},
		&
		\GTenToneHL{quantity}{q}(\GTenToneHL{alternative}{A_2},\GTenToneHL{state}{S_2})&=2.
	\end{aligned}
	\]
	
	The \GTenToneHL{alternative}{alternative}-based \GTenToneHL{revenue}{revenue} rates, in thousand USD per group, are
	\[
	\GTenToneHL{revenue}{r}(\GTenToneHL{alternative}{A_1})=\GTenToneHL{revenue}{14},
	\qquad
	\GTenToneHL{revenue}{r}(\GTenToneHL{alternative}{A_2})=20.
	\]
	
	The state-based fixed revenues, in thousand USD, are
	\[
	\GTenToneHL{revenue}{f_r}(\GTenToneHL{state}{S_1})=\GTenToneHL{revenue}{12},
	\qquad
	\GTenToneHL{revenue}{f_r}(\GTenToneHL{state}{S_2})=7.
	\]
	
	The state-based \GTenToneHL{cost}{cost} rates, in thousand USD per group, are
	\[
	\GTenToneHL{cost}{c}(\GTenToneHL{state}{S_1})=\GTenToneHL{cost}{6},
	\qquad
	\GTenToneHL{cost}{c}(\GTenToneHL{state}{S_2})=8.
	\]
	
	The state-based fixed \GTenToneHL{cost}{costs}, in thousand USD, are
	\[
	\GTenToneHL{cost}{f_c}(\GTenToneHL{state}{S_1})=\GTenToneHL{cost}{10},
	\qquad
	\GTenToneHL{cost}{f_c}(\GTenToneHL{state}{S_2})=14.
	\]
\end{multicols}

\textbf{Questions/Tasks.}
\begin{enumerate}
	\CriterionTask
	{subsec:problem11-online-training-programme-c1}{1}
	{Interpreting decision alternatives, events, consequences, and states.}

	\CriterionTask
	{subsec:problem11-online-training-programme-c3}{3}
	{Calculating profit and building the payoff table.}
\end{enumerate}

\subsubsection*{C1 - Interpreting decision alternatives, events, consequences, and states.}
\phantomsection
\label{subsec:problem11-online-training-programme-c1}

\DecisionElementsTable
{Offer the foundation programme; Offer the advanced programme}
{Corporate enrolment; Individual enrolment}
{Enrolment type revealed after the programme decision}
{Profit in thousand USD for each programme--enrolment pair}

\GTenToneHL{revenue}{Revenue} contains an \GTenToneHL{alternative}{alternative}-based variable component and a state-based
fixed component:
\[
\GTenToneHL{revenue}{R}(\GTenToneHL{alternative}{A_i},\GTenToneHL{state}{S_j})
=
\GTenToneHL{quantity}{q}(\GTenToneHL{alternative}{A_i},\GTenToneHL{state}{S_j})\GTenToneHL{revenue}{r}(\GTenToneHL{alternative}{A_i})+\GTenToneHL{revenue}{f_r}(\GTenToneHL{state}{S_j}).
\]

\GTenToneHL{cost}{Cost} contains state-based variable and fixed components:
\[
\GTenToneHL{cost}{C}(\GTenToneHL{alternative}{A_i},\GTenToneHL{state}{S_j})
=
\GTenToneHL{quantity}{q}(\GTenToneHL{alternative}{A_i},\GTenToneHL{state}{S_j})\GTenToneHL{cost}{c}(\GTenToneHL{state}{S_j})+\GTenToneHL{cost}{f_c}(\GTenToneHL{state}{S_j}).
\]

Therefore,
\[
\GTenToneHL{payoff}{P}(\GTenToneHL{alternative}{A_i},\GTenToneHL{state}{S_j})
=
\left[
\GTenToneHL{quantity}{q}(\GTenToneHL{alternative}{A_i},\GTenToneHL{state}{S_j})\GTenToneHL{revenue}{r}(\GTenToneHL{alternative}{A_i})+\GTenToneHL{revenue}{f_r}(\GTenToneHL{state}{S_j})
\right]
-
\left[
\GTenToneHL{quantity}{q}(\GTenToneHL{alternative}{A_i},\GTenToneHL{state}{S_j})\GTenToneHL{cost}{c}(\GTenToneHL{state}{S_j})+\GTenToneHL{cost}{f_c}(\GTenToneHL{state}{S_j})
\right].
\]

\secInicio[sec:problem11-online-training-programme-c1-c3][sec:two-component-models]

\subsubsection*{C3 - Calculating profit and building the payoff table.}
\phantomsection
\label{subsec:problem11-online-training-programme-c3}

For every \GTenToneHL{alternative}{alternative}--state pair, calculate
\[
\begin{aligned}
	\GTenToneHL{revenue}{R}(\GTenToneHL{alternative}{A_i},\GTenToneHL{state}{S_j})
	&=
	\GTenToneHL{quantity}{q}(\GTenToneHL{alternative}{A_i},\GTenToneHL{state}{S_j})\GTenToneHL{revenue}{r}(\GTenToneHL{alternative}{A_i})+\GTenToneHL{revenue}{f_r}(\GTenToneHL{state}{S_j}),
	\\
	\GTenToneHL{cost}{C}(\GTenToneHL{alternative}{A_i},\GTenToneHL{state}{S_j})
	&=
	\GTenToneHL{quantity}{q}(\GTenToneHL{alternative}{A_i},\GTenToneHL{state}{S_j})\GTenToneHL{cost}{c}(\GTenToneHL{state}{S_j})+\GTenToneHL{cost}{f_c}(\GTenToneHL{state}{S_j}),
	\\
	\GTenToneHL{payoff}{P}(\GTenToneHL{alternative}{A_i},\GTenToneHL{state}{S_j})
	&=
	\GTenToneHL{revenue}{R}(\GTenToneHL{alternative}{A_i},\GTenToneHL{state}{S_j})-\GTenToneHL{cost}{C}(\GTenToneHL{alternative}{A_i},\GTenToneHL{state}{S_j}).
\end{aligned}
\]

\begin{center}
	\small
	\begin{tabular}{llll}
		\toprule
		Cell
		& \(\GTenToneHL{revenue}{R}(\GTenToneHL{alternative}{A_i},\GTenToneHL{state}{S_j})\)
		& \(\GTenToneHL{cost}{C}(\GTenToneHL{alternative}{A_i},\GTenToneHL{state}{S_j})\)
		& \(\GTenToneHL{payoff}{P}(\GTenToneHL{alternative}{A_i},\GTenToneHL{state}{S_j})=\GTenToneHL{revenue}{R}(\GTenToneHL{alternative}{A_i},\GTenToneHL{state}{S_j})-\GTenToneHL{cost}{C}(\GTenToneHL{alternative}{A_i},\GTenToneHL{state}{S_j})\) \\
		\midrule
		\((\GTenToneHL{alternative}{A_1},\GTenToneHL{state}{S_1})\)
		& \(5(14)+12=82\)
		& \(5(6)+10=40\)
		& \(82-40=42\) \\
		
		\((\GTenToneHL{alternative}{A_1},\GTenToneHL{state}{S_2})\)
		& \(3(14)+7=49\)
		& \(3(8)+14=38\)
		& \(49-38=11\) \\
		
		\((\GTenToneHL{alternative}{A_2},\GTenToneHL{state}{S_1})\)
		& \(4(20)+12=92\)
		& \(4(6)+10=34\)
		& \(92-34=58\) \\
		
		\((\GTenToneHL{alternative}{A_2},\GTenToneHL{state}{S_2})\)
		& \(2(20)+7=47\)
		& \(2(8)+14=30\)
		& \(47-30=17\) \\
		\bottomrule
	\end{tabular}
\end{center}

The resulting \GTenToneHL{payoff}{payoff} table is:

\begin{center}
	\begin{tabular}{lrr}
		\toprule
		Alternative & Corporate enrolment & Individual enrolment \\
		\midrule
		Foundation programme & \GTenToneHL{rate}{42} & \GTenToneHL{payoff}{11} \\
		Advanced programme & \GTenToneHL{rate}{58} & \GTenToneHL{payoff}{17} \\
		\bottomrule
	\end{tabular}
\end{center}

\secInicio[sec:problem11-online-training-programme-c1-c3][sec:two-component-models]

% ----------------------------------------------------------------
\subsection{Community Venue Membership (C1--C3)}
\label{sec:problem12-community-venue-membership-c1-c3}

A community organisation must select a venue-membership plan before booking
conditions and permit requirements are known. \GTenToneHL{revenue}{Revenue} consists of a
state-based rate per booking and a fixed sponsorship payment determined by
the selected membership plan. Costs consist of two fixed components: one
determined by the plan and one determined by the state.

\begin{multicols}{2}
	The organisation is comparing two \GTenToneHL{alternative}{alternatives}:
	\begin{enumerate}[label=\arabic*.]
		\item \(\GTenToneHL{alternative}{A_1}\): Select the standard venue membership
		\item \(\GTenToneHL{alternative}{A_2}\): Select the expanded venue membership
	\end{enumerate}
	
	The possible \GTenToneHL{state}{states} \GTenToneHL{state}{of} \GTenToneHL{state}{nature} are:
	\begin{enumerate}[label=\arabic*.]
		\item \(\GTenToneHL{state}{S_1}\): High booking demand
		\item \(\GTenToneHL{state}{S_2}\): Low booking demand
	\end{enumerate}
	
	\columnbreak
	
	The numbers of bookings are
	\[
	\begin{aligned}
		\GTenToneHL{quantity}{q}(\GTenToneHL{alternative}{A_1},\GTenToneHL{state}{S_1})&=\GTenToneHL{quantity}{6},
		&
		\GTenToneHL{quantity}{q}(\GTenToneHL{alternative}{A_1},\GTenToneHL{state}{S_2})&=\GTenToneHL{quantity}{3},
		\\
		\GTenToneHL{quantity}{q}(\GTenToneHL{alternative}{A_2},\GTenToneHL{state}{S_1})&=\GTenToneHL{quantity}{8},
		&
		\GTenToneHL{quantity}{q}(\GTenToneHL{alternative}{A_2},\GTenToneHL{state}{S_2})&=4.
	\end{aligned}
	\]
	
	The state-based \GTenToneHL{revenue}{revenue} rates, in thousand USD per booking, are
	\[
	\GTenToneHL{revenue}{r}(\GTenToneHL{state}{S_1})=\GTenToneHL{revenue}{10},
	\qquad
	\GTenToneHL{revenue}{r}(\GTenToneHL{state}{S_2})=7.
	\]
	
	The \GTenToneHL{alternative}{alternative}-based fixed revenues, in thousand USD, are
	\[
	\GTenToneHL{revenue}{f_r}(\GTenToneHL{alternative}{A_1})=\GTenToneHL{revenue}{15},
	\qquad
	\GTenToneHL{revenue}{f_r}(\GTenToneHL{alternative}{A_2})=22.
	\]
	
	The \GTenToneHL{alternative}{alternative}-based fixed \GTenToneHL{cost}{costs}, in thousand USD, are
	\[
	\GTenToneHL{cost}{f_c}(\GTenToneHL{alternative}{A_1})=\GTenToneHL{cost}{28},
	\qquad
	\GTenToneHL{cost}{f_c}(\GTenToneHL{alternative}{A_2})=40.
	\]
	
	The state-based fixed \GTenToneHL{cost}{costs}, in thousand USD, are
	\[
	\GTenToneHL{cost}{f_c}(\GTenToneHL{state}{S_1})=\GTenToneHL{cost}{8},
	\qquad
	\GTenToneHL{cost}{f_c}(\GTenToneHL{state}{S_2})=13.
	\]
\end{multicols}

\textbf{Questions/Tasks.}
\begin{enumerate}
	\CriterionTask
	{subsec:problem12-community-venue-membership-c1}{1}
	{Interpreting decision alternatives, events, consequences, and states.}

	\CriterionTask
	{subsec:problem12-community-venue-membership-c3}{3}
	{Calculating profit and building the payoff table.}
\end{enumerate}

\subsubsection*{C1 - Interpreting decision alternatives, events, consequences, and states.}
\phantomsection
\label{subsec:problem12-community-venue-membership-c1}

\DecisionElementsTable
{Select the standard venue membership; Select the expanded venue membership}
{High booking demand; Low booking demand}
{Booking demand revealed after the membership-plan decision}
{Profit in thousand USD for each membership--demand pair}

\GTenToneHL{revenue}{Revenue} contains a state-based variable component and an \GTenToneHL{alternative}{alternative}-based
fixed component:
\[
\GTenToneHL{revenue}{R}(\GTenToneHL{alternative}{A_i},\GTenToneHL{state}{S_j})
=
\GTenToneHL{quantity}{q}(\GTenToneHL{alternative}{A_i},\GTenToneHL{state}{S_j})\GTenToneHL{revenue}{r}(\GTenToneHL{state}{S_j})+\GTenToneHL{revenue}{f_r}(\GTenToneHL{alternative}{A_i}).
\]

\GTenToneHL{cost}{Cost} contains \GTenToneHL{alternative}{alternative}-based and state-based fixed components:
\[
\GTenToneHL{cost}{C}(\GTenToneHL{alternative}{A_i},\GTenToneHL{state}{S_j})
=
\GTenToneHL{cost}{f_c}(\GTenToneHL{alternative}{A_i})+\GTenToneHL{cost}{f_c}(\GTenToneHL{state}{S_j}).
\]

Therefore,
\[
\GTenToneHL{payoff}{P}(\GTenToneHL{alternative}{A_i},\GTenToneHL{state}{S_j})
=
\left[
\GTenToneHL{quantity}{q}(\GTenToneHL{alternative}{A_i},\GTenToneHL{state}{S_j})\GTenToneHL{revenue}{r}(\GTenToneHL{state}{S_j})+\GTenToneHL{revenue}{f_r}(\GTenToneHL{alternative}{A_i})
\right]
-
\left[
\GTenToneHL{cost}{f_c}(\GTenToneHL{alternative}{A_i})+\GTenToneHL{cost}{f_c}(\GTenToneHL{state}{S_j})
\right].
\]

\secInicio[sec:problem12-community-venue-membership-c1-c3][sec:two-component-models]

\subsubsection*{C3 - Calculating profit and building the payoff table.}
\phantomsection
\label{subsec:problem12-community-venue-membership-c3}

For every \GTenToneHL{alternative}{alternative}--state pair, calculate
\[
\begin{aligned}
	\GTenToneHL{revenue}{R}(\GTenToneHL{alternative}{A_i},\GTenToneHL{state}{S_j})
	&=
	\GTenToneHL{quantity}{q}(\GTenToneHL{alternative}{A_i},\GTenToneHL{state}{S_j})\GTenToneHL{revenue}{r}(\GTenToneHL{state}{S_j})+\GTenToneHL{revenue}{f_r}(\GTenToneHL{alternative}{A_i}),
	\\
	\GTenToneHL{cost}{C}(\GTenToneHL{alternative}{A_i},\GTenToneHL{state}{S_j})
	&=
	\GTenToneHL{cost}{f_c}(\GTenToneHL{alternative}{A_i})+\GTenToneHL{cost}{f_c}(\GTenToneHL{state}{S_j}),
	\\
	\GTenToneHL{payoff}{P}(\GTenToneHL{alternative}{A_i},\GTenToneHL{state}{S_j})
	&=
	\GTenToneHL{revenue}{R}(\GTenToneHL{alternative}{A_i},\GTenToneHL{state}{S_j})-\GTenToneHL{cost}{C}(\GTenToneHL{alternative}{A_i},\GTenToneHL{state}{S_j}).
\end{aligned}
\]

\begin{center}
	\small
	\begin{tabular}{llll}
		\toprule
		Cell
		& \(\GTenToneHL{revenue}{R}(\GTenToneHL{alternative}{A_i},\GTenToneHL{state}{S_j})\)
		& \(\GTenToneHL{cost}{C}(\GTenToneHL{alternative}{A_i},\GTenToneHL{state}{S_j})\)
		& \(\GTenToneHL{payoff}{P}(\GTenToneHL{alternative}{A_i},\GTenToneHL{state}{S_j})=\GTenToneHL{revenue}{R}(\GTenToneHL{alternative}{A_i},\GTenToneHL{state}{S_j})-\GTenToneHL{cost}{C}(\GTenToneHL{alternative}{A_i},\GTenToneHL{state}{S_j})\) \\
		\midrule
		\((\GTenToneHL{alternative}{A_1},\GTenToneHL{state}{S_1})\)
		& \(6(10)+15=75\)
		& \(28+8=36\)
		& \(75-36=39\) \\
		
		\((\GTenToneHL{alternative}{A_1},\GTenToneHL{state}{S_2})\)
		& \(3(7)+15=36\)
		& \(28+13=41\)
		& \(36-41=-5\) \\
		
		\((\GTenToneHL{alternative}{A_2},\GTenToneHL{state}{S_1})\)
		& \(8(10)+22=102\)
		& \(40+8=48\)
		& \(102-48=54\) \\
		
		\((\GTenToneHL{alternative}{A_2},\GTenToneHL{state}{S_2})\)
		& \(4(7)+22=50\)
		& \(40+13=53\)
		& \(50-53=-3\) \\
		\bottomrule
	\end{tabular}
\end{center}

The resulting \GTenToneHL{payoff}{payoff} table is:

\begin{center}
	\begin{tabular}{lrr}
		\toprule
		Alternative & High booking demand & Low booking demand \\
		\midrule
		Standard venue membership & \GTenToneHL{payoff}{39} & \(-5\) \\
		Expanded venue membership & \GTenToneHL{payoff}{54} & \(-3\) \\
		\bottomrule
	\end{tabular}
\end{center}

The negative \GTenToneHL{payoff}{payoffs} under low booking demand represent losses because the
combined fixed \GTenToneHL{cost}{costs} exceed the \GTenToneHL{revenue}{revenue} earned.

\secInicio[sec:problem12-community-venue-membership-c1-c3][sec:two-component-models]

% ----------------------------------------------------------------
\subsection{Urban Delivery Fleet (C1--C3)}
\label{sec:problem13-urban-delivery-fleet-c1-c3}

An urban-delivery company must choose a fleet configuration before delivery
demand and traffic conditions are known. \GTenToneHL{revenue}{Revenue} consists of a state-based
rate per delivery batch and a fixed performance payment determined by the
state. Variable \GTenToneHL{cost}{cost} per batch combines a fleet-based operating rate and a
state-based traffic surcharge.

\begin{multicols}{2}
	The company is comparing two \GTenToneHL{alternative}{alternatives}:
	\begin{enumerate}[label=\arabic*.]
		\item \(\GTenToneHL{alternative}{A_1}\): Use a conventional van fleet
		\item \(\GTenToneHL{alternative}{A_2}\): Use an electric cargo-bike fleet
	\end{enumerate}
	
	The possible \GTenToneHL{state}{states} \GTenToneHL{state}{of} \GTenToneHL{state}{nature} are:
	\begin{enumerate}[label=\arabic*.]
		\item \(\GTenToneHL{state}{S_1}\): High demand with light traffic
		\item \(\GTenToneHL{state}{S_2}\): Low demand with heavy traffic
	\end{enumerate}
	
	\columnbreak
	
	The quantities, measured in delivery batches, are
	\[
	\begin{aligned}
		\GTenToneHL{quantity}{q}(\GTenToneHL{alternative}{A_1},\GTenToneHL{state}{S_1})&=\GTenToneHL{quantity}{9},
		&
		\GTenToneHL{quantity}{q}(\GTenToneHL{alternative}{A_1},\GTenToneHL{state}{S_2})&=\GTenToneHL{quantity}{5},
		\\
		\GTenToneHL{quantity}{q}(\GTenToneHL{alternative}{A_2},\GTenToneHL{state}{S_1})&=\GTenToneHL{quantity}{12},
		&
		\GTenToneHL{quantity}{q}(\GTenToneHL{alternative}{A_2},\GTenToneHL{state}{S_2})&=6.
	\end{aligned}
	\]
	
	The state-based \GTenToneHL{revenue}{revenue} rates, in thousand USD per batch, are
	\[
	\GTenToneHL{revenue}{r}(\GTenToneHL{state}{S_1})=\GTenToneHL{revenue}{8},
	\qquad
	\GTenToneHL{revenue}{r}(\GTenToneHL{state}{S_2})=6.
	\]
	
	The state-based fixed revenues, in thousand USD, are
	\[
	\GTenToneHL{revenue}{f_r}(\GTenToneHL{state}{S_1})=\GTenToneHL{revenue}{14},
	\qquad
	\GTenToneHL{revenue}{f_r}(\GTenToneHL{state}{S_2})=8.
	\]
	
	The \GTenToneHL{alternative}{alternative}-based \GTenToneHL{cost}{cost} rates, in thousand USD per batch, are
	\[
	\GTenToneHL{cost}{c}(\GTenToneHL{alternative}{A_1})=\GTenToneHL{cost}{3},
	\qquad
	\GTenToneHL{cost}{c}(\GTenToneHL{alternative}{A_2})=2.
	\]
	
	The state-based \GTenToneHL{cost}{cost} rates, in thousand USD per batch, are
	\[
	\GTenToneHL{cost}{c}(\GTenToneHL{state}{S_1})=\GTenToneHL{cost}{1},
	\qquad
	\GTenToneHL{cost}{c}(\GTenToneHL{state}{S_2})=3.
	\]
\end{multicols}

\textbf{Questions/Tasks.}
\begin{enumerate}
	\CriterionTask
	{subsec:problem13-urban-delivery-fleet-c1}{1}
	{Interpreting decision alternatives, events, consequences, and states.}

	\CriterionTask
	{subsec:problem13-urban-delivery-fleet-c3}{3}
	{Calculating profit and building the payoff table.}
\end{enumerate}

\subsubsection*{C1 - Interpreting decision alternatives, events, consequences, and states.}
\phantomsection
\label{subsec:problem13-urban-delivery-fleet-c1}

\DecisionElementsTable
{Use a conventional van fleet; Use an electric cargo-bike fleet}
{High demand with light traffic; Low demand with heavy traffic}
{Delivery demand and traffic conditions revealed after the fleet decision}
{Profit in thousand USD for each fleet--market-state pair}

\GTenToneHL{revenue}{Revenue} contains state-based variable and fixed components:
\[
\GTenToneHL{revenue}{R}(\GTenToneHL{alternative}{A_i},\GTenToneHL{state}{S_j})
=
\GTenToneHL{quantity}{q}(\GTenToneHL{alternative}{A_i},\GTenToneHL{state}{S_j})\GTenToneHL{revenue}{r}(\GTenToneHL{state}{S_j})+\GTenToneHL{revenue}{f_r}(\GTenToneHL{state}{S_j}).
\]

\GTenToneHL{cost}{Cost} contains \GTenToneHL{alternative}{alternative}-based and state-based variable components:
\[
\GTenToneHL{cost}{C}(\GTenToneHL{alternative}{A_i},\GTenToneHL{state}{S_j})
=
\GTenToneHL{quantity}{q}(\GTenToneHL{alternative}{A_i},\GTenToneHL{state}{S_j})
\left[
\GTenToneHL{cost}{c}(\GTenToneHL{alternative}{A_i})+\GTenToneHL{cost}{c}(\GTenToneHL{state}{S_j})
\right].
\]

Therefore,
\[
\GTenToneHL{payoff}{P}(\GTenToneHL{alternative}{A_i},\GTenToneHL{state}{S_j})
=
\left[
\GTenToneHL{quantity}{q}(\GTenToneHL{alternative}{A_i},\GTenToneHL{state}{S_j})\GTenToneHL{revenue}{r}(\GTenToneHL{state}{S_j})+\GTenToneHL{revenue}{f_r}(\GTenToneHL{state}{S_j})
\right]
-
\GTenToneHL{quantity}{q}(\GTenToneHL{alternative}{A_i},\GTenToneHL{state}{S_j})
\left[
\GTenToneHL{cost}{c}(\GTenToneHL{alternative}{A_i})+\GTenToneHL{cost}{c}(\GTenToneHL{state}{S_j})
\right].
\]

\secInicio[sec:problem13-urban-delivery-fleet-c1-c3][sec:two-component-models]

\subsubsection*{C3 - Calculating profit and building the payoff table.}
\phantomsection
\label{subsec:problem13-urban-delivery-fleet-c3}

For every \GTenToneHL{alternative}{alternative}--state pair, calculate
\[
\begin{aligned}
	\GTenToneHL{revenue}{R}(\GTenToneHL{alternative}{A_i},\GTenToneHL{state}{S_j})
	&=
	\GTenToneHL{quantity}{q}(\GTenToneHL{alternative}{A_i},\GTenToneHL{state}{S_j})\GTenToneHL{revenue}{r}(\GTenToneHL{state}{S_j})+\GTenToneHL{revenue}{f_r}(\GTenToneHL{state}{S_j}),
	\\
	\GTenToneHL{cost}{C}(\GTenToneHL{alternative}{A_i},\GTenToneHL{state}{S_j})
	&=
	\GTenToneHL{quantity}{q}(\GTenToneHL{alternative}{A_i},\GTenToneHL{state}{S_j})
	\left[
	\GTenToneHL{cost}{c}(\GTenToneHL{alternative}{A_i})+\GTenToneHL{cost}{c}(\GTenToneHL{state}{S_j})
	\right],
	\\
	\GTenToneHL{payoff}{P}(\GTenToneHL{alternative}{A_i},\GTenToneHL{state}{S_j})
	&=
	\GTenToneHL{revenue}{R}(\GTenToneHL{alternative}{A_i},\GTenToneHL{state}{S_j})-\GTenToneHL{cost}{C}(\GTenToneHL{alternative}{A_i},\GTenToneHL{state}{S_j}).
\end{aligned}
\]

\begin{center}
	\small
	\begin{tabular}{llll}
		\toprule
		Cell
		& \(\GTenToneHL{revenue}{R}(\GTenToneHL{alternative}{A_i},\GTenToneHL{state}{S_j})\)
		& \(\GTenToneHL{cost}{C}(\GTenToneHL{alternative}{A_i},\GTenToneHL{state}{S_j})\)
		& \(\GTenToneHL{payoff}{P}(\GTenToneHL{alternative}{A_i},\GTenToneHL{state}{S_j})=\GTenToneHL{revenue}{R}(\GTenToneHL{alternative}{A_i},\GTenToneHL{state}{S_j})-\GTenToneHL{cost}{C}(\GTenToneHL{alternative}{A_i},\GTenToneHL{state}{S_j})\) \\
		\midrule
		\((\GTenToneHL{alternative}{A_1},\GTenToneHL{state}{S_1})\)
		& \(9(8)+14=86\)
		& \(9(3+1)=36\)
		& \(86-36=50\) \\
		
		\((\GTenToneHL{alternative}{A_1},\GTenToneHL{state}{S_2})\)
		& \(5(6)+8=38\)
		& \(5(3+3)=30\)
		& \(38-30=8\) \\
		
		\((\GTenToneHL{alternative}{A_2},\GTenToneHL{state}{S_1})\)
		& \(12(8)+14=110\)
		& \(12(2+1)=36\)
		& \(110-36=74\) \\
		
		\((\GTenToneHL{alternative}{A_2},\GTenToneHL{state}{S_2})\)
		& \(6(6)+8=44\)
		& \(6(2+3)=30\)
		& \(44-30=14\) \\
		\bottomrule
	\end{tabular}
\end{center}

The resulting \GTenToneHL{payoff}{payoff} table is:

\begin{center}
	\begin{tabular}{lrr}
		\toprule
		Alternative & High demand, light traffic & Low demand, heavy traffic \\
		\midrule
		Conventional van fleet & \GTenToneHL{payoff}{50} & \GTenToneHL{payoff}{8} \\
		Electric cargo-bike fleet & \GTenToneHL{payoff}{74} & \GTenToneHL{payoff}{14} \\
		\bottomrule
	\end{tabular}
\end{center}

\secInicio[sec:problem13-urban-delivery-fleet-c1-c3][sec:two-component-models]

% ----------------------------------------------------------------
\subsection{Data-Service Contract (C1--C3)}
\label{sec:problem14-data-service-contract-c1-c3}

A data-services company must choose a processing method before the requirements
of a new contract are confirmed. \GTenToneHL{revenue}{Revenue} consists of a fixed payment determined
by the selected method and a fixed service-level payment determined by the
contract state. \GTenToneHL{cost}{Cost} consists of a method-based processing rate and a fixed
implementation \GTenToneHL{cost}{cost} determined by the selected method.

\begin{multicols}{2}
	The company is comparing two \GTenToneHL{alternative}{alternatives}:
	\begin{enumerate}[label=\arabic*.]
		\item \(\GTenToneHL{alternative}{A_1}\): Use standard data processing
		\item \(\GTenToneHL{alternative}{A_2}\): Use automated data processing
	\end{enumerate}
	
	The possible \GTenToneHL{state}{states} \GTenToneHL{state}{of} \GTenToneHL{state}{nature} are:
	\begin{enumerate}[label=\arabic*.]
		\item \(\GTenToneHL{state}{S_1}\): Standard-volume contract
		\item \(\GTenToneHL{state}{S_2}\): High-volume contract
	\end{enumerate}
	
	\columnbreak
	
	The quantities, measured in processing batches, are
	\[
	\begin{aligned}
		\GTenToneHL{quantity}{q}(\GTenToneHL{alternative}{A_1},\GTenToneHL{state}{S_1})&=\GTenToneHL{quantity}{4},
		&
		\GTenToneHL{quantity}{q}(\GTenToneHL{alternative}{A_1},\GTenToneHL{state}{S_2})&=\GTenToneHL{quantity}{7},
		\\
		\GTenToneHL{quantity}{q}(\GTenToneHL{alternative}{A_2},\GTenToneHL{state}{S_1})&=\GTenToneHL{quantity}{5},
		&
		\GTenToneHL{quantity}{q}(\GTenToneHL{alternative}{A_2},\GTenToneHL{state}{S_2})&=9.
	\end{aligned}
	\]
	
	The \GTenToneHL{alternative}{alternative}-based fixed revenues, in thousand USD, are
	\[
	\GTenToneHL{revenue}{f_r}(\GTenToneHL{alternative}{A_1})=\GTenToneHL{revenue}{55},
	\qquad
	\GTenToneHL{revenue}{f_r}(\GTenToneHL{alternative}{A_2})=72.
	\]
	
	The state-based fixed revenues, in thousand USD, are
	\[
	\GTenToneHL{revenue}{f_r}(\GTenToneHL{state}{S_1})=\GTenToneHL{revenue}{16},
	\qquad
	\GTenToneHL{revenue}{f_r}(\GTenToneHL{state}{S_2})=9.
	\]
	
	The \GTenToneHL{alternative}{alternative}-based \GTenToneHL{cost}{cost} rates, in thousand USD per batch, are
	\[
	\GTenToneHL{cost}{c}(\GTenToneHL{alternative}{A_1})=\GTenToneHL{cost}{6},
	\qquad
	\GTenToneHL{cost}{c}(\GTenToneHL{alternative}{A_2})=4.
	\]
	
	The \GTenToneHL{alternative}{alternative}-based fixed \GTenToneHL{cost}{costs}, in thousand USD, are
	\[
	\GTenToneHL{cost}{f_c}(\GTenToneHL{alternative}{A_1})=\GTenToneHL{cost}{18},
	\qquad
	\GTenToneHL{cost}{f_c}(\GTenToneHL{alternative}{A_2})=31.
	\]
\end{multicols}

\textbf{Questions/Tasks.}
\begin{enumerate}
	\CriterionTask
	{subsec:problem14-data-service-contract-c1}{1}
	{Interpreting decision alternatives, events, consequences, and states.}

	\CriterionTask
	{subsec:problem14-data-service-contract-c3}{3}
	{Calculating profit and building the payoff table.}
\end{enumerate}

\subsubsection*{C1 - Interpreting decision alternatives, events, consequences, and states.}
\phantomsection
\label{subsec:problem14-data-service-contract-c1}

\DecisionElementsTable
{Use standard data processing; Use automated data processing}
{Standard-volume contract; High-volume contract}
{Contract requirements confirmed after the processing-method decision}
{Profit in thousand USD for each processing-method--contract-state pair}

\GTenToneHL{revenue}{Revenue} contains \GTenToneHL{alternative}{alternative}-based and state-based fixed components:
\[
\GTenToneHL{revenue}{R}(\GTenToneHL{alternative}{A_i},\GTenToneHL{state}{S_j})
=
\GTenToneHL{revenue}{f_r}(\GTenToneHL{alternative}{A_i})+\GTenToneHL{revenue}{f_r}(\GTenToneHL{state}{S_j}).
\]

\GTenToneHL{cost}{Cost} contains \GTenToneHL{alternative}{alternative}-based variable and fixed components:
\[
\GTenToneHL{cost}{C}(\GTenToneHL{alternative}{A_i},\GTenToneHL{state}{S_j})
=
\GTenToneHL{quantity}{q}(\GTenToneHL{alternative}{A_i},\GTenToneHL{state}{S_j})\GTenToneHL{cost}{c}(\GTenToneHL{alternative}{A_i})+\GTenToneHL{cost}{f_c}(\GTenToneHL{alternative}{A_i}).
\]

Therefore,
\[
\GTenToneHL{payoff}{P}(\GTenToneHL{alternative}{A_i},\GTenToneHL{state}{S_j})
=
\left[
\GTenToneHL{revenue}{f_r}(\GTenToneHL{alternative}{A_i})+\GTenToneHL{revenue}{f_r}(\GTenToneHL{state}{S_j})
\right]
-
\left[
\GTenToneHL{quantity}{q}(\GTenToneHL{alternative}{A_i},\GTenToneHL{state}{S_j})\GTenToneHL{cost}{c}(\GTenToneHL{alternative}{A_i})+\GTenToneHL{cost}{f_c}(\GTenToneHL{alternative}{A_i})
\right].
\]

\secInicio[sec:problem14-data-service-contract-c1-c3][sec:two-component-models]

\subsubsection*{C3 - Calculating profit and building the payoff table.}
\phantomsection
\label{subsec:problem14-data-service-contract-c3}

For every \GTenToneHL{alternative}{alternative}--state pair, calculate
\[
\begin{aligned}
	\GTenToneHL{revenue}{R}(\GTenToneHL{alternative}{A_i},\GTenToneHL{state}{S_j})
	&=
	\GTenToneHL{revenue}{f_r}(\GTenToneHL{alternative}{A_i})+\GTenToneHL{revenue}{f_r}(\GTenToneHL{state}{S_j}),
	\\
	\GTenToneHL{cost}{C}(\GTenToneHL{alternative}{A_i},\GTenToneHL{state}{S_j})
	&=
	\GTenToneHL{quantity}{q}(\GTenToneHL{alternative}{A_i},\GTenToneHL{state}{S_j})\GTenToneHL{cost}{c}(\GTenToneHL{alternative}{A_i})+\GTenToneHL{cost}{f_c}(\GTenToneHL{alternative}{A_i}),
	\\
	\GTenToneHL{payoff}{P}(\GTenToneHL{alternative}{A_i},\GTenToneHL{state}{S_j})
	&=
	\GTenToneHL{revenue}{R}(\GTenToneHL{alternative}{A_i},\GTenToneHL{state}{S_j})-\GTenToneHL{cost}{C}(\GTenToneHL{alternative}{A_i},\GTenToneHL{state}{S_j}).
\end{aligned}
\]

\begin{center}
	\small
	\begin{tabular}{llll}
		\toprule
		Cell
		& \(\GTenToneHL{revenue}{R}(\GTenToneHL{alternative}{A_i},\GTenToneHL{state}{S_j})\)
		& \(\GTenToneHL{cost}{C}(\GTenToneHL{alternative}{A_i},\GTenToneHL{state}{S_j})\)
		& \(\GTenToneHL{payoff}{P}(\GTenToneHL{alternative}{A_i},\GTenToneHL{state}{S_j})=\GTenToneHL{revenue}{R}(\GTenToneHL{alternative}{A_i},\GTenToneHL{state}{S_j})-\GTenToneHL{cost}{C}(\GTenToneHL{alternative}{A_i},\GTenToneHL{state}{S_j})\) \\
		\midrule
		\((\GTenToneHL{alternative}{A_1},\GTenToneHL{state}{S_1})\)
		& \(55+16=71\)
		& \(4(6)+18=42\)
		& \(71-42=29\) \\
		
		\((\GTenToneHL{alternative}{A_1},\GTenToneHL{state}{S_2})\)
		& \(55+9=64\)
		& \(7(6)+18=60\)
		& \(64-60=4\) \\
		
		\((\GTenToneHL{alternative}{A_2},\GTenToneHL{state}{S_1})\)
		& \(72+16=88\)
		& \(5(4)+31=51\)
		& \(88-51=37\) \\
		
		\((\GTenToneHL{alternative}{A_2},\GTenToneHL{state}{S_2})\)
		& \(72+9=81\)
		& \(9(4)+31=67\)
		& \(81-67=14\) \\
		\bottomrule
	\end{tabular}
\end{center}

The resulting \GTenToneHL{payoff}{payoff} table is:

\begin{center}
	\begin{tabular}{lrr}
		\toprule
		Alternative & Standard-volume contract & High-volume contract \\
		\midrule
		Standard data processing & \GTenToneHL{payoff}{29} & \GTenToneHL{payoff}{4} \\
		Automated data processing & \GTenToneHL{payoff}{37} & \GTenToneHL{payoff}{14} \\
		\bottomrule
	\end{tabular}
\end{center}

\secInicio[sec:problem14-data-service-contract-c1-c3][sec:two-component-models]

% ================================================================
\section{Three-Component Models}
\label{sec:three-component-models}

\SectionProblemsTableOfContents{
	\SectionProblemLink{sec:problem15-regional-parcel-delivery-c1-c3}{Regional Parcel Delivery (C1--C3)} \\
	\SectionProblemLink{sec:problem16-solar-installation-programme-c1-c3}{Solar Installation Programme (C1--C3)} \\
	\SectionProblemLink{sec:problem17-hybrid-conference-production-c1-c3}{Hybrid Conference Production (C1--C3)} \\
	\SectionProblemLink{sec:problem18-cold-storage-export-contract-c1-c3}{Cold-Storage Export Contract (C1--C3)}
}

\vspace{0.8cm}

In a three-component model, \GTenToneHL{revenue}{revenue} contains exactly three nonzero
components and \GTenToneHL{cost}{cost} contains exactly three nonzero components. Each model
is constructed independently, so the \GTenToneHL{revenue}{revenue} and \GTenToneHL{cost}{cost} structures do not
need to match.

The four paired patterns in this section are:

\begin{center}
	\small
	\begin{tabular}{cll}
		\toprule
		Problem & \GTenToneHL{revenue}{Revenue} model & \GTenToneHL{cost}{Cost} model \\
		\midrule
		15
		& \(\GTenToneHL{revenue}{R}(\GTenToneHL{alternative}{A_i},\GTenToneHL{state}{S_j})
		=\GTenToneHL{quantity}{q}(\GTenToneHL{alternative}{A_i},\GTenToneHL{state}{S_j})\bigl[\GTenToneHL{revenue}{r}(\GTenToneHL{alternative}{A_i})+\GTenToneHL{revenue}{r}(\GTenToneHL{state}{S_j})\bigr]
		+\GTenToneHL{revenue}{f_r}(\GTenToneHL{alternative}{A_i})\)
		& \(\GTenToneHL{cost}{C}(\GTenToneHL{alternative}{A_i},\GTenToneHL{state}{S_j})
		=\GTenToneHL{quantity}{q}(\GTenToneHL{alternative}{A_i},\GTenToneHL{state}{S_j})\bigl[\GTenToneHL{cost}{c}(\GTenToneHL{alternative}{A_i})+\GTenToneHL{cost}{c}(\GTenToneHL{state}{S_j})\bigr]
		+\GTenToneHL{cost}{f_c}(\GTenToneHL{state}{S_j})\) \\
		
		16
		& \(\GTenToneHL{revenue}{R}(\GTenToneHL{alternative}{A_i},\GTenToneHL{state}{S_j})
		=\GTenToneHL{quantity}{q}(\GTenToneHL{alternative}{A_i},\GTenToneHL{state}{S_j})\bigl[\GTenToneHL{revenue}{r}(\GTenToneHL{alternative}{A_i})+\GTenToneHL{revenue}{r}(\GTenToneHL{state}{S_j})\bigr]
		+\GTenToneHL{revenue}{f_r}(\GTenToneHL{state}{S_j})\)
		& \(\GTenToneHL{cost}{C}(\GTenToneHL{alternative}{A_i},\GTenToneHL{state}{S_j})
		=\GTenToneHL{quantity}{q}(\GTenToneHL{alternative}{A_i},\GTenToneHL{state}{S_j})\GTenToneHL{cost}{c}(\GTenToneHL{alternative}{A_i})
		+\GTenToneHL{cost}{f_c}(\GTenToneHL{alternative}{A_i})+\GTenToneHL{cost}{f_c}(\GTenToneHL{state}{S_j})\) \\
		
		17
		& \(\GTenToneHL{revenue}{R}(\GTenToneHL{alternative}{A_i},\GTenToneHL{state}{S_j})
		=\GTenToneHL{quantity}{q}(\GTenToneHL{alternative}{A_i},\GTenToneHL{state}{S_j})\GTenToneHL{revenue}{r}(\GTenToneHL{alternative}{A_i})
		+\GTenToneHL{revenue}{f_r}(\GTenToneHL{alternative}{A_i})+\GTenToneHL{revenue}{f_r}(\GTenToneHL{state}{S_j})\)
		& \(\GTenToneHL{cost}{C}(\GTenToneHL{alternative}{A_i},\GTenToneHL{state}{S_j})
		=\GTenToneHL{quantity}{q}(\GTenToneHL{alternative}{A_i},\GTenToneHL{state}{S_j})\GTenToneHL{cost}{c}(\GTenToneHL{state}{S_j})
		+\GTenToneHL{cost}{f_c}(\GTenToneHL{alternative}{A_i})+\GTenToneHL{cost}{f_c}(\GTenToneHL{state}{S_j})\) \\
		
		18
		& \(\GTenToneHL{revenue}{R}(\GTenToneHL{alternative}{A_i},\GTenToneHL{state}{S_j})
		=\GTenToneHL{quantity}{q}(\GTenToneHL{alternative}{A_i},\GTenToneHL{state}{S_j})\GTenToneHL{revenue}{r}(\GTenToneHL{state}{S_j})
		+\GTenToneHL{revenue}{f_r}(\GTenToneHL{alternative}{A_i})+\GTenToneHL{revenue}{f_r}(\GTenToneHL{state}{S_j})\)
		& \(\GTenToneHL{cost}{C}(\GTenToneHL{alternative}{A_i},\GTenToneHL{state}{S_j})
		=\GTenToneHL{quantity}{q}(\GTenToneHL{alternative}{A_i},\GTenToneHL{state}{S_j})\bigl[\GTenToneHL{cost}{c}(\GTenToneHL{alternative}{A_i})+\GTenToneHL{cost}{c}(\GTenToneHL{state}{S_j})\bigr]
		+\GTenToneHL{cost}{f_c}(\GTenToneHL{alternative}{A_i})\) \\
		\bottomrule
	\end{tabular}
\end{center}

Thus, all four possible three-component structures---AV+SV+AF,
AV+SV+SF, AV+AF+SF, and SV+AF+SF---appear exactly once as a
\GTenToneHL{revenue}{revenue} structure and exactly once as a \GTenToneHL{cost}{cost} structure.

% ----------------------------------------------------------------
\subsection{Regional Parcel Delivery (C1--C3)}
\label{sec:problem15-regional-parcel-delivery-c1-c3}

A regional logistics company must select a parcel-delivery plan before demand
and traffic conditions are known. \GTenToneHL{revenue}{Revenue} per delivery batch combines a rate
determined by the selected plan and a rate determined by the state. The
selected plan also provides a fixed service payment. \GTenToneHL{cost}{Cost} per batch combines
a plan-based operating rate and a state-based traffic rate, together with a
fixed distribution \GTenToneHL{cost}{cost} determined by the state.

\begin{multicols}{2}
	The company is comparing two \GTenToneHL{alternative}{alternatives}:
	\begin{enumerate}[label=\arabic*.]
		\item \(\GTenToneHL{alternative}{A_1}\): Use the standard delivery plan
		\item \(\GTenToneHL{alternative}{A_2}\): Use the express delivery plan
	\end{enumerate}
	
	The possible \GTenToneHL{state}{states} \GTenToneHL{state}{of} \GTenToneHL{state}{nature} are:
	\begin{enumerate}[label=\arabic*.]
		\item \(\GTenToneHL{state}{S_1}\): High demand with light traffic
		\item \(\GTenToneHL{state}{S_2}\): Moderate demand with heavy traffic
	\end{enumerate}
	
	\columnbreak
	
	The quantities, measured in hundreds of delivery batches, are
	\[
	\begin{aligned}
		\GTenToneHL{quantity}{q}(\GTenToneHL{alternative}{A_1},\GTenToneHL{state}{S_1})&=\GTenToneHL{quantity}{8},
		&
		\GTenToneHL{quantity}{q}(\GTenToneHL{alternative}{A_1},\GTenToneHL{state}{S_2})&=\GTenToneHL{quantity}{5},
		\\
		\GTenToneHL{quantity}{q}(\GTenToneHL{alternative}{A_2},\GTenToneHL{state}{S_1})&=\GTenToneHL{quantity}{10},
		&
		\GTenToneHL{quantity}{q}(\GTenToneHL{alternative}{A_2},\GTenToneHL{state}{S_2})&=6.
	\end{aligned}
	\]
	
	The \GTenToneHL{alternative}{alternative}-based \GTenToneHL{revenue}{revenue} rates, in thousand USD per hundred
	delivery batches, are
	\[
	\GTenToneHL{revenue}{r}(\GTenToneHL{alternative}{A_1})=\GTenToneHL{revenue}{9},
	\qquad
	\GTenToneHL{revenue}{r}(\GTenToneHL{alternative}{A_2})=11.
	\]
	
	The state-based \GTenToneHL{revenue}{revenue} rates are
	\[
	\GTenToneHL{revenue}{r}(\GTenToneHL{state}{S_1})=\GTenToneHL{revenue}{3},
	\qquad
	\GTenToneHL{revenue}{r}(\GTenToneHL{state}{S_2})=2.
	\]
	
	The \GTenToneHL{alternative}{alternative}-based fixed revenues, in thousand USD, are
	\[
	\GTenToneHL{revenue}{f_r}(\GTenToneHL{alternative}{A_1})=\GTenToneHL{revenue}{16},
	\qquad
	\GTenToneHL{revenue}{f_r}(\GTenToneHL{alternative}{A_2})=24.
	\]
	
	The \GTenToneHL{alternative}{alternative}-based \GTenToneHL{cost}{cost} rates are
	\[
	\GTenToneHL{cost}{c}(\GTenToneHL{alternative}{A_1})=\GTenToneHL{cost}{4},
	\qquad
	\GTenToneHL{cost}{c}(\GTenToneHL{alternative}{A_2})=5.
	\]
	
	The state-based \GTenToneHL{cost}{cost} rates are
	\[
	\GTenToneHL{cost}{c}(\GTenToneHL{state}{S_1})=\GTenToneHL{cost}{1},
	\qquad
	\GTenToneHL{cost}{c}(\GTenToneHL{state}{S_2})=3.
	\]
	
	The state-based fixed \GTenToneHL{cost}{costs}, in thousand USD, are
	\[
	\GTenToneHL{cost}{f_c}(\GTenToneHL{state}{S_1})=\GTenToneHL{cost}{10},
	\qquad
	\GTenToneHL{cost}{f_c}(\GTenToneHL{state}{S_2})=15.
	\]
\end{multicols}

\textbf{Questions/Tasks.}
\begin{enumerate}
	\CriterionTask
	{subsec:problem15-regional-parcel-delivery-c1}{1}
	{Interpreting decision alternatives, events, consequences, and states.}

	\CriterionTask
	{subsec:problem15-regional-parcel-delivery-c3}{3}
	{Calculating profit and building the payoff table.}
\end{enumerate}

\subsubsection*{C1 - Interpreting decision alternatives, events, consequences, and states.}
\phantomsection
\label{subsec:problem15-regional-parcel-delivery-c1}

\DecisionElementsTable
{Use the standard delivery plan; Use the express delivery plan}
{High demand with light traffic; Moderate demand with heavy traffic}
{Demand and traffic conditions revealed after the delivery-plan decision}
{Profit in thousand USD for each delivery-plan--state pair}

\GTenToneHL{revenue}{Revenue} contains \GTenToneHL{alternative}{alternative}-based and state-based variable components and
an \GTenToneHL{alternative}{alternative}-based fixed component:
\[
\GTenToneHL{revenue}{R}(\GTenToneHL{alternative}{A_i},\GTenToneHL{state}{S_j})
=
\GTenToneHL{quantity}{q}(\GTenToneHL{alternative}{A_i},\GTenToneHL{state}{S_j})
\left[
\GTenToneHL{revenue}{r}(\GTenToneHL{alternative}{A_i})+\GTenToneHL{revenue}{r}(\GTenToneHL{state}{S_j})
\right]
+
\GTenToneHL{revenue}{f_r}(\GTenToneHL{alternative}{A_i}).
\]

\GTenToneHL{cost}{Cost} contains \GTenToneHL{alternative}{alternative}-based and state-based variable components and a
state-based fixed component:
\[
\GTenToneHL{cost}{C}(\GTenToneHL{alternative}{A_i},\GTenToneHL{state}{S_j})
=
\GTenToneHL{quantity}{q}(\GTenToneHL{alternative}{A_i},\GTenToneHL{state}{S_j})
\left[
\GTenToneHL{cost}{c}(\GTenToneHL{alternative}{A_i})+\GTenToneHL{cost}{c}(\GTenToneHL{state}{S_j})
\right]
+
\GTenToneHL{cost}{f_c}(\GTenToneHL{state}{S_j}).
\]

Therefore,
\[
\begin{aligned}
	\GTenToneHL{payoff}{P}(\GTenToneHL{alternative}{A_i},\GTenToneHL{state}{S_j})
	={}&
	\left\{
	\GTenToneHL{quantity}{q}(\GTenToneHL{alternative}{A_i},\GTenToneHL{state}{S_j})
	\left[
	\GTenToneHL{revenue}{r}(\GTenToneHL{alternative}{A_i})+\GTenToneHL{revenue}{r}(\GTenToneHL{state}{S_j})
	\right]
	+
	\GTenToneHL{revenue}{f_r}(\GTenToneHL{alternative}{A_i})
	\right\}
	\\
	&-
	\left\{
	\GTenToneHL{quantity}{q}(\GTenToneHL{alternative}{A_i},\GTenToneHL{state}{S_j})
	\left[
	\GTenToneHL{cost}{c}(\GTenToneHL{alternative}{A_i})+\GTenToneHL{cost}{c}(\GTenToneHL{state}{S_j})
	\right]
	+
	\GTenToneHL{cost}{f_c}(\GTenToneHL{state}{S_j})
	\right\}.
\end{aligned}
\]

\secInicio[sec:problem15-regional-parcel-delivery-c1-c3][sec:three-component-models]

\subsubsection*{C3 - Calculating profit and building the payoff table.}
\phantomsection
\label{subsec:problem15-regional-parcel-delivery-c3}

For every \GTenToneHL{alternative}{alternative}--state pair, calculate
\[
\begin{aligned}
	\GTenToneHL{revenue}{R}(\GTenToneHL{alternative}{A_i},\GTenToneHL{state}{S_j})
	&=
	\GTenToneHL{quantity}{q}(\GTenToneHL{alternative}{A_i},\GTenToneHL{state}{S_j})
	\left[
	\GTenToneHL{revenue}{r}(\GTenToneHL{alternative}{A_i})+\GTenToneHL{revenue}{r}(\GTenToneHL{state}{S_j})
	\right]
	+
	\GTenToneHL{revenue}{f_r}(\GTenToneHL{alternative}{A_i}),
	\\
	\GTenToneHL{cost}{C}(\GTenToneHL{alternative}{A_i},\GTenToneHL{state}{S_j})
	&=
	\GTenToneHL{quantity}{q}(\GTenToneHL{alternative}{A_i},\GTenToneHL{state}{S_j})
	\left[
	\GTenToneHL{cost}{c}(\GTenToneHL{alternative}{A_i})+\GTenToneHL{cost}{c}(\GTenToneHL{state}{S_j})
	\right]
	+
	\GTenToneHL{cost}{f_c}(\GTenToneHL{state}{S_j}),
	\\
	\GTenToneHL{payoff}{P}(\GTenToneHL{alternative}{A_i},\GTenToneHL{state}{S_j})
	&=
	\GTenToneHL{revenue}{R}(\GTenToneHL{alternative}{A_i},\GTenToneHL{state}{S_j})-\GTenToneHL{cost}{C}(\GTenToneHL{alternative}{A_i},\GTenToneHL{state}{S_j}).
\end{aligned}
\]

\begin{center}
	\small
	\begin{tabular}{llll}
		\toprule
		Cell
		& \(\GTenToneHL{revenue}{R}(\GTenToneHL{alternative}{A_i},\GTenToneHL{state}{S_j})\)
		& \(\GTenToneHL{cost}{C}(\GTenToneHL{alternative}{A_i},\GTenToneHL{state}{S_j})\)
		& \(\GTenToneHL{payoff}{P}(\GTenToneHL{alternative}{A_i},\GTenToneHL{state}{S_j})=\GTenToneHL{revenue}{R}(\GTenToneHL{alternative}{A_i},\GTenToneHL{state}{S_j})-\GTenToneHL{cost}{C}(\GTenToneHL{alternative}{A_i},\GTenToneHL{state}{S_j})\) \\
		\midrule
		\((\GTenToneHL{alternative}{A_1},\GTenToneHL{state}{S_1})\)
		& \(8(9+3)+16=112\)
		& \(8(4+1)+10=50\)
		& \(112-50=62\) \\
		
		\((\GTenToneHL{alternative}{A_1},\GTenToneHL{state}{S_2})\)
		& \(5(9+2)+16=71\)
		& \(5(4+3)+15=50\)
		& \(71-50=21\) \\
		
		\((\GTenToneHL{alternative}{A_2},\GTenToneHL{state}{S_1})\)
		& \(10(11+3)+24=164\)
		& \(10(5+1)+10=70\)
		& \(164-70=94\) \\
		
		\((\GTenToneHL{alternative}{A_2},\GTenToneHL{state}{S_2})\)
		& \(6(11+2)+24=102\)
		& \(6(5+3)+15=63\)
		& \(102-63=39\) \\
		\bottomrule
	\end{tabular}
\end{center}

The resulting \GTenToneHL{payoff}{payoff} table is:

\begin{center}
	\begin{tabular}{lrr}
		\toprule
		Alternative
		& High demand, light traffic
		& Moderate demand, heavy traffic \\
		\midrule
		Standard delivery plan & \GTenToneHL{payoff}{62} & \GTenToneHL{rate}{21} \\
		Express delivery plan & \GTenToneHL{payoff}{94} & \GTenToneHL{rate}{39} \\
		\bottomrule
	\end{tabular}
\end{center}

\secInicio[sec:problem15-regional-parcel-delivery-c1-c3][sec:three-component-models]

% ----------------------------------------------------------------
\subsection{Solar Installation Programme (C1--C3)}
\label{sec:problem16-solar-installation-programme-c1-c3}

A renewable-energy company must select a solar installation programme before
customer demand and permit conditions are known. \GTenToneHL{revenue}{Revenue} per installation
batch combines a programme-based rate and a state-based market rate, together
with a fixed incentive determined by the state. Costs consist of a
programme-based variable \GTenToneHL{cost}{cost}, a fixed setup \GTenToneHL{cost}{cost} determined by the programme,
and a fixed permit \GTenToneHL{cost}{cost} determined by the state.

\begin{multicols}{2}
	The company is comparing two \GTenToneHL{alternative}{alternatives}:
	\begin{enumerate}[label=\arabic*.]
		\item \(\GTenToneHL{alternative}{A_1}\): Offer the standard-panel programme
		\item \(\GTenToneHL{alternative}{A_2}\): Offer the high-efficiency-panel programme
	\end{enumerate}
	
	The possible \GTenToneHL{state}{states} \GTenToneHL{state}{of} \GTenToneHL{state}{nature} are:
	\begin{enumerate}[label=\arabic*.]
		\item \(\GTenToneHL{state}{S_1}\): Strong commercial demand
		\item \(\GTenToneHL{state}{S_2}\): Limited residential demand
	\end{enumerate}
	
	\columnbreak
	
	The quantities, measured in installation batches, are
	\[
	\begin{aligned}
		\GTenToneHL{quantity}{q}(\GTenToneHL{alternative}{A_1},\GTenToneHL{state}{S_1})&=\GTenToneHL{quantity}{7},
		&
		\GTenToneHL{quantity}{q}(\GTenToneHL{alternative}{A_1},\GTenToneHL{state}{S_2})&=\GTenToneHL{quantity}{4},
		\\
		\GTenToneHL{quantity}{q}(\GTenToneHL{alternative}{A_2},\GTenToneHL{state}{S_1})&=\GTenToneHL{quantity}{6},
		&
		\GTenToneHL{quantity}{q}(\GTenToneHL{alternative}{A_2},\GTenToneHL{state}{S_2})&=3.
	\end{aligned}
	\]
	
	The \GTenToneHL{alternative}{alternative}-based \GTenToneHL{revenue}{revenue} rates, in thousand USD per installation
	batch, are
	\[
	\GTenToneHL{revenue}{r}(\GTenToneHL{alternative}{A_1})=\GTenToneHL{revenue}{12},
	\qquad
	\GTenToneHL{revenue}{r}(\GTenToneHL{alternative}{A_2})=18.
	\]
	
	The state-based \GTenToneHL{revenue}{revenue} rates are
	\[
	\GTenToneHL{revenue}{r}(\GTenToneHL{state}{S_1})=\GTenToneHL{revenue}{4},
	\qquad
	\GTenToneHL{revenue}{r}(\GTenToneHL{state}{S_2})=2.
	\]
	
	The state-based fixed revenues, in thousand USD, are
	\[
	\GTenToneHL{revenue}{f_r}(\GTenToneHL{state}{S_1})=\GTenToneHL{revenue}{20},
	\qquad
	\GTenToneHL{revenue}{f_r}(\GTenToneHL{state}{S_2})=11.
	\]
	
	The \GTenToneHL{alternative}{alternative}-based \GTenToneHL{cost}{cost} rates, in thousand USD per installation
	batch, are
	\[
	\GTenToneHL{cost}{c}(\GTenToneHL{alternative}{A_1})=\GTenToneHL{cost}{6},
	\qquad
	\GTenToneHL{cost}{c}(\GTenToneHL{alternative}{A_2})=8.
	\]
	
	The \GTenToneHL{alternative}{alternative}-based fixed \GTenToneHL{cost}{costs}, in thousand USD, are
	\[
	\GTenToneHL{cost}{f_c}(\GTenToneHL{alternative}{A_1})=\GTenToneHL{cost}{18},
	\qquad
	\GTenToneHL{cost}{f_c}(\GTenToneHL{alternative}{A_2})=25.
	\]
	
	The state-based fixed \GTenToneHL{cost}{costs}, in thousand USD, are
	\[
	\GTenToneHL{cost}{f_c}(\GTenToneHL{state}{S_1})=\GTenToneHL{cost}{9},
	\qquad
	\GTenToneHL{cost}{f_c}(\GTenToneHL{state}{S_2})=14.
	\]
\end{multicols}

\textbf{Questions/Tasks.}
\begin{enumerate}
	\CriterionTask
	{subsec:problem16-solar-installation-programme-c1}{1}
	{Interpreting decision alternatives, events, consequences, and states.}

	\CriterionTask
	{subsec:problem16-solar-installation-programme-c3}{3}
	{Calculating profit and building the payoff table.}
\end{enumerate}

\subsubsection*{C1 - Interpreting decision alternatives, events, consequences, and states.}
\phantomsection
\label{subsec:problem16-solar-installation-programme-c1}

\DecisionElementsTable
{Offer the standard-panel programme; Offer the high-efficiency-panel programme}
{Strong commercial demand; Limited residential demand}
{Customer demand and permit conditions revealed after the programme decision}
{Profit in thousand USD for each programme--market-state pair}

\GTenToneHL{revenue}{Revenue} contains \GTenToneHL{alternative}{alternative}-based and state-based variable components and
a state-based fixed component:
\[
\GTenToneHL{revenue}{R}(\GTenToneHL{alternative}{A_i},\GTenToneHL{state}{S_j})
=
\GTenToneHL{quantity}{q}(\GTenToneHL{alternative}{A_i},\GTenToneHL{state}{S_j})
\left[
\GTenToneHL{revenue}{r}(\GTenToneHL{alternative}{A_i})+\GTenToneHL{revenue}{r}(\GTenToneHL{state}{S_j})
\right]
+
\GTenToneHL{revenue}{f_r}(\GTenToneHL{state}{S_j}).
\]

\GTenToneHL{cost}{Cost} contains an \GTenToneHL{alternative}{alternative}-based variable component and \GTenToneHL{alternative}{alternative}-based
and state-based fixed components:
\[
\GTenToneHL{cost}{C}(\GTenToneHL{alternative}{A_i},\GTenToneHL{state}{S_j})
=
\GTenToneHL{quantity}{q}(\GTenToneHL{alternative}{A_i},\GTenToneHL{state}{S_j})\GTenToneHL{cost}{c}(\GTenToneHL{alternative}{A_i})
+
\GTenToneHL{cost}{f_c}(\GTenToneHL{alternative}{A_i})
+
\GTenToneHL{cost}{f_c}(\GTenToneHL{state}{S_j}).
\]

Therefore,
\[
\begin{aligned}
	\GTenToneHL{payoff}{P}(\GTenToneHL{alternative}{A_i},\GTenToneHL{state}{S_j})
	={}&
	\left\{
	\GTenToneHL{quantity}{q}(\GTenToneHL{alternative}{A_i},\GTenToneHL{state}{S_j})
	\left[
	\GTenToneHL{revenue}{r}(\GTenToneHL{alternative}{A_i})+\GTenToneHL{revenue}{r}(\GTenToneHL{state}{S_j})
	\right]
	+
	\GTenToneHL{revenue}{f_r}(\GTenToneHL{state}{S_j})
	\right\}
	\\
	&-
	\left\{
	\GTenToneHL{quantity}{q}(\GTenToneHL{alternative}{A_i},\GTenToneHL{state}{S_j})\GTenToneHL{cost}{c}(\GTenToneHL{alternative}{A_i})
	+
	\GTenToneHL{cost}{f_c}(\GTenToneHL{alternative}{A_i})
	+
	\GTenToneHL{cost}{f_c}(\GTenToneHL{state}{S_j})
	\right\}.
\end{aligned}
\]

\secInicio[sec:problem16-solar-installation-programme-c1-c3][sec:three-component-models]

\subsubsection*{C3 - Calculating profit and building the payoff table.}
\phantomsection
\label{subsec:problem16-solar-installation-programme-c3}

For every \GTenToneHL{alternative}{alternative}--state pair, calculate
\[
\begin{aligned}
	\GTenToneHL{revenue}{R}(\GTenToneHL{alternative}{A_i},\GTenToneHL{state}{S_j})
	&=
	\GTenToneHL{quantity}{q}(\GTenToneHL{alternative}{A_i},\GTenToneHL{state}{S_j})
	\left[
	\GTenToneHL{revenue}{r}(\GTenToneHL{alternative}{A_i})+\GTenToneHL{revenue}{r}(\GTenToneHL{state}{S_j})
	\right]
	+
	\GTenToneHL{revenue}{f_r}(\GTenToneHL{state}{S_j}),
	\\
	\GTenToneHL{cost}{C}(\GTenToneHL{alternative}{A_i},\GTenToneHL{state}{S_j})
	&=
	\GTenToneHL{quantity}{q}(\GTenToneHL{alternative}{A_i},\GTenToneHL{state}{S_j})\GTenToneHL{cost}{c}(\GTenToneHL{alternative}{A_i})
	+
	\GTenToneHL{cost}{f_c}(\GTenToneHL{alternative}{A_i})
	+
	\GTenToneHL{cost}{f_c}(\GTenToneHL{state}{S_j}),
	\\
	\GTenToneHL{payoff}{P}(\GTenToneHL{alternative}{A_i},\GTenToneHL{state}{S_j})
	&=
	\GTenToneHL{revenue}{R}(\GTenToneHL{alternative}{A_i},\GTenToneHL{state}{S_j})-\GTenToneHL{cost}{C}(\GTenToneHL{alternative}{A_i},\GTenToneHL{state}{S_j}).
\end{aligned}
\]

\begin{center}
	\small
	\begin{tabular}{llll}
		\toprule
		Cell
		& \(\GTenToneHL{revenue}{R}(\GTenToneHL{alternative}{A_i},\GTenToneHL{state}{S_j})\)
		& \(\GTenToneHL{cost}{C}(\GTenToneHL{alternative}{A_i},\GTenToneHL{state}{S_j})\)
		& \(\GTenToneHL{payoff}{P}(\GTenToneHL{alternative}{A_i},\GTenToneHL{state}{S_j})=\GTenToneHL{revenue}{R}(\GTenToneHL{alternative}{A_i},\GTenToneHL{state}{S_j})-\GTenToneHL{cost}{C}(\GTenToneHL{alternative}{A_i},\GTenToneHL{state}{S_j})\) \\
		\midrule
		\((\GTenToneHL{alternative}{A_1},\GTenToneHL{state}{S_1})\)
		& \(7(12+4)+20=132\)
		& \(7(6)+18+9=69\)
		& \(132-69=63\) \\
		
		\((\GTenToneHL{alternative}{A_1},\GTenToneHL{state}{S_2})\)
		& \(4(12+2)+11=67\)
		& \(4(6)+18+14=56\)
		& \(67-56=11\) \\
		
		\((\GTenToneHL{alternative}{A_2},\GTenToneHL{state}{S_1})\)
		& \(6(18+4)+20=152\)
		& \(6(8)+25+9=82\)
		& \(152-82=70\) \\
		
		\((\GTenToneHL{alternative}{A_2},\GTenToneHL{state}{S_2})\)
		& \(3(18+2)+11=71\)
		& \(3(8)+25+14=63\)
		& \(71-63=8\) \\
		\bottomrule
	\end{tabular}
\end{center}

The resulting \GTenToneHL{payoff}{payoff} table is:

\begin{center}
	\begin{tabular}{lrr}
		\toprule
		Alternative
		& Strong commercial demand
		& Limited residential demand \\
		\midrule
		Standard-panel programme & \GTenToneHL{payoff}{63} & \GTenToneHL{payoff}{11} \\
		High-efficiency-panel programme & \GTenToneHL{payoff}{70} & \GTenToneHL{payoff}{8} \\
		\bottomrule
	\end{tabular}
\end{center}

\secInicio[sec:problem16-solar-installation-programme-c1-c3][sec:three-component-models]

% ----------------------------------------------------------------
\subsection{Hybrid Conference Production (C1--C3)}
\label{sec:problem17-hybrid-conference-production-c1-c3}

An events company must choose a hybrid-conference production package before
attendance and sponsorship conditions are known. \GTenToneHL{revenue}{Revenue} consists of a
package-based rate per event segment, a fixed booking payment determined by
the selected package, and a fixed sponsorship payment determined by the
state. Costs consist of a state-based delivery rate, a fixed setup \GTenToneHL{cost}{cost}
determined by the package, and a fixed licensing \GTenToneHL{cost}{cost} determined by the
state.

\begin{multicols}{2}
	The company is comparing two \GTenToneHL{alternative}{alternatives}:
	\begin{enumerate}[label=\arabic*.]
		\item \(\GTenToneHL{alternative}{A_1}\): Use the standard hybrid package
		\item \(\GTenToneHL{alternative}{A_2}\): Use the immersive hybrid package
	\end{enumerate}
	
	The possible \GTenToneHL{state}{states} \GTenToneHL{state}{of} \GTenToneHL{state}{nature} are:
	\begin{enumerate}[label=\arabic*.]
		\item \(\GTenToneHL{state}{S_1}\): Strong attendance and sponsorship
		\item \(\GTenToneHL{state}{S_2}\): Limited attendance and sponsorship
	\end{enumerate}
	
	\columnbreak
	
	The quantities, measured in event segments, are
	\[
	\begin{aligned}
		\GTenToneHL{quantity}{q}(\GTenToneHL{alternative}{A_1},\GTenToneHL{state}{S_1})&=\GTenToneHL{quantity}{6},
		&
		\GTenToneHL{quantity}{q}(\GTenToneHL{alternative}{A_1},\GTenToneHL{state}{S_2})&=\GTenToneHL{quantity}{4},
		\\
		\GTenToneHL{quantity}{q}(\GTenToneHL{alternative}{A_2},\GTenToneHL{state}{S_1})&=\GTenToneHL{quantity}{8},
		&
		\GTenToneHL{quantity}{q}(\GTenToneHL{alternative}{A_2},\GTenToneHL{state}{S_2})&=5.
	\end{aligned}
	\]
	
	The \GTenToneHL{alternative}{alternative}-based \GTenToneHL{revenue}{revenue} rates, in thousand USD per event segment,
	are
	\[
	\GTenToneHL{revenue}{r}(\GTenToneHL{alternative}{A_1})=\GTenToneHL{revenue}{13},
	\qquad
	\GTenToneHL{revenue}{r}(\GTenToneHL{alternative}{A_2})=16.
	\]
	
	The \GTenToneHL{alternative}{alternative}-based fixed revenues, in thousand USD, are
	\[
	\GTenToneHL{revenue}{f_r}(\GTenToneHL{alternative}{A_1})=\GTenToneHL{revenue}{14},
	\qquad
	\GTenToneHL{revenue}{f_r}(\GTenToneHL{alternative}{A_2})=22.
	\]
	
	The state-based fixed revenues, in thousand USD, are
	\[
	\GTenToneHL{revenue}{f_r}(\GTenToneHL{state}{S_1})=\GTenToneHL{revenue}{18},
	\qquad
	\GTenToneHL{revenue}{f_r}(\GTenToneHL{state}{S_2})=9.
	\]
	
	The state-based \GTenToneHL{cost}{cost} rates, in thousand USD per event segment, are
	\[
	\GTenToneHL{cost}{c}(\GTenToneHL{state}{S_1})=\GTenToneHL{cost}{5},
	\qquad
	\GTenToneHL{cost}{c}(\GTenToneHL{state}{S_2})=7.
	\]
	
	The \GTenToneHL{alternative}{alternative}-based fixed \GTenToneHL{cost}{costs}, in thousand USD, are
	\[
	\GTenToneHL{cost}{f_c}(\GTenToneHL{alternative}{A_1})=\GTenToneHL{cost}{24},
	\qquad
	\GTenToneHL{cost}{f_c}(\GTenToneHL{alternative}{A_2})=37.
	\]
	
	The state-based fixed \GTenToneHL{cost}{costs}, in thousand USD, are
	\[
	\GTenToneHL{cost}{f_c}(\GTenToneHL{state}{S_1})=\GTenToneHL{cost}{8},
	\qquad
	\GTenToneHL{cost}{f_c}(\GTenToneHL{state}{S_2})=13.
	\]
\end{multicols}

\textbf{Questions/Tasks.}
\begin{enumerate}
	\CriterionTask
	{subsec:problem17-hybrid-conference-production-c1}{1}
	{Interpreting decision alternatives, events, consequences, and states.}

	\CriterionTask
	{subsec:problem17-hybrid-conference-production-c3}{3}
	{Calculating profit and building the payoff table.}
\end{enumerate}

\subsubsection*{C1 - Interpreting decision alternatives, events, consequences, and states.}
\phantomsection
\label{subsec:problem17-hybrid-conference-production-c1}

\DecisionElementsTable
{Use the standard hybrid package; Use the immersive hybrid package}
{Strong attendance and sponsorship; Limited attendance and sponsorship}
{Attendance and sponsorship conditions revealed after the package decision}
{Profit in thousand USD for each production-package--state pair}

\GTenToneHL{revenue}{Revenue} contains an \GTenToneHL{alternative}{alternative}-based variable component and
\GTenToneHL{alternative}{alternative}-based and state-based fixed components:
\[
\GTenToneHL{revenue}{R}(\GTenToneHL{alternative}{A_i},\GTenToneHL{state}{S_j})
=
\GTenToneHL{quantity}{q}(\GTenToneHL{alternative}{A_i},\GTenToneHL{state}{S_j})\GTenToneHL{revenue}{r}(\GTenToneHL{alternative}{A_i})
+
\GTenToneHL{revenue}{f_r}(\GTenToneHL{alternative}{A_i})
+
\GTenToneHL{revenue}{f_r}(\GTenToneHL{state}{S_j}).
\]

\GTenToneHL{cost}{Cost} contains a state-based variable component and \GTenToneHL{alternative}{alternative}-based and
state-based fixed components:
\[
\GTenToneHL{cost}{C}(\GTenToneHL{alternative}{A_i},\GTenToneHL{state}{S_j})
=
\GTenToneHL{quantity}{q}(\GTenToneHL{alternative}{A_i},\GTenToneHL{state}{S_j})\GTenToneHL{cost}{c}(\GTenToneHL{state}{S_j})
+
\GTenToneHL{cost}{f_c}(\GTenToneHL{alternative}{A_i})
+
\GTenToneHL{cost}{f_c}(\GTenToneHL{state}{S_j}).
\]

Therefore,
\[
\begin{aligned}
	\GTenToneHL{payoff}{P}(\GTenToneHL{alternative}{A_i},\GTenToneHL{state}{S_j})
	={}&
	\left\{
	\GTenToneHL{quantity}{q}(\GTenToneHL{alternative}{A_i},\GTenToneHL{state}{S_j})\GTenToneHL{revenue}{r}(\GTenToneHL{alternative}{A_i})
	+
	\GTenToneHL{revenue}{f_r}(\GTenToneHL{alternative}{A_i})
	+
	\GTenToneHL{revenue}{f_r}(\GTenToneHL{state}{S_j})
	\right\}
	\\
	&-
	\left\{
	\GTenToneHL{quantity}{q}(\GTenToneHL{alternative}{A_i},\GTenToneHL{state}{S_j})\GTenToneHL{cost}{c}(\GTenToneHL{state}{S_j})
	+
	\GTenToneHL{cost}{f_c}(\GTenToneHL{alternative}{A_i})
	+
	\GTenToneHL{cost}{f_c}(\GTenToneHL{state}{S_j})
	\right\}.
\end{aligned}
\]

\secInicio[sec:problem17-hybrid-conference-production-c1-c3][sec:three-component-models]

\subsubsection*{C3 - Calculating profit and building the payoff table.}
\phantomsection
\label{subsec:problem17-hybrid-conference-production-c3}

For every \GTenToneHL{alternative}{alternative}--state pair, calculate
\[
\begin{aligned}
	\GTenToneHL{revenue}{R}(\GTenToneHL{alternative}{A_i},\GTenToneHL{state}{S_j})
	&=
	\GTenToneHL{quantity}{q}(\GTenToneHL{alternative}{A_i},\GTenToneHL{state}{S_j})\GTenToneHL{revenue}{r}(\GTenToneHL{alternative}{A_i})
	+
	\GTenToneHL{revenue}{f_r}(\GTenToneHL{alternative}{A_i})
	+
	\GTenToneHL{revenue}{f_r}(\GTenToneHL{state}{S_j}),
	\\
	\GTenToneHL{cost}{C}(\GTenToneHL{alternative}{A_i},\GTenToneHL{state}{S_j})
	&=
	\GTenToneHL{quantity}{q}(\GTenToneHL{alternative}{A_i},\GTenToneHL{state}{S_j})\GTenToneHL{cost}{c}(\GTenToneHL{state}{S_j})
	+
	\GTenToneHL{cost}{f_c}(\GTenToneHL{alternative}{A_i})
	+
	\GTenToneHL{cost}{f_c}(\GTenToneHL{state}{S_j}),
	\\
	\GTenToneHL{payoff}{P}(\GTenToneHL{alternative}{A_i},\GTenToneHL{state}{S_j})
	&=
	\GTenToneHL{revenue}{R}(\GTenToneHL{alternative}{A_i},\GTenToneHL{state}{S_j})-\GTenToneHL{cost}{C}(\GTenToneHL{alternative}{A_i},\GTenToneHL{state}{S_j}).
\end{aligned}
\]

\begin{center}
	\small
	\begin{tabular}{llll}
		\toprule
		Cell
		& \(\GTenToneHL{revenue}{R}(\GTenToneHL{alternative}{A_i},\GTenToneHL{state}{S_j})\)
		& \(\GTenToneHL{cost}{C}(\GTenToneHL{alternative}{A_i},\GTenToneHL{state}{S_j})\)
		& \(\GTenToneHL{payoff}{P}(\GTenToneHL{alternative}{A_i},\GTenToneHL{state}{S_j})=\GTenToneHL{revenue}{R}(\GTenToneHL{alternative}{A_i},\GTenToneHL{state}{S_j})-\GTenToneHL{cost}{C}(\GTenToneHL{alternative}{A_i},\GTenToneHL{state}{S_j})\) \\
		\midrule
		\((\GTenToneHL{alternative}{A_1},\GTenToneHL{state}{S_1})\)
		& \(6(13)+14+18=110\)
		& \(6(5)+24+8=62\)
		& \(110-62=48\) \\
		
		\((\GTenToneHL{alternative}{A_1},\GTenToneHL{state}{S_2})\)
		& \(4(13)+14+9=75\)
		& \(4(7)+24+13=65\)
		& \(75-65=10\) \\
		
		\((\GTenToneHL{alternative}{A_2},\GTenToneHL{state}{S_1})\)
		& \(8(16)+22+18=168\)
		& \(8(5)+37+8=85\)
		& \(168-85=83\) \\
		
		\((\GTenToneHL{alternative}{A_2},\GTenToneHL{state}{S_2})\)
		& \(5(16)+22+9=111\)
		& \(5(7)+37+13=85\)
		& \(111-85=26\) \\
		\bottomrule
	\end{tabular}
\end{center}

The resulting \GTenToneHL{payoff}{payoff} table is:

\begin{center}
	\begin{tabular}{lrr}
		\toprule
		Alternative
		& Strong attendance and sponsorship
		& Limited attendance and sponsorship \\
		\midrule
		Standard hybrid package & \GTenToneHL{payoff}{48} & \GTenToneHL{payoff}{10} \\
		Immersive hybrid package & \GTenToneHL{payoff}{83} & \GTenToneHL{payoff}{26} \\
		\bottomrule
	\end{tabular}
\end{center}

\secInicio[sec:problem17-hybrid-conference-production-c1-c3][sec:three-component-models]

% ----------------------------------------------------------------
\subsection{Cold-Storage Export Contract (C1--C3)}
\label{sec:problem18-cold-storage-export-contract-c1-c3}

A food-export company must select a cold-storage configuration before export
demand and energy-market conditions are known. \GTenToneHL{revenue}{Revenue} consists of a
state-based rate per storage batch, a fixed capacity payment determined by
the selected configuration, and a fixed export payment determined by the
state. Variable \GTenToneHL{cost}{cost} per batch combines a configuration-based operating rate
and a state-based energy surcharge. The selected configuration also produces
a fixed setup \GTenToneHL{cost}{cost}.

\begin{multicols}{2}
	The company is comparing two \GTenToneHL{alternative}{alternatives}:
	\begin{enumerate}[label=\arabic*.]
		\item \(\GTenToneHL{alternative}{A_1}\): Use the modular cold-storage facility
		\item \(\GTenToneHL{alternative}{A_2}\): Use the large cold-storage facility
	\end{enumerate}
	
	The possible \GTenToneHL{state}{states} \GTenToneHL{state}{of} \GTenToneHL{state}{nature} are:
	\begin{enumerate}[label=\arabic*.]
		\item \(\GTenToneHL{state}{S_1}\): Strong export demand
		\item \(\GTenToneHL{state}{S_2}\): Limited export demand
	\end{enumerate}
	
	\columnbreak
	
	The quantities, measured in storage batches, are
	\[
	\begin{aligned}
		\GTenToneHL{quantity}{q}(\GTenToneHL{alternative}{A_1},\GTenToneHL{state}{S_1})&=\GTenToneHL{quantity}{9},
		&
		\GTenToneHL{quantity}{q}(\GTenToneHL{alternative}{A_1},\GTenToneHL{state}{S_2})&=\GTenToneHL{quantity}{5},
		\\
		\GTenToneHL{quantity}{q}(\GTenToneHL{alternative}{A_2},\GTenToneHL{state}{S_1})&=\GTenToneHL{quantity}{12},
		&
		\GTenToneHL{quantity}{q}(\GTenToneHL{alternative}{A_2},\GTenToneHL{state}{S_2})&=7.
	\end{aligned}
	\]
	
	The state-based \GTenToneHL{revenue}{revenue} rates, in thousand USD per storage batch, are
	\[
	\GTenToneHL{revenue}{r}(\GTenToneHL{state}{S_1})=\GTenToneHL{revenue}{10},
	\qquad
	\GTenToneHL{revenue}{r}(\GTenToneHL{state}{S_2})=7.
	\]
	
	The \GTenToneHL{alternative}{alternative}-based fixed revenues, in thousand USD, are
	\[
	\GTenToneHL{revenue}{f_r}(\GTenToneHL{alternative}{A_1})=\GTenToneHL{revenue}{20},
	\qquad
	\GTenToneHL{revenue}{f_r}(\GTenToneHL{alternative}{A_2})=30.
	\]
	
	The state-based fixed revenues, in thousand USD, are
	\[
	\GTenToneHL{revenue}{f_r}(\GTenToneHL{state}{S_1})=\GTenToneHL{revenue}{16},
	\qquad
	\GTenToneHL{revenue}{f_r}(\GTenToneHL{state}{S_2})=8.
	\]
	
	The \GTenToneHL{alternative}{alternative}-based \GTenToneHL{cost}{cost} rates, in thousand USD per storage batch, are
	\[
	\GTenToneHL{cost}{c}(\GTenToneHL{alternative}{A_1})=\GTenToneHL{cost}{4},
	\qquad
	\GTenToneHL{cost}{c}(\GTenToneHL{alternative}{A_2})=3.
	\]
	
	The state-based \GTenToneHL{cost}{cost} rates are
	\[
	\GTenToneHL{cost}{c}(\GTenToneHL{state}{S_1})=\GTenToneHL{cost}{2},
	\qquad
	\GTenToneHL{cost}{c}(\GTenToneHL{state}{S_2})=4.
	\]
	
	The \GTenToneHL{alternative}{alternative}-based fixed \GTenToneHL{cost}{costs}, in thousand USD, are
	\[
	\GTenToneHL{cost}{f_c}(\GTenToneHL{alternative}{A_1})=\GTenToneHL{cost}{28},
	\qquad
	\GTenToneHL{cost}{f_c}(\GTenToneHL{alternative}{A_2})=42.
	\]
\end{multicols}

\textbf{Questions/Tasks.}
\begin{enumerate}
	\CriterionTask
	{subsec:problem18-cold-storage-export-contract-c1}{1}
	{Interpreting decision alternatives, events, consequences, and states.}

	\CriterionTask
	{subsec:problem18-cold-storage-export-contract-c3}{3}
	{Calculating profit and building the payoff table.}
\end{enumerate}

\subsubsection*{C1 - Interpreting decision alternatives, events, consequences, and states.}
\phantomsection
\label{subsec:problem18-cold-storage-export-contract-c1}

\DecisionElementsTable
{Use the modular cold-storage facility; Use the large cold-storage facility}
{Strong export demand; Limited export demand}
{Export demand and energy-market conditions revealed after the facility decision}
{Profit in thousand USD for each storage-facility--state pair}

\GTenToneHL{revenue}{Revenue} contains a state-based variable component and \GTenToneHL{alternative}{alternative}-based and
state-based fixed components:
\[
\GTenToneHL{revenue}{R}(\GTenToneHL{alternative}{A_i},\GTenToneHL{state}{S_j})
=
\GTenToneHL{quantity}{q}(\GTenToneHL{alternative}{A_i},\GTenToneHL{state}{S_j})\GTenToneHL{revenue}{r}(\GTenToneHL{state}{S_j})
+
\GTenToneHL{revenue}{f_r}(\GTenToneHL{alternative}{A_i})
+
\GTenToneHL{revenue}{f_r}(\GTenToneHL{state}{S_j}).
\]

\GTenToneHL{cost}{Cost} contains \GTenToneHL{alternative}{alternative}-based and state-based variable components and an
\GTenToneHL{alternative}{alternative}-based fixed component:
\[
\GTenToneHL{cost}{C}(\GTenToneHL{alternative}{A_i},\GTenToneHL{state}{S_j})
=
\GTenToneHL{quantity}{q}(\GTenToneHL{alternative}{A_i},\GTenToneHL{state}{S_j})
\left[
\GTenToneHL{cost}{c}(\GTenToneHL{alternative}{A_i})+\GTenToneHL{cost}{c}(\GTenToneHL{state}{S_j})
\right]
+
\GTenToneHL{cost}{f_c}(\GTenToneHL{alternative}{A_i}).
\]

Therefore,
\[
\begin{aligned}
	\GTenToneHL{payoff}{P}(\GTenToneHL{alternative}{A_i},\GTenToneHL{state}{S_j})
	={}&
	\left\{
	\GTenToneHL{quantity}{q}(\GTenToneHL{alternative}{A_i},\GTenToneHL{state}{S_j})\GTenToneHL{revenue}{r}(\GTenToneHL{state}{S_j})
	+
	\GTenToneHL{revenue}{f_r}(\GTenToneHL{alternative}{A_i})
	+
	\GTenToneHL{revenue}{f_r}(\GTenToneHL{state}{S_j})
	\right\}
	\\
	&-
	\left\{
	\GTenToneHL{quantity}{q}(\GTenToneHL{alternative}{A_i},\GTenToneHL{state}{S_j})
	\left[
	\GTenToneHL{cost}{c}(\GTenToneHL{alternative}{A_i})+\GTenToneHL{cost}{c}(\GTenToneHL{state}{S_j})
	\right]
	+
	\GTenToneHL{cost}{f_c}(\GTenToneHL{alternative}{A_i})
	\right\}.
\end{aligned}
\]

\secInicio[sec:problem18-cold-storage-export-contract-c1-c3][sec:three-component-models]

\subsubsection*{C3 - Calculating profit and building the payoff table.}
\phantomsection
\label{subsec:problem18-cold-storage-export-contract-c3}

For every \GTenToneHL{alternative}{alternative}--state pair, calculate
\[
\begin{aligned}
	\GTenToneHL{revenue}{R}(\GTenToneHL{alternative}{A_i},\GTenToneHL{state}{S_j})
	&=
	\GTenToneHL{quantity}{q}(\GTenToneHL{alternative}{A_i},\GTenToneHL{state}{S_j})\GTenToneHL{revenue}{r}(\GTenToneHL{state}{S_j})
	+
	\GTenToneHL{revenue}{f_r}(\GTenToneHL{alternative}{A_i})
	+
	\GTenToneHL{revenue}{f_r}(\GTenToneHL{state}{S_j}),
	\\
	\GTenToneHL{cost}{C}(\GTenToneHL{alternative}{A_i},\GTenToneHL{state}{S_j})
	&=
	\GTenToneHL{quantity}{q}(\GTenToneHL{alternative}{A_i},\GTenToneHL{state}{S_j})
	\left[
	\GTenToneHL{cost}{c}(\GTenToneHL{alternative}{A_i})+\GTenToneHL{cost}{c}(\GTenToneHL{state}{S_j})
	\right]
	+
	\GTenToneHL{cost}{f_c}(\GTenToneHL{alternative}{A_i}),
	\\
	\GTenToneHL{payoff}{P}(\GTenToneHL{alternative}{A_i},\GTenToneHL{state}{S_j})
	&=
	\GTenToneHL{revenue}{R}(\GTenToneHL{alternative}{A_i},\GTenToneHL{state}{S_j})-\GTenToneHL{cost}{C}(\GTenToneHL{alternative}{A_i},\GTenToneHL{state}{S_j}).
\end{aligned}
\]

\begin{center}
	\small
	\begin{tabular}{llll}
		\toprule
		Cell
		& \(\GTenToneHL{revenue}{R}(\GTenToneHL{alternative}{A_i},\GTenToneHL{state}{S_j})\)
		& \(\GTenToneHL{cost}{C}(\GTenToneHL{alternative}{A_i},\GTenToneHL{state}{S_j})\)
		& \(\GTenToneHL{payoff}{P}(\GTenToneHL{alternative}{A_i},\GTenToneHL{state}{S_j})=\GTenToneHL{revenue}{R}(\GTenToneHL{alternative}{A_i},\GTenToneHL{state}{S_j})-\GTenToneHL{cost}{C}(\GTenToneHL{alternative}{A_i},\GTenToneHL{state}{S_j})\) \\
		\midrule
		\((\GTenToneHL{alternative}{A_1},\GTenToneHL{state}{S_1})\)
		& \(9(10)+20+16=126\)
		& \(9(4+2)+28=82\)
		& \(126-82=44\) \\
		
		\((\GTenToneHL{alternative}{A_1},\GTenToneHL{state}{S_2})\)
		& \(5(7)+20+8=63\)
		& \(5(4+4)+28=68\)
		& \(63-68=-5\) \\
		
		\((\GTenToneHL{alternative}{A_2},\GTenToneHL{state}{S_1})\)
		& \(12(10)+30+16=166\)
		& \(12(3+2)+42=102\)
		& \(166-102=64\) \\
		
		\((\GTenToneHL{alternative}{A_2},\GTenToneHL{state}{S_2})\)
		& \(7(7)+30+8=87\)
		& \(7(3+4)+42=91\)
		& \(87-91=-4\) \\
		\bottomrule
	\end{tabular}
\end{center}

The resulting \GTenToneHL{payoff}{payoff} table is:

\begin{center}
	\begin{tabular}{lrr}
		\toprule
		Alternative & Strong export demand & Limited export demand \\
		\midrule
		Modular cold-storage facility & \GTenToneHL{payoff}{44} & \(-5\) \\
		Large cold-storage facility & \GTenToneHL{payoff}{64} & \(-4\) \\
		\bottomrule
	\end{tabular}
\end{center}

The negative \GTenToneHL{payoff}{payoffs} under limited export demand represent losses because
the \GTenToneHL{revenue}{revenue} earned does not cover the variable and fixed \GTenToneHL{cost}{costs}.

\secInicio[sec:problem18-cold-storage-export-contract-c1-c3][sec:three-component-models]

% ================================================================
\section{Complete Four-Component Models}
\label{sec:complete-models}

\SectionProblemsTableOfContents{
	\SectionProblemLink{sec:problem19-smart-manufacturing-contract-c1-c3}{Smart Manufacturing Contract (C1--C3)} \\
	\SectionProblemLink{sec:problem20-distribution-centre-automation-c1-c3}{Distribution Centre Automation (C1--C3)} \\
	\SectionProblemLink{sec:problem21-electric-vehicle-charging-network-c1-c3}{Electric Vehicle Charging Network (C1--C3)}
}

\vspace{0.8cm}

In a complete four-component model, both \GTenToneHL{revenue}{revenue} and \GTenToneHL{cost}{cost} contain all four
possible components:

\begin{enumerate}[label=\arabic*.]
	\item an \GTenToneHL{alternative}{alternative}-based variable component;
	\item a state-based variable component;
	\item an \GTenToneHL{alternative}{alternative}-based fixed component;
	\item a state-based fixed component.
\end{enumerate}

Therefore, the complete \GTenToneHL{revenue}{revenue} and \GTenToneHL{cost}{cost} models are
\[
\begin{aligned}
	\GTenToneHL{revenue}{R}(\GTenToneHL{alternative}{A_i},\GTenToneHL{state}{S_j})
	&=
	\GTenToneHL{quantity}{q}(\GTenToneHL{alternative}{A_i},\GTenToneHL{state}{S_j})
	\left[
	\GTenToneHL{revenue}{r}(\GTenToneHL{alternative}{A_i})+\GTenToneHL{revenue}{r}(\GTenToneHL{state}{S_j})
	\right]
	+
	\GTenToneHL{revenue}{f_r}(\GTenToneHL{alternative}{A_i})
	+
	\GTenToneHL{revenue}{f_r}(\GTenToneHL{state}{S_j}),
	\\
	\GTenToneHL{cost}{C}(\GTenToneHL{alternative}{A_i},\GTenToneHL{state}{S_j})
	&=
	\GTenToneHL{quantity}{q}(\GTenToneHL{alternative}{A_i},\GTenToneHL{state}{S_j})
	\left[
	\GTenToneHL{cost}{c}(\GTenToneHL{alternative}{A_i})+\GTenToneHL{cost}{c}(\GTenToneHL{state}{S_j})
	\right]
	+
	\GTenToneHL{cost}{f_c}(\GTenToneHL{alternative}{A_i})
	+
	\GTenToneHL{cost}{f_c}(\GTenToneHL{state}{S_j}).
\end{aligned}
\]

The \GTenToneHL{payoff}{payoff} for each \GTenToneHL{alternative}{alternative}--state pair is consequently
\[
\begin{aligned}
	\GTenToneHL{payoff}{P}(\GTenToneHL{alternative}{A_i},\GTenToneHL{state}{S_j})
	={}&
	\left\{
	\GTenToneHL{quantity}{q}(\GTenToneHL{alternative}{A_i},\GTenToneHL{state}{S_j})
	\left[
	\GTenToneHL{revenue}{r}(\GTenToneHL{alternative}{A_i})+\GTenToneHL{revenue}{r}(\GTenToneHL{state}{S_j})
	\right]
	+
	\GTenToneHL{revenue}{f_r}(\GTenToneHL{alternative}{A_i})
	+
	\GTenToneHL{revenue}{f_r}(\GTenToneHL{state}{S_j})
	\right\}
	\\
	&-
	\left\{
	\GTenToneHL{quantity}{q}(\GTenToneHL{alternative}{A_i},\GTenToneHL{state}{S_j})
	\left[
	\GTenToneHL{cost}{c}(\GTenToneHL{alternative}{A_i})+\GTenToneHL{cost}{c}(\GTenToneHL{state}{S_j})
	\right]
	+
	\GTenToneHL{cost}{f_c}(\GTenToneHL{alternative}{A_i})
	+
	\GTenToneHL{cost}{f_c}(\GTenToneHL{state}{S_j})
	\right\}.
\end{aligned}
\]

The three problems in this section use \(2\times2\), \(3\times2\), and
\(2\times3\) \GTenToneHL{alternative}{alternative}--state structures, respectively.

% ----------------------------------------------------------------
\subsection{Smart Manufacturing Contract (C1--C3)}
\label{sec:problem19-smart-manufacturing-contract-c1-c3}

A manufacturing company must select a production system before product demand
and input-market conditions are known. \GTenToneHL{revenue}{Revenue} per production batch combines a
rate determined by the selected system and a rate determined by the state.
\GTenToneHL{revenue}{Revenue} also includes a fixed contract payment determined by the system and a
fixed market incentive determined by the state.

Variable \GTenToneHL{cost}{cost} per batch combines a system-based operating rate and a
state-based material surcharge. Fixed \GTenToneHL{cost}{cost} consists of a setup \GTenToneHL{cost}{cost} determined
by the selected system and a compliance \GTenToneHL{cost}{cost} determined by the state.

\begin{multicols}{2}
	The company is comparing two \GTenToneHL{alternative}{alternatives}:
	\begin{enumerate}[label=\arabic*.]
		\item \(\GTenToneHL{alternative}{A_1}\): Use the conventional assembly system
		\item \(\GTenToneHL{alternative}{A_2}\): Use the automated assembly system
	\end{enumerate}
	
	The possible \GTenToneHL{state}{states} \GTenToneHL{state}{of} \GTenToneHL{state}{nature} are:
	\begin{enumerate}[label=\arabic*.]
		\item \(\GTenToneHL{state}{S_1}\): Strong demand with favourable input conditions
		\item \(\GTenToneHL{state}{S_2}\): Limited demand with expensive inputs
	\end{enumerate}
	
	The quantities, measured in production batches, are
	\[
	\begin{aligned}
		\GTenToneHL{quantity}{q}(\GTenToneHL{alternative}{A_1},\GTenToneHL{state}{S_1})&=\GTenToneHL{quantity}{8},
		&
		\GTenToneHL{quantity}{q}(\GTenToneHL{alternative}{A_1},\GTenToneHL{state}{S_2})&=\GTenToneHL{quantity}{5},
		\\
		\GTenToneHL{quantity}{q}(\GTenToneHL{alternative}{A_2},\GTenToneHL{state}{S_1})&=\GTenToneHL{quantity}{11},
		&
		\GTenToneHL{quantity}{q}(\GTenToneHL{alternative}{A_2},\GTenToneHL{state}{S_2})&=6.
	\end{aligned}
	\]
	
	The \GTenToneHL{alternative}{alternative}-based \GTenToneHL{revenue}{revenue} rates, in thousand USD per production
	batch, are
	\[
	\GTenToneHL{revenue}{r}(\GTenToneHL{alternative}{A_1})=\GTenToneHL{revenue}{12},
	\qquad
	\GTenToneHL{revenue}{r}(\GTenToneHL{alternative}{A_2})=15.
	\]
	
	The state-based \GTenToneHL{revenue}{revenue} rates, in thousand USD per production batch, are
	\[
	\GTenToneHL{revenue}{r}(\GTenToneHL{state}{S_1})=\GTenToneHL{revenue}{4},
	\qquad
	\GTenToneHL{revenue}{r}(\GTenToneHL{state}{S_2})=2.
	\]
	
	\columnbreak
	
	The \GTenToneHL{alternative}{alternative}-based fixed revenues, in thousand USD, are
	\[
	\GTenToneHL{revenue}{f_r}(\GTenToneHL{alternative}{A_1})=\GTenToneHL{revenue}{18},
	\qquad
	\GTenToneHL{revenue}{f_r}(\GTenToneHL{alternative}{A_2})=28.
	\]
	
	The state-based fixed revenues, in thousand USD, are
	\[
	\GTenToneHL{revenue}{f_r}(\GTenToneHL{state}{S_1})=\GTenToneHL{revenue}{14},
	\qquad
	\GTenToneHL{revenue}{f_r}(\GTenToneHL{state}{S_2})=7.
	\]
	
	The \GTenToneHL{alternative}{alternative}-based \GTenToneHL{cost}{cost} rates, in thousand USD per production
	batch, are
	\[
	\GTenToneHL{cost}{c}(\GTenToneHL{alternative}{A_1})=\GTenToneHL{cost}{5},
	\qquad
	\GTenToneHL{cost}{c}(\GTenToneHL{alternative}{A_2})=6.
	\]
	
	The state-based \GTenToneHL{cost}{cost} rates, in thousand USD per production batch, are
	\[
	\GTenToneHL{cost}{c}(\GTenToneHL{state}{S_1})=\GTenToneHL{cost}{2},
	\qquad
	\GTenToneHL{cost}{c}(\GTenToneHL{state}{S_2})=4.
	\]
	
	The \GTenToneHL{alternative}{alternative}-based fixed \GTenToneHL{cost}{costs}, in thousand USD, are
	\[
	\GTenToneHL{cost}{f_c}(\GTenToneHL{alternative}{A_1})=\GTenToneHL{cost}{24},
	\qquad
	\GTenToneHL{cost}{f_c}(\GTenToneHL{alternative}{A_2})=38.
	\]
	
	The state-based fixed \GTenToneHL{cost}{costs}, in thousand USD, are
	\[
	\GTenToneHL{cost}{f_c}(\GTenToneHL{state}{S_1})=\GTenToneHL{cost}{9},
	\qquad
	\GTenToneHL{cost}{f_c}(\GTenToneHL{state}{S_2})=15.
	\]
\end{multicols}

\textbf{Questions/Tasks.}
\begin{enumerate}
	\CriterionTask
	{subsec:problem19-smart-manufacturing-contract-c1}{1}
	{Interpreting decision alternatives, events, consequences, and states.}

	\CriterionTask
	{subsec:problem19-smart-manufacturing-contract-c3}{3}
	{Calculating profit and building the payoff table.}
\end{enumerate}

\subsubsection*{C1 - Interpreting decision alternatives, events, consequences, and states.}
\phantomsection
\label{subsec:problem19-smart-manufacturing-contract-c1}

\DecisionElementsTable
{Use the conventional assembly system; Use the automated assembly system}
{Strong demand with favourable input conditions; Limited demand with expensive inputs}
{Demand and input-market conditions revealed after the production-system decision}
{Profit in thousand USD for each production-system--state pair}

\GTenToneHL{revenue}{Revenue} contains \GTenToneHL{alternative}{alternative}-based and state-based variable components and
\GTenToneHL{alternative}{alternative}-based and state-based fixed components:
\[
\GTenToneHL{revenue}{R}(\GTenToneHL{alternative}{A_i},\GTenToneHL{state}{S_j})
=
\GTenToneHL{quantity}{q}(\GTenToneHL{alternative}{A_i},\GTenToneHL{state}{S_j})
\left[
\GTenToneHL{revenue}{r}(\GTenToneHL{alternative}{A_i})+\GTenToneHL{revenue}{r}(\GTenToneHL{state}{S_j})
\right]
+
\GTenToneHL{revenue}{f_r}(\GTenToneHL{alternative}{A_i})
+
\GTenToneHL{revenue}{f_r}(\GTenToneHL{state}{S_j}).
\]

\GTenToneHL{cost}{Cost} also contains \GTenToneHL{alternative}{alternative}-based and state-based variable components and
\GTenToneHL{alternative}{alternative}-based and state-based fixed components:
\[
\GTenToneHL{cost}{C}(\GTenToneHL{alternative}{A_i},\GTenToneHL{state}{S_j})
=
\GTenToneHL{quantity}{q}(\GTenToneHL{alternative}{A_i},\GTenToneHL{state}{S_j})
\left[
\GTenToneHL{cost}{c}(\GTenToneHL{alternative}{A_i})+\GTenToneHL{cost}{c}(\GTenToneHL{state}{S_j})
\right]
+
\GTenToneHL{cost}{f_c}(\GTenToneHL{alternative}{A_i})
+
\GTenToneHL{cost}{f_c}(\GTenToneHL{state}{S_j}).
\]

Therefore,
\[
\begin{aligned}
	\GTenToneHL{payoff}{P}(\GTenToneHL{alternative}{A_i},\GTenToneHL{state}{S_j})
	={}&
	\left\{
	\GTenToneHL{quantity}{q}(\GTenToneHL{alternative}{A_i},\GTenToneHL{state}{S_j})
	\left[
	\GTenToneHL{revenue}{r}(\GTenToneHL{alternative}{A_i})+\GTenToneHL{revenue}{r}(\GTenToneHL{state}{S_j})
	\right]
	+
	\GTenToneHL{revenue}{f_r}(\GTenToneHL{alternative}{A_i})
	+
	\GTenToneHL{revenue}{f_r}(\GTenToneHL{state}{S_j})
	\right\}
	\\
	&-
	\left\{
	\GTenToneHL{quantity}{q}(\GTenToneHL{alternative}{A_i},\GTenToneHL{state}{S_j})
	\left[
	\GTenToneHL{cost}{c}(\GTenToneHL{alternative}{A_i})+\GTenToneHL{cost}{c}(\GTenToneHL{state}{S_j})
	\right]
	+
	\GTenToneHL{cost}{f_c}(\GTenToneHL{alternative}{A_i})
	+
	\GTenToneHL{cost}{f_c}(\GTenToneHL{state}{S_j})
	\right\}.
\end{aligned}
\]

\secInicio[sec:problem19-smart-manufacturing-contract-c1-c3][sec:complete-models]

\subsubsection*{C3 - Calculating profit and building the payoff table.}
\phantomsection
\label{subsec:problem19-smart-manufacturing-contract-c3}

For every \GTenToneHL{alternative}{alternative}--state pair, calculate
\[
\begin{aligned}
	\GTenToneHL{revenue}{R}(\GTenToneHL{alternative}{A_i},\GTenToneHL{state}{S_j})
	&=
	\GTenToneHL{quantity}{q}(\GTenToneHL{alternative}{A_i},\GTenToneHL{state}{S_j})
	\left[
	\GTenToneHL{revenue}{r}(\GTenToneHL{alternative}{A_i})+\GTenToneHL{revenue}{r}(\GTenToneHL{state}{S_j})
	\right]
	+
	\GTenToneHL{revenue}{f_r}(\GTenToneHL{alternative}{A_i})
	+
	\GTenToneHL{revenue}{f_r}(\GTenToneHL{state}{S_j}),
	\\
	\GTenToneHL{cost}{C}(\GTenToneHL{alternative}{A_i},\GTenToneHL{state}{S_j})
	&=
	\GTenToneHL{quantity}{q}(\GTenToneHL{alternative}{A_i},\GTenToneHL{state}{S_j})
	\left[
	\GTenToneHL{cost}{c}(\GTenToneHL{alternative}{A_i})+\GTenToneHL{cost}{c}(\GTenToneHL{state}{S_j})
	\right]
	+
	\GTenToneHL{cost}{f_c}(\GTenToneHL{alternative}{A_i})
	+
	\GTenToneHL{cost}{f_c}(\GTenToneHL{state}{S_j}),
	\\
	\GTenToneHL{payoff}{P}(\GTenToneHL{alternative}{A_i},\GTenToneHL{state}{S_j})
	&=
	\GTenToneHL{revenue}{R}(\GTenToneHL{alternative}{A_i},\GTenToneHL{state}{S_j})-\GTenToneHL{cost}{C}(\GTenToneHL{alternative}{A_i},\GTenToneHL{state}{S_j}).
\end{aligned}
\]

\begin{center}
	\small
	\begin{tabular}{llll}
		\toprule
		Cell
		& \(\GTenToneHL{revenue}{R}(\GTenToneHL{alternative}{A_i},\GTenToneHL{state}{S_j})\)
		& \(\GTenToneHL{cost}{C}(\GTenToneHL{alternative}{A_i},\GTenToneHL{state}{S_j})\)
		& \(\GTenToneHL{payoff}{P}(\GTenToneHL{alternative}{A_i},\GTenToneHL{state}{S_j})=\GTenToneHL{revenue}{R}(\GTenToneHL{alternative}{A_i},\GTenToneHL{state}{S_j})-\GTenToneHL{cost}{C}(\GTenToneHL{alternative}{A_i},\GTenToneHL{state}{S_j})\) \\
		\midrule
		\((\GTenToneHL{alternative}{A_1},\GTenToneHL{state}{S_1})\)
		& \(8(12+4)+18+14=160\)
		& \(8(5+2)+24+9=89\)
		& \(160-89=71\) \\
		
		\((\GTenToneHL{alternative}{A_1},\GTenToneHL{state}{S_2})\)
		& \(5(12+2)+18+7=95\)
		& \(5(5+4)+24+15=84\)
		& \(95-84=11\) \\
		
		\((\GTenToneHL{alternative}{A_2},\GTenToneHL{state}{S_1})\)
		& \(11(15+4)+28+14=251\)
		& \(11(6+2)+38+9=135\)
		& \(251-135=116\) \\
		
		\((\GTenToneHL{alternative}{A_2},\GTenToneHL{state}{S_2})\)
		& \(6(15+2)+28+7=137\)
		& \(6(6+4)+38+15=113\)
		& \(137-113=24\) \\
		\bottomrule
	\end{tabular}
\end{center}

The resulting \GTenToneHL{payoff}{payoff} table is:

\begin{center}
	\begin{tabular}{lrr}
		\toprule
		Alternative
		& Strong demand, favourable inputs
		& Limited demand, expensive inputs \\
		\midrule
		Conventional assembly system & \GTenToneHL{payoff}{71} & \GTenToneHL{payoff}{11} \\
		Automated assembly system & \GTenToneHL{payoff}{116} & \GTenToneHL{payoff}{24} \\
		\bottomrule
	\end{tabular}
\end{center}

\secInicio[sec:problem19-smart-manufacturing-contract-c1-c3][sec:complete-models]

% ----------------------------------------------------------------
\subsection{Distribution Centre Automation (C1--C3)}
\label{sec:problem20-distribution-centre-automation-c1-c3}

A retail company must select a distribution-centre configuration before order
volume and energy-market conditions are known. \GTenToneHL{revenue}{Revenue} per order batch combines
a rate determined by the selected configuration and a rate determined by the
state. \GTenToneHL{revenue}{Revenue} also includes a fixed service payment determined by the
configuration and a fixed delivery incentive determined by the state.

Variable \GTenToneHL{cost}{cost} per order batch combines a configuration-based operating rate and
a state-based energy surcharge. Fixed \GTenToneHL{cost}{cost} consists of an installation \GTenToneHL{cost}{cost}
determined by the configuration and a compliance \GTenToneHL{cost}{cost} determined by the state.

\begin{multicols}{2}
	The company is comparing three \GTenToneHL{alternative}{alternatives}:
	\begin{enumerate}[label=\arabic*.]
		\item \(\GTenToneHL{alternative}{A_1}\): Use the manual distribution configuration
		\item \(\GTenToneHL{alternative}{A_2}\): Use the semi-automated distribution configuration
		\item \(\GTenToneHL{alternative}{A_3}\): Use the robotic distribution configuration
	\end{enumerate}
	
	The possible \GTenToneHL{state}{states} \GTenToneHL{state}{of} \GTenToneHL{state}{nature} are:
	\begin{enumerate}[label=\arabic*.]
		\item \(\GTenToneHL{state}{S_1}\): High order volume with favourable energy prices
		\item \(\GTenToneHL{state}{S_2}\): Moderate order volume with expensive energy
	\end{enumerate}
	
	The quantities, measured in hundreds of order batches, are
	\[
	\begin{aligned}
		\GTenToneHL{quantity}{q}(\GTenToneHL{alternative}{A_1},\GTenToneHL{state}{S_1})&=\GTenToneHL{quantity}{8},
		&
		\GTenToneHL{quantity}{q}(\GTenToneHL{alternative}{A_1},\GTenToneHL{state}{S_2})&=\GTenToneHL{quantity}{5},
		\\
		\GTenToneHL{quantity}{q}(\GTenToneHL{alternative}{A_2},\GTenToneHL{state}{S_1})&=\GTenToneHL{quantity}{11},
		&
		\GTenToneHL{quantity}{q}(\GTenToneHL{alternative}{A_2},\GTenToneHL{state}{S_2})&=\GTenToneHL{quantity}{7},
		\\
		\GTenToneHL{quantity}{q}(\GTenToneHL{alternative}{A_3},\GTenToneHL{state}{S_1})&=\GTenToneHL{quantity}{14},
		&
		\GTenToneHL{quantity}{q}(\GTenToneHL{alternative}{A_3},\GTenToneHL{state}{S_2})&=8.
	\end{aligned}
	\]
	
	The \GTenToneHL{alternative}{alternative}-based \GTenToneHL{revenue}{revenue} rates, in thousand USD per hundred
	order batches, are
	\[
	\begin{aligned}
		\GTenToneHL{revenue}{r}(\GTenToneHL{alternative}{A_1})&=\GTenToneHL{revenue}{12},
		&
		\GTenToneHL{revenue}{r}(\GTenToneHL{alternative}{A_2})&=\GTenToneHL{revenue}{15},
		&
		\GTenToneHL{revenue}{r}(\GTenToneHL{alternative}{A_3})&=18.
	\end{aligned}
	\]
	
	The state-based \GTenToneHL{revenue}{revenue} rates are
	\[
	\GTenToneHL{revenue}{r}(\GTenToneHL{state}{S_1})=\GTenToneHL{revenue}{4},
	\qquad
	\GTenToneHL{revenue}{r}(\GTenToneHL{state}{S_2})=2.
	\]
	
	\columnbreak
	
	The \GTenToneHL{alternative}{alternative}-based fixed revenues, in thousand USD, are
	\[
	\begin{aligned}
		\GTenToneHL{revenue}{f_r}(\GTenToneHL{alternative}{A_1})&=\GTenToneHL{revenue}{16},
		&
		\GTenToneHL{revenue}{f_r}(\GTenToneHL{alternative}{A_2})&=\GTenToneHL{revenue}{24},
		&
		\GTenToneHL{revenue}{f_r}(\GTenToneHL{alternative}{A_3})&=35.
	\end{aligned}
	\]
	
	The state-based fixed revenues, in thousand USD, are
	\[
	\GTenToneHL{revenue}{f_r}(\GTenToneHL{state}{S_1})=\GTenToneHL{revenue}{15},
	\qquad
	\GTenToneHL{revenue}{f_r}(\GTenToneHL{state}{S_2})=8.
	\]
	
	The \GTenToneHL{alternative}{alternative}-based \GTenToneHL{cost}{cost} rates, in thousand USD per hundred
	order batches, are
	\[
	\begin{aligned}
		\GTenToneHL{cost}{c}(\GTenToneHL{alternative}{A_1})&=\GTenToneHL{cost}{5},
		&
		\GTenToneHL{cost}{c}(\GTenToneHL{alternative}{A_2})&=\GTenToneHL{cost}{6},
		&
		\GTenToneHL{cost}{c}(\GTenToneHL{alternative}{A_3})&=8.
	\end{aligned}
	\]
	
	The state-based \GTenToneHL{cost}{cost} rates are
	\[
	\GTenToneHL{cost}{c}(\GTenToneHL{state}{S_1})=\GTenToneHL{cost}{2},
	\qquad
	\GTenToneHL{cost}{c}(\GTenToneHL{state}{S_2})=4.
	\]
	
	The \GTenToneHL{alternative}{alternative}-based fixed \GTenToneHL{cost}{costs}, in thousand USD, are
	\[
	\begin{aligned}
		\GTenToneHL{cost}{f_c}(\GTenToneHL{alternative}{A_1})&=\GTenToneHL{cost}{22},
		&
		\GTenToneHL{cost}{f_c}(\GTenToneHL{alternative}{A_2})&=\GTenToneHL{cost}{34},
		&
		\GTenToneHL{cost}{f_c}(\GTenToneHL{alternative}{A_3})&=50.
	\end{aligned}
	\]
	
	The state-based fixed \GTenToneHL{cost}{costs}, in thousand USD, are
	\[
	\GTenToneHL{cost}{f_c}(\GTenToneHL{state}{S_1})=\GTenToneHL{cost}{9},
	\qquad
	\GTenToneHL{cost}{f_c}(\GTenToneHL{state}{S_2})=16.
	\]
\end{multicols}

\textbf{Questions/Tasks.}
\begin{enumerate}
	\CriterionTask
	{subsec:problem20-distribution-centre-automation-c1}{1}
	{Interpreting decision alternatives, events, consequences, and states.}

	\CriterionTask
	{subsec:problem20-distribution-centre-automation-c3}{3}
	{Calculating profit and building the payoff table.}
\end{enumerate}

\subsubsection*{C1 - Interpreting decision alternatives, events, consequences, and states.}
\phantomsection
\label{subsec:problem20-distribution-centre-automation-c1}

\DecisionElementsTable
{Use the manual distribution configuration; Use the semi-automated distribution configuration; Use the robotic distribution configuration}
{High order volume with favourable energy prices; Moderate order volume with expensive energy}
{Order volume and energy-market conditions revealed after the distribution-configuration decision}
{Profit in thousand USD for each distribution-configuration--state pair}

\GTenToneHL{revenue}{Revenue} contains \GTenToneHL{alternative}{alternative}-based and state-based variable components and
\GTenToneHL{alternative}{alternative}-based and state-based fixed components:
\[
\GTenToneHL{revenue}{R}(\GTenToneHL{alternative}{A_i},\GTenToneHL{state}{S_j})
=
\GTenToneHL{quantity}{q}(\GTenToneHL{alternative}{A_i},\GTenToneHL{state}{S_j})
\left[
\GTenToneHL{revenue}{r}(\GTenToneHL{alternative}{A_i})+\GTenToneHL{revenue}{r}(\GTenToneHL{state}{S_j})
\right]
+
\GTenToneHL{revenue}{f_r}(\GTenToneHL{alternative}{A_i})
+
\GTenToneHL{revenue}{f_r}(\GTenToneHL{state}{S_j}).
\]

\GTenToneHL{cost}{Cost} also contains \GTenToneHL{alternative}{alternative}-based and state-based variable components and
\GTenToneHL{alternative}{alternative}-based and state-based fixed components:
\[
\GTenToneHL{cost}{C}(\GTenToneHL{alternative}{A_i},\GTenToneHL{state}{S_j})
=
\GTenToneHL{quantity}{q}(\GTenToneHL{alternative}{A_i},\GTenToneHL{state}{S_j})
\left[
\GTenToneHL{cost}{c}(\GTenToneHL{alternative}{A_i})+\GTenToneHL{cost}{c}(\GTenToneHL{state}{S_j})
\right]
+
\GTenToneHL{cost}{f_c}(\GTenToneHL{alternative}{A_i})
+
\GTenToneHL{cost}{f_c}(\GTenToneHL{state}{S_j}).
\]

Therefore,
\[
\begin{aligned}
	\GTenToneHL{payoff}{P}(\GTenToneHL{alternative}{A_i},\GTenToneHL{state}{S_j})
	={}&
	\left\{
	\GTenToneHL{quantity}{q}(\GTenToneHL{alternative}{A_i},\GTenToneHL{state}{S_j})
	\left[
	\GTenToneHL{revenue}{r}(\GTenToneHL{alternative}{A_i})+\GTenToneHL{revenue}{r}(\GTenToneHL{state}{S_j})
	\right]
	+
	\GTenToneHL{revenue}{f_r}(\GTenToneHL{alternative}{A_i})
	+
	\GTenToneHL{revenue}{f_r}(\GTenToneHL{state}{S_j})
	\right\}
	\\
	&-
	\left\{
	\GTenToneHL{quantity}{q}(\GTenToneHL{alternative}{A_i},\GTenToneHL{state}{S_j})
	\left[
	\GTenToneHL{cost}{c}(\GTenToneHL{alternative}{A_i})+\GTenToneHL{cost}{c}(\GTenToneHL{state}{S_j})
	\right]
	+
	\GTenToneHL{cost}{f_c}(\GTenToneHL{alternative}{A_i})
	+
	\GTenToneHL{cost}{f_c}(\GTenToneHL{state}{S_j})
	\right\}.
\end{aligned}
\]

\secInicio[sec:problem20-distribution-centre-automation-c1-c3][sec:complete-models]

\subsubsection*{C3 - Calculating profit and building the payoff table.}
\phantomsection
\label{subsec:problem20-distribution-centre-automation-c3}

For every \GTenToneHL{alternative}{alternative}--state pair, calculate
\[
\begin{aligned}
	\GTenToneHL{revenue}{R}(\GTenToneHL{alternative}{A_i},\GTenToneHL{state}{S_j})
	&=
	\GTenToneHL{quantity}{q}(\GTenToneHL{alternative}{A_i},\GTenToneHL{state}{S_j})
	\left[
	\GTenToneHL{revenue}{r}(\GTenToneHL{alternative}{A_i})+\GTenToneHL{revenue}{r}(\GTenToneHL{state}{S_j})
	\right]
	+
	\GTenToneHL{revenue}{f_r}(\GTenToneHL{alternative}{A_i})
	+
	\GTenToneHL{revenue}{f_r}(\GTenToneHL{state}{S_j}),
	\\
	\GTenToneHL{cost}{C}(\GTenToneHL{alternative}{A_i},\GTenToneHL{state}{S_j})
	&=
	\GTenToneHL{quantity}{q}(\GTenToneHL{alternative}{A_i},\GTenToneHL{state}{S_j})
	\left[
	\GTenToneHL{cost}{c}(\GTenToneHL{alternative}{A_i})+\GTenToneHL{cost}{c}(\GTenToneHL{state}{S_j})
	\right]
	+
	\GTenToneHL{cost}{f_c}(\GTenToneHL{alternative}{A_i})
	+
	\GTenToneHL{cost}{f_c}(\GTenToneHL{state}{S_j}),
	\\
	\GTenToneHL{payoff}{P}(\GTenToneHL{alternative}{A_i},\GTenToneHL{state}{S_j})
	&=
	\GTenToneHL{revenue}{R}(\GTenToneHL{alternative}{A_i},\GTenToneHL{state}{S_j})-\GTenToneHL{cost}{C}(\GTenToneHL{alternative}{A_i},\GTenToneHL{state}{S_j}).
\end{aligned}
\]

\begin{center}
	\small
	\begin{tabular}{llll}
		\toprule
		Cell
		& \(\GTenToneHL{revenue}{R}(\GTenToneHL{alternative}{A_i},\GTenToneHL{state}{S_j})\)
		& \(\GTenToneHL{cost}{C}(\GTenToneHL{alternative}{A_i},\GTenToneHL{state}{S_j})\)
		& \(\GTenToneHL{payoff}{P}(\GTenToneHL{alternative}{A_i},\GTenToneHL{state}{S_j})=\GTenToneHL{revenue}{R}(\GTenToneHL{alternative}{A_i},\GTenToneHL{state}{S_j})-\GTenToneHL{cost}{C}(\GTenToneHL{alternative}{A_i},\GTenToneHL{state}{S_j})\) \\
		\midrule
		\((\GTenToneHL{alternative}{A_1},\GTenToneHL{state}{S_1})\)
		& \(8(12+4)+16+15=159\)
		& \(8(5+2)+22+9=87\)
		& \(159-87=72\) \\
		
		\((\GTenToneHL{alternative}{A_1},\GTenToneHL{state}{S_2})\)
		& \(5(12+2)+16+8=94\)
		& \(5(5+4)+22+16=83\)
		& \(94-83=11\) \\
		
		\((\GTenToneHL{alternative}{A_2},\GTenToneHL{state}{S_1})\)
		& \(11(15+4)+24+15=248\)
		& \(11(6+2)+34+9=131\)
		& \(248-131=117\) \\
		
		\((\GTenToneHL{alternative}{A_2},\GTenToneHL{state}{S_2})\)
		& \(7(15+2)+24+8=151\)
		& \(7(6+4)+34+16=120\)
		& \(151-120=31\) \\
		
		\((\GTenToneHL{alternative}{A_3},\GTenToneHL{state}{S_1})\)
		& \(14(18+4)+35+15=358\)
		& \(14(8+2)+50+9=199\)
		& \(358-199=159\) \\
		
		\((\GTenToneHL{alternative}{A_3},\GTenToneHL{state}{S_2})\)
		& \(8(18+2)+35+8=203\)
		& \(8(8+4)+50+16=162\)
		& \(203-162=41\) \\
		\bottomrule
	\end{tabular}
\end{center}

The resulting \GTenToneHL{payoff}{payoff} table is:

\begin{center}
	\begin{tabular}{lrr}
		\toprule
		Alternative
		& High order volume
		& Moderate order volume \\
		\midrule
		Manual distribution configuration & \GTenToneHL{payoff}{72} & \GTenToneHL{rate}{11} \\
		Semi-automated distribution configuration & \GTenToneHL{payoff}{117} & \GTenToneHL{rate}{31} \\
		Robotic distribution configuration & \GTenToneHL{payoff}{159} & \GTenToneHL{rate}{41} \\
		\bottomrule
	\end{tabular}
\end{center}

\secInicio[sec:problem20-distribution-centre-automation-c1-c3][sec:complete-models]

% ----------------------------------------------------------------
\subsection{Electric Vehicle Charging Network (C1--C3)}
\label{sec:problem21-electric-vehicle-charging-network-c1-c3}

An energy company must select an electric vehicle charging-network design
before charging demand and electricity-market conditions are known. \GTenToneHL{revenue}{Revenue} per
charging-session batch combines a rate determined by the selected network and a
rate determined by the state. \GTenToneHL{revenue}{Revenue} also includes a fixed service payment
determined by the network and a fixed utilisation incentive determined by the
state.

Variable \GTenToneHL{cost}{cost} per charging-session batch combines a network-based equipment
rate and a state-based electricity rate. Fixed \GTenToneHL{cost}{cost} consists of an
infrastructure \GTenToneHL{cost}{cost} determined by the network and a maintenance \GTenToneHL{cost}{cost} determined
by the state.

\begin{multicols}{2}
	The company is comparing two \GTenToneHL{alternative}{alternatives}:
	\begin{enumerate}[label=\arabic*.]
		\item \(\GTenToneHL{alternative}{A_1}\): Build the neighbourhood charging network
		\item \(\GTenToneHL{alternative}{A_2}\): Build the rapid-charging network
	\end{enumerate}
	
	The possible \GTenToneHL{state}{states} \GTenToneHL{state}{of} \GTenToneHL{state}{nature} are:
	\begin{enumerate}[label=\arabic*.]
		\item \(\GTenToneHL{state}{S_1}\): High utilisation with favourable electricity prices
		\item \(\GTenToneHL{state}{S_2}\): Moderate utilisation with standard electricity prices
		\item \(\GTenToneHL{state}{S_3}\): Low utilisation with expensive electricity
	\end{enumerate}
	
	The quantities, measured in hundreds of charging-session batches, are
	\[
	\begin{aligned}
		\GTenToneHL{quantity}{q}(\GTenToneHL{alternative}{A_1},\GTenToneHL{state}{S_1})&=\GTenToneHL{quantity}{10},
		&
		\GTenToneHL{quantity}{q}(\GTenToneHL{alternative}{A_1},\GTenToneHL{state}{S_2})&=\GTenToneHL{quantity}{7},
		&
		\GTenToneHL{quantity}{q}(\GTenToneHL{alternative}{A_1},\GTenToneHL{state}{S_3})&=\GTenToneHL{quantity}{4},
		\\
		\GTenToneHL{quantity}{q}(\GTenToneHL{alternative}{A_2},\GTenToneHL{state}{S_1})&=\GTenToneHL{quantity}{14},
		&
		\GTenToneHL{quantity}{q}(\GTenToneHL{alternative}{A_2},\GTenToneHL{state}{S_2})&=\GTenToneHL{quantity}{9},
		&
		\GTenToneHL{quantity}{q}(\GTenToneHL{alternative}{A_2},\GTenToneHL{state}{S_3})&=5.
	\end{aligned}
	\]
	
	The \GTenToneHL{alternative}{alternative}-based \GTenToneHL{revenue}{revenue} rates, in thousand USD per hundred
	charging-session batches, are
	\[
	\GTenToneHL{revenue}{r}(\GTenToneHL{alternative}{A_1})=\GTenToneHL{revenue}{11},
	\qquad
	\GTenToneHL{revenue}{r}(\GTenToneHL{alternative}{A_2})=16.
	\]
	
	The state-based \GTenToneHL{revenue}{revenue} rates are
	\[
	\begin{aligned}
		\GTenToneHL{revenue}{r}(\GTenToneHL{state}{S_1})&=\GTenToneHL{revenue}{4},
		&
		\GTenToneHL{revenue}{r}(\GTenToneHL{state}{S_2})&=\GTenToneHL{revenue}{3},
		&
		\GTenToneHL{revenue}{r}(\GTenToneHL{state}{S_3})&=2.
	\end{aligned}
	\]
	
	\columnbreak
	
	The \GTenToneHL{alternative}{alternative}-based fixed revenues, in thousand USD, are
	\[
	\GTenToneHL{revenue}{f_r}(\GTenToneHL{alternative}{A_1})=\GTenToneHL{revenue}{18},
	\qquad
	\GTenToneHL{revenue}{f_r}(\GTenToneHL{alternative}{A_2})=30.
	\]
	
	The state-based fixed revenues, in thousand USD, are
	\[
	\begin{aligned}
		\GTenToneHL{revenue}{f_r}(\GTenToneHL{state}{S_1})&=\GTenToneHL{revenue}{16},
		&
		\GTenToneHL{revenue}{f_r}(\GTenToneHL{state}{S_2})&=\GTenToneHL{revenue}{10},
		&
		\GTenToneHL{revenue}{f_r}(\GTenToneHL{state}{S_3})&=5.
	\end{aligned}
	\]
	
	The \GTenToneHL{alternative}{alternative}-based \GTenToneHL{cost}{cost} rates, in thousand USD per hundred
	charging-session batches, are
	\[
	\GTenToneHL{cost}{c}(\GTenToneHL{alternative}{A_1})=\GTenToneHL{cost}{4},
	\qquad
	\GTenToneHL{cost}{c}(\GTenToneHL{alternative}{A_2})=7.
	\]
	
	The state-based \GTenToneHL{cost}{cost} rates are
	\[
	\begin{aligned}
		\GTenToneHL{cost}{c}(\GTenToneHL{state}{S_1})&=\GTenToneHL{cost}{2},
		&
		\GTenToneHL{cost}{c}(\GTenToneHL{state}{S_2})&=\GTenToneHL{cost}{3},
		&
		\GTenToneHL{cost}{c}(\GTenToneHL{state}{S_3})&=5.
	\end{aligned}
	\]
	
	The \GTenToneHL{alternative}{alternative}-based fixed \GTenToneHL{cost}{costs}, in thousand USD, are
	\[
	\GTenToneHL{cost}{f_c}(\GTenToneHL{alternative}{A_1})=\GTenToneHL{cost}{25},
	\qquad
	\GTenToneHL{cost}{f_c}(\GTenToneHL{alternative}{A_2})=45.
	\]
	
	The state-based fixed \GTenToneHL{cost}{costs}, in thousand USD, are
	\[
	\begin{aligned}
		\GTenToneHL{cost}{f_c}(\GTenToneHL{state}{S_1})&=\GTenToneHL{cost}{8},
		&
		\GTenToneHL{cost}{f_c}(\GTenToneHL{state}{S_2})&=\GTenToneHL{cost}{12},
		&
		\GTenToneHL{cost}{f_c}(\GTenToneHL{state}{S_3})&=18.
	\end{aligned}
	\]
\end{multicols}

\textbf{Questions/Tasks.}
\begin{enumerate}
	\CriterionTask
	{subsec:problem21-electric-vehicle-charging-network-c1}{1}
	{Interpreting decision alternatives, events, consequences, and states.}

	\CriterionTask
	{subsec:problem21-electric-vehicle-charging-network-c3}{3}
	{Calculating profit and building the payoff table.}
\end{enumerate}

\subsubsection*{C1 - Interpreting decision alternatives, events, consequences, and states.}
\phantomsection
\label{subsec:problem21-electric-vehicle-charging-network-c1}

\DecisionElementsTable
{Build the neighbourhood charging network; Build the rapid-charging network}
{High utilisation with favourable electricity prices; Moderate utilisation with standard electricity prices; Low utilisation with expensive electricity}
{Charging demand and electricity-market conditions revealed after the network-design decision}
{Profit in thousand USD for each charging-network--state pair}

\GTenToneHL{revenue}{Revenue} contains \GTenToneHL{alternative}{alternative}-based and state-based variable components and
\GTenToneHL{alternative}{alternative}-based and state-based fixed components:
\[
\GTenToneHL{revenue}{R}(\GTenToneHL{alternative}{A_i},\GTenToneHL{state}{S_j})
=
\GTenToneHL{quantity}{q}(\GTenToneHL{alternative}{A_i},\GTenToneHL{state}{S_j})
\left[
\GTenToneHL{revenue}{r}(\GTenToneHL{alternative}{A_i})+\GTenToneHL{revenue}{r}(\GTenToneHL{state}{S_j})
\right]
+
\GTenToneHL{revenue}{f_r}(\GTenToneHL{alternative}{A_i})
+
\GTenToneHL{revenue}{f_r}(\GTenToneHL{state}{S_j}).
\]

\GTenToneHL{cost}{Cost} also contains \GTenToneHL{alternative}{alternative}-based and state-based variable components and
\GTenToneHL{alternative}{alternative}-based and state-based fixed components:
\[
\GTenToneHL{cost}{C}(\GTenToneHL{alternative}{A_i},\GTenToneHL{state}{S_j})
=
\GTenToneHL{quantity}{q}(\GTenToneHL{alternative}{A_i},\GTenToneHL{state}{S_j})
\left[
\GTenToneHL{cost}{c}(\GTenToneHL{alternative}{A_i})+\GTenToneHL{cost}{c}(\GTenToneHL{state}{S_j})
\right]
+
\GTenToneHL{cost}{f_c}(\GTenToneHL{alternative}{A_i})
+
\GTenToneHL{cost}{f_c}(\GTenToneHL{state}{S_j}).
\]

Therefore,
\[
\begin{aligned}
	\GTenToneHL{payoff}{P}(\GTenToneHL{alternative}{A_i},\GTenToneHL{state}{S_j})
	={}&
	\left\{
	\GTenToneHL{quantity}{q}(\GTenToneHL{alternative}{A_i},\GTenToneHL{state}{S_j})
	\left[
	\GTenToneHL{revenue}{r}(\GTenToneHL{alternative}{A_i})+\GTenToneHL{revenue}{r}(\GTenToneHL{state}{S_j})
	\right]
	+
	\GTenToneHL{revenue}{f_r}(\GTenToneHL{alternative}{A_i})
	+
	\GTenToneHL{revenue}{f_r}(\GTenToneHL{state}{S_j})
	\right\}
	\\
	&-
	\left\{
	\GTenToneHL{quantity}{q}(\GTenToneHL{alternative}{A_i},\GTenToneHL{state}{S_j})
	\left[
	\GTenToneHL{cost}{c}(\GTenToneHL{alternative}{A_i})+\GTenToneHL{cost}{c}(\GTenToneHL{state}{S_j})
	\right]
	+
	\GTenToneHL{cost}{f_c}(\GTenToneHL{alternative}{A_i})
	+
	\GTenToneHL{cost}{f_c}(\GTenToneHL{state}{S_j})
	\right\}.
\end{aligned}
\]

\secInicio[sec:problem21-electric-vehicle-charging-network-c1-c3][sec:complete-models]

\subsubsection*{C3 - Calculating profit and building the payoff table.}
\phantomsection
\label{subsec:problem21-electric-vehicle-charging-network-c3}

For every \GTenToneHL{alternative}{alternative}--state pair, calculate
\[
\begin{aligned}
	\GTenToneHL{revenue}{R}(\GTenToneHL{alternative}{A_i},\GTenToneHL{state}{S_j})
	&=
	\GTenToneHL{quantity}{q}(\GTenToneHL{alternative}{A_i},\GTenToneHL{state}{S_j})
	\left[
	\GTenToneHL{revenue}{r}(\GTenToneHL{alternative}{A_i})+\GTenToneHL{revenue}{r}(\GTenToneHL{state}{S_j})
	\right]
	+
	\GTenToneHL{revenue}{f_r}(\GTenToneHL{alternative}{A_i})
	+
	\GTenToneHL{revenue}{f_r}(\GTenToneHL{state}{S_j}),
	\\
	\GTenToneHL{cost}{C}(\GTenToneHL{alternative}{A_i},\GTenToneHL{state}{S_j})
	&=
	\GTenToneHL{quantity}{q}(\GTenToneHL{alternative}{A_i},\GTenToneHL{state}{S_j})
	\left[
	\GTenToneHL{cost}{c}(\GTenToneHL{alternative}{A_i})+\GTenToneHL{cost}{c}(\GTenToneHL{state}{S_j})
	\right]
	+
	\GTenToneHL{cost}{f_c}(\GTenToneHL{alternative}{A_i})
	+
	\GTenToneHL{cost}{f_c}(\GTenToneHL{state}{S_j}),
	\\
	\GTenToneHL{payoff}{P}(\GTenToneHL{alternative}{A_i},\GTenToneHL{state}{S_j})
	&=
	\GTenToneHL{revenue}{R}(\GTenToneHL{alternative}{A_i},\GTenToneHL{state}{S_j})-\GTenToneHL{cost}{C}(\GTenToneHL{alternative}{A_i},\GTenToneHL{state}{S_j}).
\end{aligned}
\]

\begin{center}
	\small
	\begin{tabular}{llll}
		\toprule
		Cell
		& \(\GTenToneHL{revenue}{R}(\GTenToneHL{alternative}{A_i},\GTenToneHL{state}{S_j})\)
		& \(\GTenToneHL{cost}{C}(\GTenToneHL{alternative}{A_i},\GTenToneHL{state}{S_j})\)
		& \(\GTenToneHL{payoff}{P}(\GTenToneHL{alternative}{A_i},\GTenToneHL{state}{S_j})=\GTenToneHL{revenue}{R}(\GTenToneHL{alternative}{A_i},\GTenToneHL{state}{S_j})-\GTenToneHL{cost}{C}(\GTenToneHL{alternative}{A_i},\GTenToneHL{state}{S_j})\) \\
		\midrule
		\((\GTenToneHL{alternative}{A_1},\GTenToneHL{state}{S_1})\)
		& \(10(11+4)+18+16=184\)
		& \(10(4+2)+25+8=93\)
		& \(184-93=91\) \\
		
		\((\GTenToneHL{alternative}{A_1},\GTenToneHL{state}{S_2})\)
		& \(7(11+3)+18+10=126\)
		& \(7(4+3)+25+12=86\)
		& \(126-86=40\) \\
		
		\((\GTenToneHL{alternative}{A_1},\GTenToneHL{state}{S_3})\)
		& \(4(11+2)+18+5=75\)
		& \(4(4+5)+25+18=79\)
		& \(75-79=-4\) \\
		
		\((\GTenToneHL{alternative}{A_2},\GTenToneHL{state}{S_1})\)
		& \(14(16+4)+30+16=326\)
		& \(14(7+2)+45+8=179\)
		& \(326-179=147\) \\
		
		\((\GTenToneHL{alternative}{A_2},\GTenToneHL{state}{S_2})\)
		& \(9(16+3)+30+10=211\)
		& \(9(7+3)+45+12=147\)
		& \(211-147=64\) \\
		
		\((\GTenToneHL{alternative}{A_2},\GTenToneHL{state}{S_3})\)
		& \(5(16+2)+30+5=125\)
		& \(5(7+5)+45+18=123\)
		& \(125-123=2\) \\
		\bottomrule
	\end{tabular}
\end{center}

The resulting \GTenToneHL{payoff}{payoff} table is:

\begin{center}
	\begin{tabular}{lrrr}
		\toprule
		Alternative
		& High utilisation
		& Moderate utilisation
		& Low utilisation \\
		\midrule
		Neighbourhood charging network & \GTenToneHL{payoff}{91} & \GTenToneHL{rate}{40} & \(-4\) \\
		Rapid-charging network & \GTenToneHL{payoff}{147} & \GTenToneHL{rate}{64} & \GTenToneHL{payoff}{2} \\
		\bottomrule
	\end{tabular}
\end{center}

The negative \GTenToneHL{payoff}{payoff} for the neighbourhood charging network under low
utilisation represents a loss because the \GTenToneHL{revenue}{revenue} earned does not cover the
variable and fixed \GTenToneHL{cost}{costs}.

\secInicio[sec:problem21-electric-vehicle-charging-network-c1-c3][sec:complete-models]

\end{document}
